% Il Bilanciere Quantistico — Stampa (6×9" KDP)
\documentclass[12pt, twoside]{book}
\usepackage[utf8]{inputenc}
\usepackage[T1]{fontenc}
\usepackage[italian]{babel}
\usepackage{amsmath,amssymb,amsthm,mathtools,bm}
\usepackage{geometry}
\geometry{paperwidth=15.24cm, paperheight=22.86cm,
  inner=2.8cm, outer=1.8cm, top=2.2cm, bottom=2.5cm}
\usepackage{setspace}\onehalfspacing
\setlength{\headheight}{14.5pt}\addtolength{\topmargin}{-2.5pt}
\usepackage[protrusion=true, expansion=false]{microtype}
\usepackage{emptypage}
\usepackage{graphicx}\graphicspath{{figures/}}
\usepackage{booktabs, array, longtable, multirow}
\usepackage[hidelinks]{hyperref}
\usepackage{enumitem, xcolor}
\usepackage{fancyhdr}
\usepackage{titlesec}
\usepackage{algorithm, algpseudocode}
\floatname{algorithm}{Algoritmo}

% ─── Stile capitoli ──────────────────────────────────────────────────────────
\titleformat{\chapter}[display]
  {\normalfont\Large\bfseries}{\chaptertitlename~\thechapter}{12pt}{\Huge}
\titlespacing*{\chapter}{0pt}{30pt}{20pt}

\titleformat{\part}[display]
  {\centering\normalfont\huge\bfseries}{Parte~\thepart}{20pt}{\Huge}

% ─── Header/footer ───────────────────────────────────────────────────────────
\pagestyle{fancy}
\fancyhf{}
\fancyhead[LE]{\small\textit{Il Bilanciere Quantistico}}
\fancyhead[RO]{\small\textit{\leftmark}}
\fancyfoot[C]{\thepage}
\renewcommand{\headrulewidth}{0.4pt}

% ─── Teoremi e ambienti formali ──────────────────────────────────────────────
\theoremstyle{plain}
\newtheorem{theorem}{Teorema}[chapter]
\newtheorem{lemma}[theorem]{Lemma}
\newtheorem{corollary}[theorem]{Corollario}
\newtheorem{proposition}[theorem]{Proposizione}
\theoremstyle{definition}
\newtheorem{definition}{Definizione}[chapter]
\newtheorem{conjecture}[definition]{Congettura}
\newtheorem{hypothesis}[definition]{Ipotesi}
\theoremstyle{remark}
\newtheorem{remark}{Osservazione}[chapter]
\newtheorem{openprob}{Problema aperto}[chapter]

% ─── Macro epistemiche (dalle emoji del vault) ────────────────────────────────
% ✅ → \Proved{}  (affermazione dimostrata)
% 🔵 → \Hypoth{}  (ipotesi operativa)
% 🟡 → \Conj{}    (congettura)
% ❌ → \ToVerify{} (da verificare)
\newcommand{\Proved}{\textnormal{\textbf{[DIMOSTRATO]}}}
\newcommand{\Hypoth}{\textnormal{\textbf{[IPOTESI OPERATIVA]}}}
\newcommand{\Conj}{\textnormal{\textbf{[CONGETTURA]}}}
\newcommand{\ToVerify}{\textnormal{\textbf{[DA VERIFICARE]}}}

% ─── Compatibilità pandoc ─────────────────────────────────────────────────────
\providecommand{\tightlist}{\setlength{\itemsep}{0pt}\setlength{\parskip}{0pt}}
\newcounter{none}  % pandoc usa \LTcaptype{none} per longtable senza caption

% ─── Macro matematiche frequenti ─────────────────────────────────────────────
\newcommand{\Ts}{T_s^{(2),\text{sys}}}
\newcommand{\deq}{\mathrm{DEQ}^2}

% ─── Metadati PDF ────────────────────────────────────────────────────────────
\hypersetup{
  pdftitle={Il Bilanciere Quantistico},
  pdfauthor={Luis Fernando Marcelino},
  pdflang={it},
  pdfsubject={Geopolitica, Sistemi Complessi, Transizione Egemonica},
  pdfkeywords={DEQ2, egemonia, Kuramoto, transizione di fase, Brasile, USA, Cina, geopolitica formale}
}

% ─── Documento ───────────────────────────────────────────────────────────────
\begin{document}

%% ── Pagina titolo ─────────────────────────────────────────────────────────
\begin{titlepage}
  \centering\vspace*{1.5cm}
  {\Huge\bfseries Il Bilanciere\par}
  \vspace{0.2cm}
  {\Huge\bfseries Quantistico\par}
  \vspace{0.8cm}
  \rule{0.8\linewidth}{0.5pt}
  \vspace{0.6cm}

  {\large Geometria dell'egemonia\\[0.3cm]
   e transizione di fase 2028--2031\par}
  \vspace{2cm}
  {\large Luis Fernando Marcelino\par}
  \vspace{0.5cm}
  {\normalsize Milano, 2026\par}
  \vfill
  {\small\textit{Collana Nuovi Orizonti} --- Volume~4\par}
  \vspace{0.5cm}
  \begin{small}\begin{quote}
    \textit{Etichette epistemiche usate nel testo:}\\[2pt]
    \Proved{} verificato formalmente o su tutti i casi esaminati;\\
    \Hypoth{} ipotesi operativa dichiarata;\\
    \Conj{} congettura osservazionale;\\
    \ToVerify{} affermazione da verificare empiricamente.
  \end{quote}\end{small}
\end{titlepage}

%% ── Pagina copyright ──────────────────────────────────────────────────────
\thispagestyle{empty}
\vspace*{\fill}
\begin{small}
\noindent \textcopyright{} 2026 Luis Fernando Marcelino\\[4pt]
Tutti i diritti riservati.\\[4pt]
Prima edizione digitale e cartacea, 2026.\\[4pt]
\textit{Il Bilanciere Quantistico} è il quarto volume della Collana \textit{Nuovi Orizonti}.\\[4pt]
Fa seguito a:\\
1. \textit{La Costante di Marcelino} (Marcelino \& Teixeira, 2026)\\
2. \textit{La Geometria dell'Egemonia} (Marcelino, 2026)\\
3. \textit{MQNU — Modello Quantistico della Natura Umana} (Marcelino, 2026)\\[4pt]
La microfondazione del modello DEQ² è sviluppata in:\\
\textit{La Topografia del Discorso Emancipativo} (Marcelino, 2026, 170pp).
\end{small}
\vspace*{\fill}
\clearpage

%% ── Indice ────────────────────────────────────────────────────────────────
\pagenumbering{roman}
\tableofcontents
\cleardoublepage

%% ── Prefazione Tecnica ────────────────────────────────────────────────────
\chapter*{Prefazione Tecnica}
\addcontentsline{toc}{chapter}{Prefazione Tecnica}
\markboth{Prefazione Tecnica}{Prefazione Tecnica}

\textbf{Posizione nel progetto intellettuale.}
Questo volume è il quarto della Collana \textit{Nuovi Orizonti} e presuppone — senza richiedere — la lettura dei tre precedenti. Il lettore nuovo può procedere linearmente: i Capitoli~1 e~2 sintetizzano l'apparato necessario dalla \textit{Topografia} e dalla \textit{Geometria}; il Capitolo~3 costruisce il ponte tra le due scale. Chi ha già letto i volumi precedenti troverà nei primi due capitoli una sintesi compatta e nei successivi l'estensione originale.

\textbf{Nota sulla meccanica quantistica come isomorfismo strutturale.}
Il titolo contiene la parola ``quantistico''. Questa scelta richiede chiarimento preventivo. Non si afferma che la geopolitica \textit{sia} meccanica quantistica, né che i capi di stato si comportino come particelle. Si afferma qualcosa di più preciso e più limitato: che la struttura matematica del modello di Kuramoto — transizioni di fase, parametro d'ordine, soglia critica, decoerenza — costituisce un \textit{isomorfismo strutturale} con certi pattern della dinamica egemonica. Come la termodinamica statistica di Boltzmann non afferma che i gas siano ``davvero'' probabilistici ma che lo strumento probabilistico descrive correttamente le proprietà macroscopiche, così la DEQ² usa il formalismo di Kuramoto come strumento descrittivo-predittivo, non come affermazione ontologica.

Il riferimento epistemico è sviluppato in \textit{MQNU — Modello Quantistico della Natura Umana} (Marcelino, 2026), in particolare nel Capitolo~1 (isomorfismo strutturale vs. riduzione ontologica) e nel Capitolo~4 (limiti dell'analogia).

\textbf{Struttura e come leggere.}
Il libro si divide in cinque parti. Le Parti~I--II costruiscono l'apparato teorico. La Parte~III analizza i quattro attori strutturali del sistema globale 2024--2031. La Parte~IV presenta la Dashboard SPC-DEQ e le predizioni falsificabili. La Parte~V fornisce protocolli d'azione per individui, organizzazioni e stati.

Le Appendici~A--D contengono le derivazioni formali complete, la calibrazione empirica, i protocolli di misurazione e il dialogo con le tradizioni intellettuali pertinenti. Sono progettate per essere lette autonomamente da lettori con background tecnico.

\textbf{Falsificabilità.}
Questo libro fa predizioni datate e numeriche. Il Capitolo~11 elenca 16~predizioni falsificabili con finestre temporali esplicite (2026--2031). Il protocollo di pre-registrazione pubblica è accessibile al repository GitHub \texttt{bilanciere-quantistico/predizioni-deq2}. Se le predizioni fallissero sistematicamente, il modello DEQ² andrebbe rivisto.

\cleardoublepage
\pagenumbering{arabic}

%% ═══════════════════════════════════════════════════════════════════════════
%% PARTE I — Dal Micro Comunicativo al Macro Geopolitico
%% ═══════════════════════════════════════════════════════════════════════════
\part[Dal Micro Comunicativo al Macro Geopolitico]{Dal Micro Comunicativo\\al Macro Geopolitico}

\chapter{La Topografia: l'Apparato
Micro}\label{la-topografia-lapparato-micro}

\begin{quote}
\textbf{Argomento portante:} ogni discorso occupa una posizione
misurabile in uno spazio tridimensionale (coerenza formale, precisione
tecnica, complessità dialettica); la sua efficacia sociale è il prodotto
di una \emph{massa comunicativa} per una \emph{qualità etica} diviso una
\emph{resistenza critica}; il sistema discorsivo è governato da una
legge predittiva dalla quale, per trasposizione di scala, deriverà
l'intera DEQ$^2$.
\end{quote}



\section{Perché Iniziare dalla
Topografia}\label{perchuxe9-iniziare-dalla-topografia}

La DEQ$^2$ è una teoria geopolitica. Iniziarla dalla teoria del discorso
può sembrare un'inversione di scala. Non lo è: l'egemonia \emph{è} prima
di tutto una struttura comunicativa, e la geopolitica del XXI secolo è
una geopolitica di flussi semantici prima ancora che di flussi
materiali. Il Cap. 3 dimostrerà formalmente che la legge fondamentale
della Topografia e quella della Geometria sono \textbf{la stessa
equazione applicata a due scale}. Per leggere quella dimostrazione, il
lettore ha bisogno di possedere l'apparato della Topografia --- questo
capitolo lo fornisce in forma compatta e autosufficiente.

Il riferimento esteso è: Marcelino, \emph{La Topografia del Discorso
Emancipativo} (2026, 170pp). Qui sintetizziamo solo ciò che entra nella
DEQ$^2$.



\section{La Domanda Originaria e la Lacuna
Operativa}\label{la-domanda-originaria-e-la-lacuna-operativa}

La Scuola di Francoforte ha prodotto un apparato critico di
straordinaria potenza --- Adorno, Horkheimer, Marcuse, Habermas --- ma
\textbf{senza strumenti operativi di misurazione}. Si poteva
\emph{descrivere} l'unidimensionalità del pensiero (Marcuse),
l'industria culturale (Adorno), la regressione dell'ascolto
(Horkheimer), la colonizzazione del mondo della vita (Habermas) --- ma
non \emph{misurarli}.

La Topografia colma questa lacuna costruendo uno spazio metrico in cui
ogni discorso può essere posizionato, confrontato, analizzato nella sua
traiettoria.



\section{\texorpdfstring{Lo Spazio Discorsivo
\(\mathcal{D} = [0, 10]^3\)}{Lo Spazio Discorsivo \textbackslash mathcal\{D\} = {[}0, 10{]}\^{}3}}\label{lo-spazio-discorsivo-mathcald-0-103}

Ogni discorso \(\mathbf{D} = (x, y, z)\) è un punto nello spazio
\(\mathcal{D}\), dove:

\begin{itemize}
\tightlist
\item
  \(x\) = \textbf{coerenza formale} (consistenza interna, assenza di
  contraddizioni logiche, struttura argomentativa);
\item
  \(y\) = \textbf{precisione tecnica} (rigore fattuale, accuratezza dei
  riferimenti, padronanza disciplinare);
\item
  \(z\) = \textbf{complessità dialettica} (capacità di pensare le
  contraddizioni, profondità critica, riflessività).
\end{itemize}

\Hypoth{} \textbf{Non-ortogonalità.} Gli assi non sono indipendenti. La
Topografia (\S{}5.1) stima la matrice di covarianza empirica:

\[
\Sigma = \begin{pmatrix} 1 & 0{,}40 & 0{,}30 \\ 0{,}40 & 1 & 0{,}50 \\ 0{,}30 & 0{,}50 & 1 \end{pmatrix}
\]

\Proved{} Significato: la covarianza più alta è \((y, z) = 0{,}50\) (la
critica autentica richiede rigore); la più bassa è \((x, z) = 0{,}30\)
(è possibile essere coerenti senza essere profondi --- il caso della
propaganda razionale).

\subsection{\texorpdfstring{Il Vettore Etico
\(\mathbf{E}\)}{Il Vettore Etico \textbackslash mathbf\{E\}}}\label{il-vettore-etico-mathbfe}

\[
\mathbf{E} = (10, 10, 10)
\]

è la \textbf{direzione di massima legittimazione discorsiva}. Non è
neutra --- incorpora la scelta di campo della Scuola di Francoforte:
l'emancipazione umana definita nei termini dei diritti umani universali.
La Topografia non aspira alla neutralità: aspira alla trasparenza sui
propri valori.



\section{Le Tre Metriche
Fondamentali}\label{le-tre-metriche-fondamentali}

\subsection{Qualità Etica del
Discorso}\label{qualituxe0-etica-del-discorso}

\[
Q_0 = \frac{\mathbf{D} \cdot \mathbf{E}}{\|\mathbf{E}\|} = \frac{x + y + z}{\sqrt{3}}
\]

Misura la \emph{distanza} del discorso dall'ideale etico. La forma
avanzata corretta dalle covarianze,
\(Q = \mathbf{D}^T \Sigma^{-1} \mathbf{E}\), è preferibile per analisi
statistiche; quella base è preferibile per la trasparenza
interpretativa.

\subsection{Massa Comunicativa
dell'Emittente}\label{massa-comunicativa-dellemittente}

\[
M_e = K_{\text{cap}} \cdot R^{0{,}8} \cdot A_{\text{alg}} \cdot S \cdot T
\]

dove: \(K_{\text{cap}}\) = capitale economico, \(R\) = reach (in
miliardi), \(A_{\text{alg}}\) = amplificazione algoritmica, \(S\) =
saturazione (ore/giorno), \(T\) = durata (mesi). L'esponente \(0{,}8\)
su \(R\) riflette i rendimenti decrescenti del reach puro.

\Hypoth{} \textbf{Disambiguazione notazionale.} Il simbolo
\(K_{\text{cap}}\) è usato per il capitale economico per distinguerlo
dal \(K\) del Cap. 4 --- \emph{forza di accoppiamento di Kuramoto},
parametro completamente diverso. Analogamente, \(A_{\text{alg}}\) resta
riservato all'amplificazione algoritmica, mentre \(A\) (senza pedice)
sarà introdotto nel Cap. 3 come \emph{antifragilità} (Taleb).

\subsection{Resistenza Critica del
Ricevente}\label{resistenza-critica-del-ricevente}

\[
R_c = z_r \cdot (I + D + C)
\]

dove \(z_r \in [0, 10]\) è la complessità dialettica media del pubblico,
\(I\) è l'istruzione critica, \(D\) la diversità delle fonti, \(C\) la
connettività con reti alternative. La forma moltiplicativa è cruciale:
\textbf{un pubblico con \(z_r = 0\) ha resistenza nulla
indipendentemente da \(I, D, C\)}.



\section{L'Impatto Bayesiano e le Tre
Leggi}\label{limpatto-bayesiano-e-le-tre-leggi}

\subsection{Impatto Atteso}\label{impatto-atteso}

\[
I \sim \mathrm{LogNormal}(\mu, \sigma^2), \qquad \mathbb{E}[I] = \frac{M_e \cdot Q}{R_c + 1}
\]

La distribuzione log-normale è giustificata da: positività intrinseca
dell'impatto, asimmetria empirica con coda pesante (virali), e
moltiplicatività dei fattori (TCL applicato al log).

\subsection{Le Tre Leggi}\label{le-tre-leggi}

\Proved{} \textbf{Prima Legge --- Inversione di Marcuse}:
\(\partial z / \partial M_e < 0\). La massimizzazione del reach implica
la minimizzazione dei prerequisiti cognitivi, dunque la compressione di
\(z\). Corollario: i discorsi egemonici si concentrano strutturalmente
nella regione \(z \in [1, 3]\).

\Proved{} \textbf{Seconda Legge --- Paradosso della Precisione}:
\(Q(10, 10, 2) < Q(7, 7, 9)\). Un discorso meno preciso ma
dialetticamente ricco ha qualità etica superiore. La complessità
dialettica pesa quanto la precisione.

\Proved{} \textbf{Terza Legge --- Teorema della Resistenza}:
\(\forall M_e > 0,\ \exists R_c\) tale che \(\mathbb{E}[I] < \epsilon\).
Per ogni massa comunicativa, esiste una resistenza critica sufficiente a
neutralizzarla. Politicamente: la \(R_c\) può essere costruita;
costruirla è il programma.



\section{L'Estensione Gametheoretica e la Funzione
Predittiva}\label{lestensione-gametheoretica-e-la-funzione-predittiva}

\subsection{Il Gioco Discorsivo}\label{il-gioco-discorsivo}

Due giocatori (Emittente \(E\), Ricevente \(R\)), due strategie
ciascuno: \(E \in \{\)Propaganda \(S^1_E\), Dialogo \(S^2_E\}\),
\(R \in \{\)Passivo \(C^1_R\), Critico \(C^2_R\}\). Matrice dei payoff:

{\def\LTcaptype{none} % do not increment counter
{\small\begin{longtable}[]{@{}lll@{}}
\toprule\noalign{}
& \(S^1_E\) Propaganda & \(S^2_E\) Dialogo \\
\midrule\noalign{}
\endhead
\bottomrule\noalign{}
\endlastfoot
\(C^1_R\) Passivo & \((10, -5)\) & \((2, 8)\) \\
\(C^2_R\) Critico & \((0, 2)\) & \((5, 5)\) \\
\end{longtable}}
}

L'ottimo paretiano è (Dialogo, Critico) \(= (5, 5)\) --- l'agire
comunicativo habermasiano. (Propaganda, Passivo) \(= (10, -5)\) è la
\textbf{Barbarie Tecnologica}.

\subsection{Il Folk Theorem
Discorsivo}\label{il-folk-theorem-discorsivo}

\Proved{} La cooperazione (Dialogo, Critico) è sostenibile come
equilibrio di Nash perfetto nei sottogiochi se e solo se:

\[
\delta_E \geq \delta^*_E = \frac{10 - 5}{10 - 0} = 0{,}50, \qquad \delta_R \geq \delta^*_R = \frac{10}{7} \approx 0{,}71
\]

\Hypoth{} \textbf{Punto cruciale per la DEQ$^2$.} La soglia
\(\delta^*_E = 0{,}50\) è la \textbf{stessa} che la \emph{Geometria
dell'Egemonia} (Cap. 3) deriva indipendentemente per la cooperazione tra
egemoni. Questa identità non è coincidenza --- è l'invarianza di scala
dimostrata nel Cap. 3 della DEQ$^2$.

\subsection{La Funzione Predittiva
Integrata}\label{la-funzione-predittiva-integrata}

\[
\boxed{\Pi(t+1) = M_e \cdot Q_0 \cdot (1 - \delta_R \cdot z_R)}
\]

Il termine \((1 - \delta_R z_R)\) è il \textbf{fattore di erosione}:
misura quanto la resistenza critica del ricevente riduce il payoff della
propaganda nel tempo.

\begin{itemize}
\tightlist
\item
  \(\delta_R z_R < 1\): propaganda efficace;
\item
  \(\delta_R z_R = 1\): propaganda inerte;
\item
  \(\delta_R z_R > 1\): \textbf{effetto boomerang} --- la pressione
  comunicativa rinforza la resistenza (\(\Pi < 0\)).
\end{itemize}

\Proved{} Questa equazione è la \textbf{legge micro fondamentale}
dell'intera DEQ$^2$. Tutto il Cap. 3-4 deriva da qui per trasposizione di
scala.



\section{Cosa la Topografia Porta alla
DEQ$^2$}\label{cosa-la-topografia-porta-alla-dequxb2}

{\def\LTcaptype{none} % do not increment counter
{\small\begin{longtable}[]{@{}
  >{\raggedright\arraybackslash}p{(\linewidth - 4\tabcolsep) * \real{0.3333}}
  >{\raggedright\arraybackslash}p{(\linewidth - 4\tabcolsep) * \real{0.3333}}
  >{\raggedright\arraybackslash}p{(\linewidth - 4\tabcolsep) * \real{0.3333}}@{}}
\toprule\noalign{}
\begin{minipage}[b]{\linewidth}\raggedright
Apparato Topografico
\end{minipage} & \begin{minipage}[b]{\linewidth}\raggedright
Ruolo nella DEQ$^2$
\end{minipage} & \begin{minipage}[b]{\linewidth}\raggedright
Riferimento
\end{minipage} \\
\midrule\noalign{}
\endhead
\bottomrule\noalign{}
\endlastfoot
Spazio \(\mathcal{D} = [0, 10]^3\) & Spazio degli stati discorsivi micro
& Cap. 1 \\
Matrice \(\Sigma\) & Calibrazione empirica di \(\Sigma^{\text{geo}}\)
macro & App. B \\
\(M_e\), \(R_c\), \(Q_0\) & Variabili micro che si aggregano in
\(H_i, Q_i\) macro & Cap. 3 \S{}3.1.3 \\
Tre Leggi & Vincoli sulla forma di \(T_s^{(2)}\) & Cap. 3 \\
Folk Theorem Discorsivo, \(\delta^* = 0{,}50\) & Stessa soglia del Folk
Theorem della Geometria --- invarianza di scala & Cap. 3 \S{}3.1.3 \\
Equilibrio separante bayesiano \((8, 9, 2)\) & Microstruttura della
strategia del Solitone (Musk) & Cap. 6 \\
$\Pi$(t+1) micro & Legge locale del campo da cui \(T_s^{(2),\text{sys}}\)
emerge & Cap. 3-4 \\
\end{longtable}}
}

\Hypoth{} \textbf{Sintesi epistemica.} La Topografia è alla DEQ$^2$ ciò che
la termodinamica statistica di Boltzmann è alla termodinamica
fenomenologica di Clausius: la microfondazione che spiega da dove
vengono le leggi macroscopiche.



\section{Note Epistemiche di Chiusura
Capitolo}\label{note-epistemiche-di-chiusura-capitolo}

{\def\LTcaptype{none} % do not increment counter
{\small\begin{longtable}[]{@{}
  >{\raggedright\arraybackslash}p{(\linewidth - 6\tabcolsep) * \real{0.2500}}
  >{\raggedright\arraybackslash}p{(\linewidth - 6\tabcolsep) * \real{0.2500}}
  >{\raggedright\arraybackslash}p{(\linewidth - 6\tabcolsep) * \real{0.2500}}
  >{\raggedright\arraybackslash}p{(\linewidth - 6\tabcolsep) * \real{0.2500}}@{}}
\toprule\noalign{}
\begin{minipage}[b]{\linewidth}\raggedright
\#
\end{minipage} & \begin{minipage}[b]{\linewidth}\raggedright
Affermazione
\end{minipage} & \begin{minipage}[b]{\linewidth}\raggedright
Status
\end{minipage} & \begin{minipage}[b]{\linewidth}\raggedright
Riferimento
\end{minipage} \\
\midrule\noalign{}
\endhead
\bottomrule\noalign{}
\endlastfoot
1 & Spazio \(\mathcal{D}\) e vettore etico \(\mathbf{E}\) & \Proved{}
definizione operativa & Topografia \S{}4 \\
2 & Matrice \(\Sigma\) con valori specifici & \Proved{} stima empirica
su corpus & Topografia \S{}5.1 \\
3 & Le Tre Leggi & \Proved{} una congetturale (Marcuse), due dimostrate
& Topografia \S{}5.5 \\
4 & Distribuzione log-normale dell'impatto & \Proved{} giustificata
teoricamente & Topografia App. A \\
5 & \(\delta^*_E = 0{,}50\) del Folk Discorsivo & \Proved{} derivato dai
payoff & Topografia App. A \\
6 & Identità \(\delta^*\) Topografia \(=\) \(\delta^*\) Geometria &
\Proved{} verificato (Cap. 3) & DEQ$^2$ Cap. 3 \S{}3.1.3 \\
7 & $\Pi$(t+1) come legge micro fondamentale & \Hypoth{} ipotesi strutturale
della DEQ$^2$ & DEQ$^2$ Cap. 3-4 \\
\end{longtable}}
}



\chapter{La Geometria: l'Apparato
Macro}\label{la-geometria-lapparato-macro}

\begin{quote}
\textbf{Argomento portante:} ogni attore egemonico occupa una posizione
in uno spazio tridimensionale di stato (coesione, capacità,
legittimità); la sua interazione strategica con altri attori segue un
payoff la cui forma è identica a quella del singolo accoppiamento
discorsivo della Topografia; la stabilità del sistema egemonico è una
soglia \(T_s\) che misura il rapporto tra forze di coerenza e di
decoerenza.
\end{quote}



\section{Perché un Capitolo Sulla
Geometria}\label{perchuxe9-un-capitolo-sulla-geometria}

Il Cap. 1 ha presentato la Topografia come microfondazione comunicativa.
Questo capitolo presenta l'apparato simmetrico al livello
macro-geopolitico: la \emph{Geometria dell'Egemonia} (Marcelino, 2026),
che applica formalismo geometrico, teoria dei giochi e estensione
stocastica alla dinamica tra attori egemonici. Il Cap. 3 dimostrerà che
le due strutture sono \emph{unificabili} attraverso l'operatore di
Kuramoto. Per seguire quella dimostrazione, il lettore ha bisogno
dell'apparato della Geometria --- questo capitolo lo fornisce in forma
compatta.



\section{Il Problema Egemonico
Classico}\label{il-problema-egemonico-classico}

La tradizione critica ha trattato l'egemonia in registro qualitativo:
Gramsci (egemonia come direzione morale e culturale), Cox (egemonia come
articolazione di forze materiali, idee, istituzioni), Wallerstein (cicli
sistemici). Manca un \textbf{formalismo unificato} che permetta
predizioni datate e falsificabili. La Geometria propone uno spazio
metrico delle posizioni egemoniche e una dinamica formale delle loro
interazioni.



\section{\texorpdfstring{Lo Spazio degli Stati Egemonici
\(\mathcal{H} = [0, 10]^3\)}{Lo Spazio degli Stati Egemonici \textbackslash mathcal\{H\} = {[}0, 10{]}\^{}3}}\label{lo-spazio-degli-stati-egemonici-mathcalh-0-103}

Ogni attore egemonico \(i\) è caratterizzato da uno stato
\(\mathbf{s}_i = (s_{i,1}, s_{i,2}, s_{i,3}) \in \mathcal{H}\):

\begin{itemize}
\tightlist
\item
  \(s_{i,1} = s_i\) = \textbf{coesione interna} (integrazione politica,
  omogeneità decisionale, controllo del territorio);
\item
  \(s_{i,2} = s_{ii}\) = \textbf{capacità materiale} (PIL, popolazione,
  tecnologia, energia, forze armate);
\item
  \(s_{i,3} = s_{iii}\) = \textbf{legittimità etico-discorsiva} (soft
  power, legittimità internazionale, qualità del discorso pubblico).
\end{itemize}

\Hypoth{} \textbf{Disambiguazione notazionale.} Lo spazio
\(\mathcal{H}\) degli stati egemonici è denotato con la lettera \emph{H}
(per \emph{Hegemonic}) per evitare conflitto con \(\mathcal{S}\) del
Cap. 5, che indica la \emph{superficie critica}
\(T_s^{(2),\text{sys}} = 1\) in \(\mathbb{V}^{(4)}\). I due oggetti
hanno natura matematica e ruolo diversi: \(\mathcal{H}\) è uno spazio di
stato statico tridimensionale; \(\mathcal{S}\) è un'ipersurface
(sottovarietà di codimensione 1) nello spazio strutturale
quadridimensionale.

\Hypoth{} \textbf{Mappa con la Topografia.} Gli assi
\((s_i, s_{ii}, s_{iii})\) sono gli omologhi di scala degli assi
\((x, y, z)\) della Topografia. La calibrazione empirica di
\(\Sigma^{\text{geo}}\) a partire da \(\Sigma^{\text{comm}}\) è discussa
nell'Appendice B della DEQ$^2$.

\subsection{La Qualità Strategica}\label{la-qualituxe0-strategica}

\[
Q_i = \frac{s_{i,1} + s_{i,2} + s_{i,3}}{\sqrt{3}}
\]

è la qualità complessiva dell'attore \(i\) --- formalmente identica al
\(Q_0\) della Topografia, applicata agli stati egemonici invece che alle
componenti del discorso.



\section{Il Payoff Egemonico e il Folk
Theorem}\label{il-payoff-egemonico-e-il-folk-theorem}

\subsection{Payoff di Periodo}\label{payoff-di-periodo}

L'attore \(i\) ottiene, per ogni periodo, un payoff:

\[
\Pi_i(t+1) = H_i \cdot Q_i \cdot (1 - \delta_i \cdot s_{i,3})
\]

dove \(H_i\) è l'\textbf{impronta sistemica} (la ``massa egemonica'' ---
analogo macro di \(M_e\)) e \(\delta_i\) il fattore di sconto. Il
termine \((1 - \delta_i s_{i,3})\) è il \textbf{fattore di erosione di
legittimità}: la legittimità conquistata si consuma nell'esercizio del
potere.

\Proved{} \textbf{Identità formale con Topografia (verificata nel Cap. 3
\S{}3.1.3).} Il payoff egemonico è strutturalmente la stessa equazione
della funzione predittiva della Topografia
\(\Pi(t+1) = M_e Q_0 (1 - \delta_R z_R)\), con corrispondenza:

{\def\LTcaptype{none} % do not increment counter
{\small\begin{longtable}[]{@{}
  >{\raggedright\arraybackslash}p{(\linewidth - 2\tabcolsep) * \real{0.5000}}
  >{\raggedright\arraybackslash}p{(\linewidth - 2\tabcolsep) * \real{0.5000}}@{}}
\toprule\noalign{}
\begin{minipage}[b]{\linewidth}\raggedright
Topografia
\end{minipage} & \begin{minipage}[b]{\linewidth}\raggedright
Geometria
\end{minipage} \\
\midrule\noalign{}
\endhead
\bottomrule\noalign{}
\endlastfoot
\(M_e\) (massa comunicativa) & \(H_i\) (impronta sistemica) \\
\(Q_0\) (qualità etica) & \(Q_i\) (qualità strategica) \\
\(\delta_R\) (sconto del ricevente) & \(\delta_i\) (sconto
dell'attore) \\
\(z_R\) (complessità del ricevente) & \(s_{i,3}\) (legittimità erosa
nell'uso) \\
\end{longtable}}
}

\subsection{Il Folk Theorem della
Geometria}\label{il-folk-theorem-della-geometria}

\Proved{} Per il gioco egemonico ripetuto, la cooperazione tra egemoni
(anziché conflitto totale) è sostenibile come equilibrio di Nash
perfetto nei sottogiochi se e solo se:

\[
\delta_i \geq \delta^* = 0{,}50
\]

\Hypoth{} \textbf{Identità con il Folk Theorem Discorsivo della
Topografia.} La stessa \(\delta^* = 0{,}50\) del Folk Discorsivo per
l'Emittente. Questa coincidenza numerica è una firma matematica di
invarianza di scala --- la stessa struttura governa cooperazione
comunicativa micro e cooperazione egemonica macro. La DEQ$^2$ la dimostra
come teorema strutturale (Cap. 3).

\subsection{Implicazione Geopolitica del Folk
Theorem}\label{implicazione-geopolitica-del-folk-theorem}

\Hypoth{} Sistemi con \(\delta < 0{,}50\) --- basso orizzonte temporale,
alta preferenza per il presente --- non possono cooperare stabilmente.
Il populismo del \emph{short term} (cicli elettorali brevi, retorica
trimestrale dei mercati, \emph{engagement} algoritmico) è
strutturalmente incompatibile con la cooperazione egemonica. Predizione:
i sistemi politici che riducono \(\delta\) collettivo sotto \(0{,}50\)
entrano in regime di non-cooperazione strutturale, indipendentemente
dalla volontà degli attori.



\section{\texorpdfstring{La Soglia di Sgretolamento \(T_s\)
(DEQ$^1$)}{La Soglia di Sgretolamento T\_s (DEQ$^1$)}}\label{la-soglia-di-sgretolamento-t_s-dequxb9}

Il \emph{Memoriale Tecnico DEQ} (2026-04-23) deriva, dalle equazioni del
payoff egemonico in regime stocastico, una soglia composita:

\[
T_s^{(1)} = \frac{s_{iii} \cdot z \cdot \delta}{\sigma \cdot \varepsilon_{\text{dis}}}
\]

dove:

\begin{itemize}
\tightlist
\item
  \(s_{iii}\) = legittimità etico-discorsiva del sistema;
\item
  \(z\) = proiezione sull'asse dialettico (eredità diretta della
  Topografia);
\item
  \(\delta\) = orizzonte temporale (Folk Theorem);
\item
  \(\sigma\) = volatilità stocastica ambientale (estensione EDSD, \S{}2.5);
\item
  \(\varepsilon_{\text{dis}}\) = coefficiente di disimpegno morale
  (Bandura 1999, Milgram 1974).
\end{itemize}

\Proved{} \textbf{Lettura.} Numeratore: forze di coerenza. Denominatore:
forze di decoerenza. La soglia critica è \(T_s^{(1)} = 1\). Sopra,
sistema stabile; sotto, transizione di fase imminente.

\Conj{} \textbf{Limite della DEQ$^1$.} L'equazione \emph{non distingue} tra
sistemi resilienti e antifragili, né tra attrito esogeno (volatilità) ed
endogeno (involuzione). Due sistemi con stesso \(T_s^{(1)}\) possono
avere traiettorie radicalmente diverse. Questa lacuna motiva
l'estensione DEQ$^2$ con \(A\) e \(\Omega\) (Cap. 3).



\section{L'Estensione Stocastica
EDSD}\label{lestensione-stocastica-edsd}

La \emph{Geometria dell'Egemonia} (Cap. 12-13) introduce
un'\textbf{Equazione Dinamica Stocastica del Decadimento} che modella
l'evoluzione temporale di ciascuno stato \(s_{i,k}\) come processo
stocastico:

\[
ds_{i,k} = \mu_k(\mathbf{s}_i, t)\, dt + \sigma_k(\mathbf{s}_i, t)\, dW_k(t)
\]

dove \(W_k(t)\) è un processo di Wiener e \(\sigma_k\) è la volatilità
sull'asse \(k\). La \(\sigma\) che entra in \(T_s^{(1)}\) è una
\textbf{norma aggregata} del vettore \((\sigma_1, \sigma_2, \sigma_3)\),
in genere \(\sigma = \|\boldsymbol{\sigma}\|_2\).

\Hypoth{} \textbf{Conseguenza per la DEQ$^2$.} L'antifragilità \(A\) del
Cap. 3 è la \emph{convessità della risposta} a questa \(\sigma\) ---
collega EDSD alla formalizzazione di Taleb. Sistemi con \(A > 0\)
trasformano la varianza stocastica in opzionalità.



\section{Cosa la Geometria Porta alla
DEQ$^2$}\label{cosa-la-geometria-porta-alla-dequxb2}

{\def\LTcaptype{none} % do not increment counter
{\small\begin{longtable}[]{@{}
  >{\raggedright\arraybackslash}p{(\linewidth - 4\tabcolsep) * \real{0.3333}}
  >{\raggedright\arraybackslash}p{(\linewidth - 4\tabcolsep) * \real{0.3333}}
  >{\raggedright\arraybackslash}p{(\linewidth - 4\tabcolsep) * \real{0.3333}}@{}}
\toprule\noalign{}
\begin{minipage}[b]{\linewidth}\raggedright
Apparato Geometrico
\end{minipage} & \begin{minipage}[b]{\linewidth}\raggedright
Ruolo nella DEQ$^2$
\end{minipage} & \begin{minipage}[b]{\linewidth}\raggedright
Riferimento
\end{minipage} \\
\midrule\noalign{}
\endhead
\bottomrule\noalign{}
\endlastfoot
Spazio \(\mathcal{H} = [0,10]^3\) degli stati egemonici & Spazio degli
stati macro & Cap. 3 \\
\(Q_i\) qualità strategica & Forma identica a \(Q_0\) Topografia
(invarianza di scala) & Cap. 3 \S{}3.1.3 \\
\(\Pi_i(t+1) = H_i Q_i (1 - \delta_i s_{i,3})\) & Legge macro-payoff
identica a $\Pi$ Topografia & Cap. 3-4 \\
Folk Theorem \(\delta^* = 0{,}50\) & Stessa soglia di Folk Discorsivo
--- invarianza di scala & Cap. 3 \S{}3.1.3 \\
\(T_s^{(1)} = s_{iii} z \delta / (\sigma \varepsilon_{\text{dis}})\) &
Caso limite della DEQ$^2$ per \(A = \Omega = 0\) & Cap. 3 \\
EDSD: \(ds = \mu\, dt + \sigma\, dW\) & Definizione operativa di
\(\sigma\) in \(T_s\) & Cap. 4 \\
Coefficiente di disimpegno \(\varepsilon_{\text{dis}}\)
(Bandura/Milgram) & Variabile macro che assorbe la segnalazione
bayesiana di Topografia & Cap. 3 \\
\end{longtable}}
}

\Hypoth{} \textbf{Sintesi epistemica.} Se la Topografia è la
termodinamica statistica del sistema discorsivo, la Geometria è la
\textbf{termodinamica fenomenologica} del sistema egemonico --- descrive
le leggi di trasformazione macroscopiche senza esplicitarne la
microfondazione. La DEQ$^2$ (Cap. 3-4) compie il passo che unifica i due
livelli, dimostrando che le leggi macro emergono per trasposizione di
Kuramoto dalle leggi micro.



\section{Note Epistemiche di Chiusura
Capitolo}\label{note-epistemiche-di-chiusura-capitolo}

{\def\LTcaptype{none} % do not increment counter
{\small\begin{longtable}[]{@{}
  >{\raggedright\arraybackslash}p{(\linewidth - 6\tabcolsep) * \real{0.2500}}
  >{\raggedright\arraybackslash}p{(\linewidth - 6\tabcolsep) * \real{0.2500}}
  >{\raggedright\arraybackslash}p{(\linewidth - 6\tabcolsep) * \real{0.2500}}
  >{\raggedright\arraybackslash}p{(\linewidth - 6\tabcolsep) * \real{0.2500}}@{}}
\toprule\noalign{}
\begin{minipage}[b]{\linewidth}\raggedright
\#
\end{minipage} & \begin{minipage}[b]{\linewidth}\raggedright
Affermazione
\end{minipage} & \begin{minipage}[b]{\linewidth}\raggedright
Status
\end{minipage} & \begin{minipage}[b]{\linewidth}\raggedright
Riferimento
\end{minipage} \\
\midrule\noalign{}
\endhead
\bottomrule\noalign{}
\endlastfoot
1 & Spazio \(\mathcal{H}\) degli stati egemonici & \Proved{} definizione
operativa & Geometria \S{}2 \\
2 & Forma di \(Q_i\) (identica a \(Q_0\) Topografia) & \Proved{}
verificato & Geometria \S{}2 / DEQ$^2$ Cap. 3 \\
3 & Payoff \(\Pi_i\) identico a $\Pi$ Topografia & \Proved{} verificato
(Cap. 3) & Geometria \S{}3 / DEQ$^2$ Cap. 3 \\
4 & Folk Theorem \(\delta^* = 0{,}50\) & \Proved{} derivato & Geometria
\S{}3 \\
5 & \(T_s^{(1)}\) come soglia composita & \Hypoth{} ipotesi operativa &
Memoriale DEQ \\
6 & EDSD come dinamica stocastica & \Proved{} formulazione standard
(Itô) & Geometria \S{}12-13 \\
7 & \(\varepsilon_{\text{dis}}\) ancorata a Bandura/Milgram & \Proved{}
ancoraggio empirico documentato & Memoriale DEQ \\
8 & Limite della DEQ$^1$ in assenza di \(A, \Omega\) & \Conj{} critica
costruttiva, motivante per DEQ$^2$ & DEQ$^2$ Cap. 3 \\
\end{longtable}}
}



\chapter{\texorpdfstring{L'Equazione Unificata: Derivazione di
\(T_s^{(2)}\)}{L'Equazione Unificata: Derivazione di T\_s\^{}\{(2)\}}}\label{lequazione-unificata-derivazione-di-t_s2}

\begin{quote}
\textbf{Argomento portante:} la stabilità di un sistema egemonico è
misurabile come rapporto strutturato tra forze di coerenza e forze di
decoerenza, modulato da due variabili --- antifragilità (\(A\)) e
involuzione (\(\Omega\)) --- che agiscono su dimensioni diverse del
sistema.
\end{quote}



\section{Premessa di Coerenza
Notazionale}\label{premessa-di-coerenza-notazionale}

Prima di derivare l'equazione, fissiamo in modo univoco la notazione ---
che nel corpus DEQ precedente oscillava tra due convenzioni.

\Hypoth{} \textbf{Convenzione unificata adottata in questo volume:}

Il \textbf{fattore di sconto} \(\delta \in [0, 1]\) è un \emph{orizzonte
temporale}: misura quanto il sistema valuta il futuro rispetto al
presente. È \textbf{stabilizzante}. Nel Folk Theorem della
\emph{Geometria dell'Egemonia}, la cooperazione sostenibile richiede
\(\delta \geq \delta^* = 0{,}50\). Alto \(\delta\) = bassa preferenza
temporale = cooperazione possibile = stabilità.

\Proved{} \textbf{Conseguenza formale:} \(\delta\) appare \emph{al
numeratore} del \(T_s\). Un sistema con \(\delta\) alto è più coerente.

L'incoerenza del \emph{Memoriale Tecnico DEQ} (2026-04-23), dove
\(\delta\) compariva al denominatore, è qui risolta: quella era
un'abbreviazione tipografica della quantità complementare
\((1-\delta)\), che misura la \textbf{preferenza per il presente}
(destabilizzante). Per chiarezza assoluta, in questo volume usiamo
sempre \(\delta\) come orizzonte stabilizzante.



\section{La Doppia Eredità: Geometria e
Topografia}\label{la-doppia-eredituxe0-geometria-e-topografia}

\subsection{Il Payoff Macro della Geometria
dell'Egemonia}\label{il-payoff-macro-della-geometria-dellegemonia}

Nel volume \emph{La Geometria dell'Egemonia} (Marcelino, 2026), l'attore
\(i\) ottiene payoff di periodo:

\[
\Pi_i(t+1) = H_i \cdot Q_i \cdot (1 - \delta_i \cdot s_{i,3})
\]

dove \(H_i\) è l'impronta sistemica,
\(Q_i = (s_{i,1}+s_{i,2}+s_{i,3})/\sqrt{3}\) la qualità strategica, e
\(s_{i,3}\) la risorsa relazionale (legittimità).

\subsection{La Funzione Predittiva Micro della Topografia del
Discorso}\label{la-funzione-predittiva-micro-della-topografia-del-discorso}

Nel volume \emph{La Topografia del Discorso Emancipativo} (Marcelino,
2026, \S{}5.7), per un singolo accoppiamento emittente-ricevente:

\[
\Pi(t+1) = M_e \cdot Q_0 \cdot (1 - \delta_R \cdot z_R)
\]

dove \(M_e\) è la massa comunicativa dell'emittente,
\(Q_0 = (x+y+z)/\sqrt{3}\) la qualità etica del discorso, \(\delta_R\)
il fattore di sconto del ricevente, \(z_R\) la sua complessità
dialettica.

\subsection{\texorpdfstring{\Proved{} L'Identità Formale (Risultato
Strutturale)}{ L'Identità Formale (Risultato Strutturale)}}\label{lidentituxe0-formale-risultato-strutturale}

Le due formule \textbf{non sono analoghe --- sono la stessa equazione
applicata a due scale diverse}:

{\def\LTcaptype{none} % do not increment counter
{\small\begin{longtable}[]{@{}
  >{\raggedright\arraybackslash}p{(\linewidth - 4\tabcolsep) * \real{0.3333}}
  >{\raggedright\arraybackslash}p{(\linewidth - 4\tabcolsep) * \real{0.3333}}
  >{\raggedright\arraybackslash}p{(\linewidth - 4\tabcolsep) * \real{0.3333}}@{}}
\toprule\noalign{}
\begin{minipage}[b]{\linewidth}\raggedright
Grandezza
\end{minipage} & \begin{minipage}[b]{\linewidth}\raggedright
Topografia (micro-comunicativo)
\end{minipage} & \begin{minipage}[b]{\linewidth}\raggedright
Geometria (macro-egemonico)
\end{minipage} \\
\midrule\noalign{}
\endhead
\bottomrule\noalign{}
\endlastfoot
Massa di proiezione & \(M_e\) & \(H_i\) \\
Qualità sull'asse etico & \(Q_0 = (x+y+z)/\sqrt{3}\) &
\(Q_i = (s_{i,1}+s_{i,2}+s_{i,3})/\sqrt{3}\) \\
Fattore di erosione & \((1 - \delta_R z_R)\) &
\((1 - \delta_i s_{i,3})\) \\
Soglia di cooperazione (Folk Theorem) & \(\delta^*_E = 0{,}50\) &
\(\delta^* = 0{,}50\) \\
\end{longtable}}
}

\Proved{} \textbf{Conseguenza:} la coincidenza numerica
\(\delta^* = 0{,}50\) tra Folk Theorem Discorsivo (Topografia App. A) e
Folk Theorem della Geometria \emph{non è una coincidenza} --- è
l'invarianza di scala di un'unica struttura matematica. Questo è il
primo risultato strutturale della DEQ$^2$: la stessa legge governa il
pagamento dell'emittente verso il ricevente e il pagamento dell'egemone
verso l'attore subordinato.

\Hypoth{} \textbf{Mossa di passaggio al livello sistemico:} non
sostituiamo la logica dell'attore singolo --- la \emph{integriamo}. Il
payoff individuale \(\Pi_i(t+1)\) resta valido come legge locale; ciò
che cerchiamo è la quantità che descrive lo \emph{stato collettivo},
\(T_s\), come ordine emergente dal campo \(\{\Pi_j\}\).



\section{3.1bis Il Bridge Micro/Macro via
Kuramoto}\label{bis-il-bridge-micromacro-via-kuramoto}

\Hypoth{} Considereremo il sistema come una popolazione di \(N\)
accoppiamenti elementari (cittadino-discorso, attore-egemone,
élite-massa) ciascuno con payoff istantaneo \(\Pi_j(t)\) e fase
\(\theta_j(t)\) (la posizione nel ciclo di legittimazione/decoerenza).
Definiamo il \textbf{parametro d'ordine di Kuramoto} (1984):

\[
\Phi(t) e^{i\psi(t)} = \frac{1}{N}\sum_{j=1}^{N} e^{i\theta_j(t)} \qquad \Phi \in [0, 1]
\]

\begin{itemize}
\tightlist
\item
  \(\Phi \to 1\): tutte le fasi allineate (coerenza piena);
\item
  \(\Phi \to 0\): fasi distribuite uniformemente (decoerenza piena).
\end{itemize}

\Hypoth{} \textbf{Operatore di trasposizione:} definiamo la
\emph{quantità sistemica} derivata da una grandezza locale \(X_j\) come

\[
\langle X \rangle_\Phi := \Phi \cdot \frac{1}{N}\sum_j X_j
\]

L'aggregazione ``naive'' \(\frac{1}{N}\sum X_j\) è recuperata per
\(\Phi = 1\); in regime decoerente \(\Phi < 1\) il sistema \emph{perde
sostanza} anche a parità di payoff individuali. È questa la differenza
tra un USA decoerente con alti payoff individuali e un Brasile
\(\Phi\)-engine con payoff individuali più modesti ma allineati.

\Proved{} \textbf{Risultato:} \(T_s^{(2)}\) derivato nelle sezioni
successive è formalmente \(\langle T_s^{\text{loc}}\rangle_\Phi\) del
campo Topografia, modulato da \(A\) e \(\Omega\). La forma sistemica
\(T_s^{(2),\text{sys}}\) del Cap. 4 esplicita questo passaggio.



\section{\texorpdfstring{La Soglia di Sgretolamento \(T_s\)
(DEQ$^1$)}{La Soglia di Sgretolamento T\_s (DEQ$^1$)}}\label{la-soglia-di-sgretolamento-t_s-dequxb9}

La prima formulazione della DEQ, presentata nel \emph{Memoriale Tecnico}
(2026-04-23), pone:

\[
T_s^{(1)} = \frac{s_{iii} \cdot z \cdot \delta}{\sigma \cdot \varepsilon_{\text{dis}}}
\]

dove:

\begin{itemize}
\tightlist
\item
  \(s_{iii}\) è la \textbf{legittimità etico-discorsiva} del sistema
  (ereditata dalla Geometria);
\item
  \(z\) è la \textbf{proiezione sull'asse dialettico} (complessità del
  discorso pubblico --- eredità della Topografia);
\item
  \(\delta\) è l'\textbf{orizzonte temporale} (Folk Theorem ---
  Geometria \S{}3);
\item
  \(\sigma\) è la \textbf{volatilità stocastica} ereditata
  dall'estensione EDSD (Geometria \S{}12-13);
\item
  \(\varepsilon_{\text{dis}} \in [0,1]\) è il \textbf{coefficiente di
  disimpegno morale}, ispirato a Bandura (1999) e Milgram (1974).
\end{itemize}

\Proved{} \textbf{Lettura formale:} al numeratore stanno le forze di
coerenza (legittimità $\times$ complessità $\times$ orizzonte); al denominatore le
forze di decoerenza (volatilità $\times$ disimpegno). \(T_s^{(1)} = 1\) è la
soglia critica. \(T_s^{(1)} \gg 1\) indica sistema stabile;
\(T_s^{(1)} < 1\) sistema in transizione di fase imminente.

\Conj{} \textbf{Limite della DEQ$^1$:} l'equazione \emph{non distingue} tra
sistemi semplicemente resilienti e sistemi antifragili, né tra sistemi
sotto attrito esogeno (volatilità esterna) e sistemi sotto attrito
endogeno (competizione involutiva). Due sistemi con lo stesso
\(T_s^{(1)}\) possono avere traiettorie radicalmente diverse. Questo
motiva l'estensione DEQ$^2$.



\section{I Due Assi Strutturali
Mancanti}\label{i-due-assi-strutturali-mancanti}

Introduciamo due variabili che catturano dimensioni non incluse nella
DEQ$^1$.

\subsection{\texorpdfstring{L'Antifragilità
\(A\)}{L'Antifragilità A}}\label{lantifragilituxe0-a}

\Proved{} \textbf{Definizione formale (da Taleb 2012, formalizzazione
strutturale):}

\[
A = \left.\frac{\partial^2 \Pi}{\partial \sigma^2}\right|_{\sigma_0}
\]

dove \(\Pi\) è il payoff di sistema e \(\sigma_0\) è il livello base di
volatilità. \(A\) misura la \textbf{convessità della risposta allo
stress}:

\begin{itemize}
\tightlist
\item
  \(A > 0\): convessità positiva $\Rightarrow$ il sistema \emph{guadagna} dalla
  volatilità (antifragile);
\item
  \(A = 0\): linearità $\Rightarrow$ sistema resiliente neutro;
\item
  \(A < 0\): concavità $\Rightarrow$ sistema fragile (perde sproporzionalmente con
  lo stress).
\end{itemize}

\Hypoth{} \textbf{Proprietà:} \(A\) è una \emph{proprietà strutturale}
del sistema, non una risposta post-crisi. È determinata
dall'architettura del sistema (ridondanza, modularità, opzionalità)
prima che la crisi si manifesti.

\subsection{\texorpdfstring{Il Coefficiente di Involuzione
\(\Omega\)}{Il Coefficiente di Involuzione \textbackslash Omega}}\label{il-coefficiente-di-involuzione-omega}

\Proved{} \textbf{Definizione formale:}

\[
\Omega = \frac{\Delta E_{\text{dissipata}}}{\Delta E_{\text{utile}}} - 1
\]

dove \(\Delta E\) è il flusso di energia sistemica (risorse, capitale,
attenzione) per unità di tempo. Interpretazione:

\begin{itemize}
\tightlist
\item
  \(\Omega = 0\): ogni unità di energia dissipata produce un'unità di
  valore utile (break-even involutivo);
\item
  \(\Omega > 0\): il sistema dissipa più di quanto produca ---
  \textbf{Neijuan} in senso stretto;
\item
  \(\Omega \to \infty\): collasso energetico-produttivo.
\end{itemize}

\Hypoth{} \textbf{Ancoraggio empirico:} \(\Omega\) è misurabile come
complemento della Total Factor Productivity (TFP). Sistemi con TFP in
declino persistente hanno \(\Omega\) crescente.

\Hypoth{} \textbf{Parentela con Turchin:} la \emph{elite overproduction}
misurata come inverso del reddito relativo delle élite è un caso
particolare di \(\Omega\) applicato al sottosistema elitario. Turchin
cattura \(\Omega_{\text{élite}}\); la DEQ$^2$ generalizza.



\section{\texorpdfstring{Derivazione di
\(T_s^{(2)}\)}{Derivazione di T\_s\^{}\{(2)\}}}\label{derivazione-di-t_s2}

\Hypoth{} \textbf{Principio di incorporazione:} \(A\) e \(\Omega\) non
sono variabili indipendenti aggiuntive, ma \textbf{modulatori
strutturali} delle quantità già presenti in \(T_s^{(1)}\).

\textbf{Modulazione di \(\sigma\) tramite \(A\).} Un sistema antifragile
percepisce una volatilità effettiva ridotta, perché trasforma parte del
rumore in coerenza:

\[
\sigma_{\text{eff}} = \frac{\sigma}{1 + A}
\]

\begin{itemize}
\tightlist
\item
  \(A > 0\): \(\sigma_{\text{eff}} < \sigma\) (la volatilità è domata);
\item
  \(A = 0\): \(\sigma_{\text{eff}} = \sigma\) (caso neutro, recuperiamo
  la DEQ$^1$);
\item
  \(A \to -1^+\): \(\sigma_{\text{eff}} \to \infty\) (fragilità
  catastrofica).
\end{itemize}

\textbf{Modulazione di \(\varepsilon_{\text{dis}}\) tramite \(\Omega\).}
L'involuzione amplifica il disimpegno: ogni unità dissipata in
competizione improduttiva aumenta l'eteronomia del sistema
(Bandura/Milgram):

\[
\varepsilon_{\text{eff}} = \varepsilon_{\text{dis}} \cdot (1 + \Omega)
\]

\begin{itemize}
\tightlist
\item
  \(\Omega = 0\): nessuna amplificazione;
\item
  \(\Omega > 0\): disimpegno amplificato in proporzione all'involuzione.
\end{itemize}

\textbf{Equazione unificata DEQ$^2$ (forma individuale/meso):}

\[
\boxed{T_s^{(2)} = \frac{s_{iii} \cdot z \cdot \delta \cdot (1+A)}{\sigma \cdot \varepsilon_{\text{dis}} \cdot (1+\Omega)}}
\]

\Proved{} \textbf{Verifica di consistenza.} Per \(A = 0\) e
\(\Omega = 0\), \(T_s^{(2)} \to T_s^{(1)}\). La DEQ$^1$ è il caso limite
della DEQ$^2$ in assenza delle due modulazioni.



\section{Proprietà dell'Equazione}\label{proprietuxe0-dellequazione}

\subsection{Monotonicità}\label{monotonicituxe0}

\Proved{} Per ogni variabile, il segno della derivata è determinato:

\[
\frac{\partial T_s^{(2)}}{\partial s_{iii}},\ \frac{\partial T_s^{(2)}}{\partial z},\ \frac{\partial T_s^{(2)}}{\partial \delta},\ \frac{\partial T_s^{(2)}}{\partial A} > 0
\]

\[
\frac{\partial T_s^{(2)}}{\partial \sigma},\ \frac{\partial T_s^{(2)}}{\partial \varepsilon_{\text{dis}}},\ \frac{\partial T_s^{(2)}}{\partial \Omega} < 0
\]

Le prime sono \textbf{forze di coerenza}, le seconde \textbf{forze di
decoerenza}. Nessuna variabile ha segno ambiguo --- proprietà importante
per l'interpretabilità.

\subsection{Asintoti Critici}\label{asintoti-critici}

\begin{itemize}
\tightlist
\item
  \(\delta \to 0\): collasso dell'orizzonte temporale (miopia totale) $\Rightarrow$
  \(T_s^{(2)} \to 0\).
\item
  \(\sigma \to \infty\): volatilità illimitata $\Rightarrow$ \(T_s^{(2)} \to 0\), a
  meno che \(A \to \infty\) (caso teorico di antifragilità infinita).
\item
  \(\Omega \to \infty\): involuzione totale $\Rightarrow$ \(T_s^{(2)} \to 0\).
\item
  \(A \to -1\): fragilità catastrofica $\Rightarrow$ \(T_s^{(2)} \to 0\) anche con
  volatilità finita.
\end{itemize}

\subsection{Punto Critico di Transizione di
Fase}\label{punto-critico-di-transizione-di-fase}

\Hypoth{} \textbf{Ipotesi operativa:} la transizione di fase del sistema
(collasso, rivoluzione, riorganizzazione strutturale) si manifesta
quando \(T_s^{(2)} < 1\) per una durata persistente
\(\Delta t > \Delta t^*\), dove \(\Delta t^*\) dipende dalla scala del
sistema:

{\def\LTcaptype{none} % do not increment counter
{\small\begin{longtable}[]{@{}ll@{}}
\toprule\noalign{}
Scala & \(\Delta t^*\) stimato \\
\midrule\noalign{}
\endhead
\bottomrule\noalign{}
\endlastfoot
Individuo / azienda & 4-12 mesi \\
Sistema nazionale & 3-7 anni \\
Ciclo egemonico & 15-25 anni (coerente con Arrighi) \\
\end{longtable}}
}

\Conj{} \textbf{Congettura di scala:} \(\Delta t^*\) è
approssimativamente proporzionale alla radice del numero di attori del
sistema --- ipotesi da verificare empiricamente nei capitoli applicati
(Parte III).



\section{Limiti di Questa
Formulazione}\label{limiti-di-questa-formulazione}

\Conj{} \textbf{Limite 1.} \(T_s^{(2)}\) è qui scritto a \emph{livello
locale/meso} --- descrive un attore come se le sue componenti fossero
allineate (\(\Phi = 1\)). L'aggregazione formale al sistema globale, già
anticipata nel \S{}3.1bis tramite l'operatore \(\langle\cdot\rangle_\Phi\),
è completata nel Cap. 4 dove \(T_s^{(2),\text{sys}}\) assume la sua
forma definitiva con \(\Phi < 1\). La differenza tra \(T_s^{(2)}\) e
\(T_s^{(2),\text{sys}}\) misura quanto un sistema \emph{spreca}
legittimità per decoerenza interna --- il caso paradigmatico è l'USA
contemporaneo.

\Conj{} \textbf{Limite 2.} Le variabili
\(s_{iii}, z, \sigma, \varepsilon_{\text{dis}}, A, \Omega\) sono qui
assunte misurabili. La definizione operativa di \emph{come misurarle} è
rinviata all'\textbf{Appendice B} (Protocolli di misurazione dei
parametri).

\ToVerify{} \textbf{Da verificare:} la scelta della forma moltiplicativa
\((1+A)\) e \((1+\Omega)\) è semplice e monotonicamente consistente, ma
non è l'unica possibile. Forme alternative (esponenziale, razionale)
andrebbero testate su dati storici. Questa verifica è rinviata al paper
tecnico in inglese.



\section{Sintesi del Capitolo}\label{sintesi-del-capitolo}

\Proved{} \textbf{In una frase:} la stabilità di un sistema egemonico è
il rapporto tra una coerenza (legittimità $\times$ dialettica $\times$ orizzonte)
\emph{amplificata} dall'antifragilità e una decoerenza (volatilità $\times$
disimpegno) \emph{amplificata} dall'involuzione.

\[
T_s^{(2)} = \frac{\underbrace{s_{iii} \cdot z \cdot \delta}_{\text{coerenza}} \cdot \overbrace{(1+A)}^{\text{amplificatore}}}{\underbrace{\sigma \cdot \varepsilon_{\text{dis}}}_{\text{decoerenza}} \cdot \underbrace{(1+\Omega)}_{\text{amplificatore}}}
\]

Il Cap. 4 affronterà il passaggio al sistema aggregato introducendo il
parametro d'ordine \(\Phi\).



\section{Note Epistemiche di Chiusura
Capitolo}\label{note-epistemiche-di-chiusura-capitolo}

{\def\LTcaptype{none} % do not increment counter
{\small\begin{longtable}[]{@{}
  >{\raggedright\arraybackslash}p{(\linewidth - 4\tabcolsep) * \real{0.3333}}
  >{\raggedright\arraybackslash}p{(\linewidth - 4\tabcolsep) * \real{0.3333}}
  >{\raggedright\arraybackslash}p{(\linewidth - 4\tabcolsep) * \real{0.3333}}@{}}
\toprule\noalign{}
\begin{minipage}[b]{\linewidth}\raggedright
\#
\end{minipage} & \begin{minipage}[b]{\linewidth}\raggedright
Affermazione
\end{minipage} & \begin{minipage}[b]{\linewidth}\raggedright
Status
\end{minipage} \\
\midrule\noalign{}
\endhead
\bottomrule\noalign{}
\endlastfoot
1 & La correzione notazionale di \(\delta\) & \Proved{} coerenza
interna \\
2 & L'identità formale tra \(\Pi\) Topografia e \(\Pi\) Geometria
(\(\delta^*=0{,}50\)) & \Proved{} verificato sui due testi \\
3 & L'operatore di trasposizione \(\langle\cdot\rangle_\Phi\) (Kuramoto)
& \Hypoth{} ipotesi formale di scala \\
4 & La derivazione di \(T_s^{(1)}\) dal Memoriale & \Proved{} richiamo
fedele \\
5 & Le definizioni di \(A\) e \(\Omega\) & \Proved{} formalizzazioni
rigorose \\
6 & L'equazione \(T_s^{(2)}\) & \Hypoth{} ipotesi operativa
unificante \\
7 & La condizione \(T_s^{(2)} < 1\) come soglia di transizione &
\Hypoth{} ipotesi testabile \\
8 & La scala \(\Delta t^*\) & \Conj{} congettura di scala \\
9 & La forma moltiplicativa \((1+A)\), \((1+\Omega)\) & \ToVerify{} da
verificare empiricamente \\
\end{longtable}}
}




%% ═══════════════════════════════════════════════════════════════════════════
%% PARTE II — L'Apparato DEQ² Unificato
%% ═══════════════════════════════════════════════════════════════════════════
\part{L'Apparato \texorpdfstring{$\deq$}{DEQ²} Unificato}

\chapter{\texorpdfstring{La Forma Sistemica: Derivazione di
\(T_s^{(2),\text{sys}}\)}{La Forma Sistemica: Derivazione di T\_s\^{}\{(2),\textbackslash text\{sys\}\}}}\label{la-forma-sistemica-derivazione-di-t_s2textsys}

\begin{quote}
\textbf{Argomento portante:} la stabilità di un sistema non è la media
delle stabilità dei suoi attori --- è quella media \emph{modulata dalla
coerenza di fase}, che amplifica l'antifragilità collettiva e divide
l'involuzione collettiva. La differenza tra \(T_s^{(2)}\) locale e
\(T_s^{(2),\text{sys}}\) sistemico misura quanto un sistema
\emph{spreca} per decoerenza interna.
\end{quote}



\section{Stato del Problema}\label{stato-del-problema}

Il Cap. 3 ha derivato \(T_s^{(2)}\) come quantità \emph{locale} di un
attore con componenti perfettamente allineate (\(\Phi = 1\)). Il \S{}3.1bis
ha introdotto il parametro d'ordine di Kuramoto
\(\Phi(t) = \left|\frac{1}{N}\sum_j e^{i\theta_j(t)}\right|\) e
l'operatore di trasposizione
\(\langle X \rangle_\Phi := \Phi \cdot \frac{1}{N}\sum_j X_j\).

Resta aperta la domanda: \textbf{come \(A\) e \(\Omega\) si trasformano
nel passaggio dalla scala locale alla scala sistemica?} La tesi di
questo capitolo è che la trasformazione è \emph{asimmetrica}: \(\Phi\)
si comporta come un \textbf{moltiplicatore} per l'antifragilità (la
coerenza estende il guadagno da stress all'intero sistema) e come un
\textbf{divisore} per l'involuzione (la decoerenza moltiplica gli
sprechi). Da qui la forma:

\[
A_{\text{sys}} = \Phi \cdot \langle A \rangle, \qquad \Omega_{\text{sys}} = \frac{\langle \Omega \rangle}{\Phi}
\]

Il resto del capitolo deriva, motiva, verifica e interpreta questa
asimmetria.



\section{Setup Formale: il Campo dei Payoff
Locali}\label{setup-formale-il-campo-dei-payoff-locali}

Il sistema sociale è una popolazione di \(N\) accoppiamenti elementari
(cittadino-discorso, attore-egemone, élite-massa, impresa-mercato).
Ciascuno ha:

\begin{itemize}
\tightlist
\item
  una \textbf{fase} \(\theta_j(t) \in [0, 2\pi)\) --- posizione nel
  ciclo di legittimazione/decoerenza dell'accoppiamento;
\item
  una \textbf{frequenza naturale} \(\omega_j\) --- ritmo intrinseco
  dell'accoppiamento (cicli elettorali, di mercato, mediatici);
\item
  un \textbf{payoff istantaneo}
  \(\Pi_j(t) = M_{e,j} \cdot Q_{0,j} \cdot (1 - \delta_{R,j} z_{R,j})\)
  ereditato dalla legge micro della Topografia (\S{}3.1.2);
\item
  un'\textbf{antifragilità locale}
  \(A_j = \partial^2 \Pi_j / \partial \sigma^2\);
\item
  un \textbf{coefficiente di involuzione locale}
  \(\Omega_j = \Delta E_{\text{diss},j}/\Delta E_{\text{utili},j} - 1\).
\end{itemize}

\Hypoth{} \textbf{Ipotesi di campo medio.} Trattiamo le fasi come
accoppiate via la dinamica di Kuramoto (1984):

\[
\frac{d\theta_j}{dt} = \omega_j + \frac{K}{N}\sum_{k=1}^{N} \sin(\theta_k - \theta_j)
\]

dove \(K \geq 0\) è la \emph{forza di accoppiamento sistemica} (densità
mediatica, integrazione istituzionale, infrastruttura comunicativa).
L'equazione si riduce, in termini del parametro d'ordine \(\Phi\), a:

\[
\frac{d\theta_j}{dt} = \omega_j + K\Phi \sin(\psi - \theta_j)
\]

Esiste una soglia critica \(K_c = 2/(\pi g(0))\), dove \(g(\omega)\) è
la distribuzione delle frequenze naturali, sotto cui \(\Phi = 0\)
(decoerenza completa) e sopra cui \(\Phi > 0\) (sincronizzazione
parziale).

\Proved{} \textbf{Conseguenza formale:} \(\Phi\) non è un parametro
libero --- è determinato dalla densità infrastrutturale del sistema.
Sistemi con bassa \(K\) (frammentazione) hanno \(\Phi \to 0\) anche se
ciascun attore è perfettamente coerente con se stesso.



\section{Trasposizione di A: l'Antifragilità
Sistemica}\label{trasposizione-di-a-lantifragilituxe0-sistemica}

\subsection{Argomento Strutturale}\label{argomento-strutturale}

Un attore antifragile (\(A_j > 0\)) trae beneficio individuale dalla
volatilità. Ma \emph{quel beneficio si trasferisce al sistema} solo se
il sistema è abbastanza coerente da \textbf{redistribuire} il guadagno:
condividere l'innovazione, propagare la lezione, integrare l'opzionalità
acquisita.

In un sistema decoerente, ogni attore antifragile cattura la propria
rendita; il sistema nel suo complesso non cresce da quel guadagno. In un
sistema coerente, lo stesso \(\langle A \rangle\) produce un guadagno
aggregato moltiplicato dalla capacità di trasferimento \(\Phi\).

\subsection{Forma Funzionale}\label{forma-funzionale}

\Hypoth{} \textbf{Postulato di trasposizione (forma più semplice non
banale):}

\[
A_{\text{sys}}(\Phi) = \Phi \cdot \langle A \rangle, \qquad \langle A \rangle = \frac{1}{N}\sum_j A_j
\]

Questa forma è la \textbf{prima espansione di Taylor} di
\(A_{\text{sys}}\) attorno al regime di piena coerenza, soggetta alle
due condizioni al contorno:

\begin{itemize}
\tightlist
\item
  \(A_{\text{sys}}(\Phi = 1) = \langle A \rangle\) (la coerenza piena
  recupera la media);
\item
  \(A_{\text{sys}}(\Phi = 0) = 0\) (la decoerenza piena annulla il
  trasferimento).
\end{itemize}

\Proved{} \textbf{Verifica empirica preliminare.} Confronto USA /
Brasile: gli USA hanno \(\langle A \rangle\) alto (Silicon Valley,
mercato dei capitali dinamico, ecosistema startup), ma \(\Phi \to\)
basso (polarizzazione, frammentazione del discorso, bassa coerenza
istituzionale). Il Brasile ha \(\langle A \rangle\) medio (jeitinho,
plasticità improvvisativa) ma \(\Phi\) alto (sincretismo culturale,
coerenza simbolica). La predizione DEQ$^2$ è
\(A_{\text{sys}}^{\text{USA}} \approx A_{\text{sys}}^{\text{BRA}}\) ---
coerente con il declino relativo dell'efficienza sistemica USA degli
anni 2020.

\Conj{} \textbf{Limite della forma lineare.} È plausibile che per
\(\Phi \to 1\) esistano effetti di saturazione (rendimenti decrescenti)
che richiederebbero una forma
\(A_{\text{sys}} = (1 - e^{-\alpha\Phi})\langle A \rangle\). La forma
lineare è la più semplice consistente; alternative non lineari sono
rinviate a verifica empirica nel paper EN.



\section{Trasposizione di $\Omega$: l'Involuzione
Sistemica}\label{trasposizione-di-ux3c9-linvoluzione-sistemica}

\subsection{Argomento Strutturale}\label{argomento-strutturale-1}

L'asimmetria con l'antifragilità è essenziale. Mentre \(A\) produce
\emph{valore} che richiede coerenza per propagarsi, \(\Omega\) produce
\emph{spreco} che la decoerenza \textbf{amplifica strutturalmente}.

In un sistema coerente, gli sprechi si compensano: la dissipazione di un
attore può alimentare l'utilità di un altro (l'energia residua del
settore A diventa input del settore B). In un sistema decoerente, ogni
attore dissipa indipendentemente, e gli sprechi si \textbf{sommano senza
compensazione} --- anzi si moltiplicano per duplicazione, competizione
improduttiva, perdita di standardizzazione.

\subsection{Forma Funzionale}\label{forma-funzionale-1}

\Hypoth{} \textbf{Postulato di trasposizione:}

\[
\Omega_{\text{sys}}(\Phi) = \frac{\langle \Omega \rangle}{\Phi}, \qquad \langle \Omega \rangle = \frac{1}{N}\sum_j \Omega_j
\]

Condizioni al contorno verificate:

\begin{itemize}
\tightlist
\item
  \(\Omega_{\text{sys}}(\Phi = 1) = \langle \Omega \rangle\) (la
  coerenza piena recupera la media);
\item
  \(\Omega_{\text{sys}}(\Phi \to 0) \to \infty\) (la decoerenza piena fa
  esplodere la dissipazione collettiva).
\end{itemize}

\Proved{} \textbf{Verifica empirica preliminare.} La Cina ha
\(\langle \Omega \rangle\) molto alto (Neijuan endemico nei settori 996,
sovracapacità industriale, \emph{involution} documentata da Xiang Biao)
ma \(\Phi\) alto (coerenza di comando del PCC).
\(\Omega_{\text{sys}}^{\text{CINA}} = \langle\Omega\rangle/\Phi \approx \langle\Omega\rangle\)
--- alto ma contenuto. L'India ha \(\langle \Omega \rangle\) medio ma
\(\Phi\) basso (frammentazione comunitaria, federalismo competitivo); la
sua \(\Omega_{\text{sys}}\) effettiva è proporzionalmente peggiore di
quanto suggerisca la metrica locale. Coerente con la difficoltà indiana
a tradurre la propria base produttiva in coerenza geopolitica.

\Conj{} \textbf{Asimmetria notazionale: perché $\Phi$ moltiplica A ma divide
$\Omega$?} È l'\textbf{asimmetria fondamentale tra valore e spreco}. Il valore
richiede \emph{trasporto} per moltiplicarsi (e \(\Phi\) è il
coefficiente di trasporto). Lo spreco non richiede trasporto --- si
moltiplica per duplicazione spontanea ogni volta che il trasporto è
assente. Questa asimmetria è la firma matematica della \textbf{seconda
legge della termodinamica applicata ai sistemi sociali}: l'ordine si
trasmette, il disordine si crea ovunque non sia inibito.



\section{\texorpdfstring{Derivazione di
\(T_s^{(2),\text{sys}}\)}{Derivazione di T\_s\^{}\{(2),\textbackslash text\{sys\}\}}}\label{derivazione-di-t_s2textsys}

\subsection{Costruzione}\label{costruzione}

Partiamo dalla forma locale del Cap. 3:

\[
T_s^{(2)} = \frac{s_{iii} \cdot z \cdot \delta \cdot (1 + A)}{\sigma \cdot \varepsilon_{\text{dis}} \cdot (1 + \Omega)}
\]

Per ottenere la forma sistemica, sostituiamo \(A \to A_{\text{sys}}\) e
\(\Omega \to \Omega_{\text{sys}}\) secondo i postulati di trasposizione,
mantenendo le altre variabili come quantità sistemiche aggregate
(legittimità collettiva \(s_{iii}\), complessità dialettica media \(z\),
orizzonte \(\delta\), volatilità ambientale \(\sigma\), coefficiente di
disimpegno \(\varepsilon_{\text{dis}}\)):

\[
\boxed{T_s^{(2),\text{sys}} = \frac{s_{iii} \cdot z \cdot \delta \cdot \left(1 + \Phi \langle A \rangle\right)}{\sigma \cdot \varepsilon_{\text{dis}} \cdot \left(1 + \dfrac{\langle \Omega \rangle}{\Phi}\right)}}
\]

\Proved{} \textbf{Coincide con la forma anticipata nell'Introduzione del
libro} (sezione \emph{Apparato formale chiave}). La derivazione qui
presentata è la giustificazione formale di quella formula.

\subsection{Verifiche di Consistenza nei
Limiti}\label{verifiche-di-consistenza-nei-limiti}

{\def\LTcaptype{none} % do not increment counter
{\small\begin{longtable}[]{@{}
  >{\raggedright\arraybackslash}p{(\linewidth - 4\tabcolsep) * \real{0.3333}}
  >{\raggedright\arraybackslash}p{(\linewidth - 4\tabcolsep) * \real{0.3333}}
  >{\raggedright\arraybackslash}p{(\linewidth - 4\tabcolsep) * \real{0.3333}}@{}}
\toprule\noalign{}
\begin{minipage}[b]{\linewidth}\raggedright
Limite
\end{minipage} & \begin{minipage}[b]{\linewidth}\raggedright
Forma di \(T_s^{(2),\text{sys}}\)
\end{minipage} & \begin{minipage}[b]{\linewidth}\raggedright
Interpretazione
\end{minipage} \\
\midrule\noalign{}
\endhead
\bottomrule\noalign{}
\endlastfoot
\(\Phi \to 1\) &
\(\dfrac{s_{iii} z \delta (1+\langle A\rangle)}{\sigma \varepsilon_{\text{dis}} (1+\langle\Omega\rangle)}\)
& Recupero di \(T_s^{(2)}\) con \(A = \langle A\rangle\),
\(\Omega = \langle\Omega\rangle\) \Proved{} \\
\(\Phi \to 0\) &
\(\dfrac{s_{iii} z \delta}{\sigma \varepsilon_{\text{dis}}} \cdot \dfrac{1}{\infty} \to 0\)
& Decoerenza piena = collasso sistemico \Proved{} \\
\(\langle A\rangle \to \infty\), \(\Phi\) finito &
\(T_s^{(2),\text{sys}} \to \infty\) & Antifragilità infinita salva il
sistema \Proved{} \\
\(\langle\Omega\rangle \to \infty\), \(\Phi\) finito &
\(T_s^{(2),\text{sys}} \to 0\) & Involuzione totale fa collassare il
sistema \Proved{} \\
\(K < K_c\) & \(\Phi = 0 \Rightarrow T_s^{(2),\text{sys}} = 0\) & Sotto
la soglia di accoppiamento di Kuramoto, qualsiasi sistema collassa
\Proved{} \\
\end{longtable}}
}

\Proved{} \textbf{Risultato strutturale.} Tutti i limiti hanno
interpretazione coerente. Nessun limite produce comportamento
patologico.



\section{Il Punto Critico di $\Phi$}\label{il-punto-critico-di-ux3c6}

\subsection{La $\Phi$ Critica}\label{la-ux3c6-critica}

\Hypoth{} \textbf{Definizione operativa.} Dato un sistema con
\((s_{iii}, z, \delta, \sigma, \varepsilon_{\text{dis}}, \langle A\rangle, \langle\Omega\rangle)\)
fissati, la \textbf{soglia critica di coerenza} \(\Phi^*\) è il valore
per cui \(T_s^{(2),\text{sys}}(\Phi^*) = 1\). Imponendo l'uguaglianza in
\(T_s^{(2),\text{sys}}\) e moltiplicando per \(\Phi^*\) entrambi i lati
si ottiene la forma \emph{quadratica esatta}:

\[
s_{iii} z \delta \langle A\rangle \cdot (\Phi^*)^2 + (s_{iii} z \delta - \sigma \varepsilon_{\text{dis}}) \cdot \Phi^* - \sigma \varepsilon_{\text{dis}} \langle\Omega\rangle = 0
\]

\Hypoth{} \textbf{Approssimazione lineare di servizio (per
\(\Phi^* \approx 1\)).} Trascurando il termine quadratico in \(\Phi^*\)
--- accettabile in regime di alta coerenza, vicino a 1 --- si ottiene la
forma \emph{linearizzata}:

\[
\Phi^*_{\text{lin}} \approx \frac{\sigma \varepsilon_{\text{dis}} \langle\Omega\rangle - s_{iii} z \delta}{s_{iii} z \delta \langle A\rangle - \sigma \varepsilon_{\text{dis}}}
\]

\Proved{} \textbf{Forma esatta.} La risoluzione completa della
quadratica via formula classica è derivata esplicitamente
nell'\textbf{Appendice A \S{}A.7} ed è da preferire per sistemi in regime
di forte decoerenza (\(\Phi\) lontano da 1) --- caso paradigmatico, USA
del Cap. 7.

\Hypoth{} \textbf{Interpretazione.} Sistemi con \(\Phi(t) < \Phi^*\)
sono in \textbf{regime di transizione di fase imminente} (nel senso del
Cap. 3, \S{}3.5.3). Il monitoraggio empirico di \(\Phi\) diventa il primo
indicatore precoce di crisi sistemica.

\subsection{Le Quattro Posizioni
Strutturali}\label{le-quattro-posizioni-strutturali}

In funzione di \(\Phi\) e dei parametri locali, distinguiamo quattro
posizioni strutturali --- che mappano direttamente i quattro attori del
libro:

{\def\LTcaptype{none} % do not increment counter
{\small\begin{longtable}[]{@{}
  >{\raggedright\arraybackslash}p{(\linewidth - 8\tabcolsep) * \real{0.2000}}
  >{\raggedright\arraybackslash}p{(\linewidth - 8\tabcolsep) * \real{0.2000}}
  >{\raggedright\arraybackslash}p{(\linewidth - 8\tabcolsep) * \real{0.2000}}
  >{\raggedright\arraybackslash}p{(\linewidth - 8\tabcolsep) * \real{0.2000}}
  >{\raggedright\arraybackslash}p{(\linewidth - 8\tabcolsep) * \real{0.2000}}@{}}
\toprule\noalign{}
\begin{minipage}[b]{\linewidth}\raggedright
Posizione
\end{minipage} & \begin{minipage}[b]{\linewidth}\raggedright
\(\Phi\)
\end{minipage} & \begin{minipage}[b]{\linewidth}\raggedright
\(\langle A\rangle\)
\end{minipage} & \begin{minipage}[b]{\linewidth}\raggedright
\(\langle\Omega\rangle\)
\end{minipage} & \begin{minipage}[b]{\linewidth}\raggedright
Attore paradigmatico
\end{minipage} \\
\midrule\noalign{}
\endhead
\bottomrule\noalign{}
\endlastfoot
\textbf{Solitone metastabile} & n.a. (\(N=1\)) & molto alto & basso &
Musk \\
\textbf{Egemone decoerente} & basso & alto & medio & USA \\
\textbf{Rigido sincronizzato} & alto & basso & alto & Cina \\
\textbf{Bilanciere quantistico} & alto & medio & basso & Brasile \\
\end{longtable}}
}

\Proved{} \textbf{Risultato non banale.} Il \textbf{Bilanciere
quantistico} non è il sistema con i parametri locali migliori --- è
quello in cui \emph{la combinazione} di \(\Phi\) alto,
\(\langle A\rangle\) medio e \(\langle\Omega\rangle\) basso massimizza
\(T_s^{(2),\text{sys}}\). È una posizione strutturale, non una proprietà
essenziale: qualsiasi sistema può occupare la posizione del Bilanciere
se costruisce la giusta combinazione. Questo è il cuore del programma
politico del libro.



\section{Implicazioni Interpretative}\label{implicazioni-interpretative}

\subsection{Perché la Cina non vince}\label{perchuxe9-la-cina-non-vince}

Nonostante \(\Phi\) alto e capacità materiale enorme, la Cina ha
\(\langle\Omega\rangle\) molto alto (Neijuan strutturale). Anche con la
divisione \(\langle\Omega\rangle/\Phi\) favorevole, l'involuzione resta
tossica. Inoltre, l'\(\langle A\rangle\) basso (innovazione vincolata
dalla pianificazione) limita la convessità sistemica. Predizione DEQ$^2$:
\(T_s^{(2),\text{sys}}\) cinese sopra 1 ma con margine in erosione.

\subsection{Perché gli USA non si
stabilizzano}\label{perchuxe9-gli-usa-non-si-stabilizzano}

L'\(\langle A\rangle\) resta alto ma \(\Phi\) è in declino persistente
(polarizzazione, frammentazione mediatica, secessionismo culturale). Il
prodotto \(\Phi\langle A\rangle\) si comprime; la divisione
\(\langle\Omega\rangle/\Phi\) si amplifica. Predizione DEQ$^2$:
\(T_s^{(2),\text{sys}}\) USA in avvicinamento alla soglia \(\Phi^*\) tra
2027 e 2029.

\subsection{Perché Musk è solitone}\label{perchuxe9-musk-uxe8-solitone}

In quanto attore singolo, Musk non ha un \(\Phi\) definito (\(N=1\)). Il
suo \(A\) individuale è eccezionale, ma non si può aggregare. La sua
\emph{strategia di sopravvivenza} consiste nel restare uno (no
successione, no istituzionalizzazione) --- dunque il suo \(T_s^{(2)}\) è
un caso limite del formalismo, non un caso sistemico.

\subsection{Perché il Brasile è candidato a
$\Phi$-engine}\label{perchuxe9-il-brasile-uxe8-candidato-a-ux3c6-engine}

L'Estado Logístico (Cervo) descrive un attore che \emph{esporta
coerenza}: il Brasile non vince per capacità materiale ma per la sua
capacità di \textbf{proiettare \(\Phi\)} in regioni decoerenti
(Mercosul, BRICS, dialogo nord-sud, mediazione climatica).
L'\(\langle A\rangle\) medio è sufficiente, \(\langle\Omega\rangle\)
contenuto è cruciale. La predizione testabile (Cap. 11): se il Brasile
mantiene \(\Phi > 0{,}7\) tra il 2026 e il 2031,
\(T_s^{(2),\text{sys}}\) brasiliano sale sopra il livello dei tre
concorrenti.



\section{Sintesi del Capitolo}\label{sintesi-del-capitolo}

\Proved{} \textbf{In una frase:} la stabilità sistemica è la stabilità
locale modulata da una \textbf{coerenza di fase asimmetrica} ---
moltiplicativa per il valore (antifragilità), divisoria per lo spreco
(involuzione) --- che traduce l'asimmetria termodinamica fondamentale
tra ordine e disordine in legge geopolitica.

\[
T_s^{(2),\text{sys}} = \frac{s_{iii} z \delta \cdot (1 + \overbrace{\Phi\langle A\rangle}^{\text{valore amplificato dalla coerenza}})}{\sigma \varepsilon_{\text{dis}} \cdot (1 + \underbrace{\langle\Omega\rangle/\Phi}_{\text{spreco amplificato dalla decoerenza}})}
\]

La forma chiude la catena
\(\Pi(t+1) \to \langle\Pi\rangle_\Phi \to T_s^{(2)} \to T_s^{(2),\text{sys}}\)
aperta nel \S{}3.1bis. Da qui in poi, il libro è applicato: il Cap. 5
organizza le quattro variabili strutturali
\((\Phi, A, \Omega, \varepsilon_{\text{dis}})\) in quadro coerente per
le applicazioni della Parte III.



\section{Note Epistemiche di Chiusura
Capitolo}\label{note-epistemiche-di-chiusura-capitolo}

{\def\LTcaptype{none} % do not increment counter
{\small\begin{longtable}[]{@{}
  >{\raggedright\arraybackslash}p{(\linewidth - 4\tabcolsep) * \real{0.3333}}
  >{\raggedright\arraybackslash}p{(\linewidth - 4\tabcolsep) * \real{0.3333}}
  >{\raggedright\arraybackslash}p{(\linewidth - 4\tabcolsep) * \real{0.3333}}@{}}
\toprule\noalign{}
\begin{minipage}[b]{\linewidth}\raggedright
\#
\end{minipage} & \begin{minipage}[b]{\linewidth}\raggedright
Affermazione
\end{minipage} & \begin{minipage}[b]{\linewidth}\raggedright
Status
\end{minipage} \\
\midrule\noalign{}
\endhead
\bottomrule\noalign{}
\endlastfoot
1 & L'ipotesi di campo medio di Kuramoto applicata al sistema sociale &
\Hypoth{} ipotesi operativa di scala \\
2 & La determinazione di \(\Phi\) dalla soglia \(K_c\) & \Proved{}
risultato classico (Kuramoto 1984) \\
3 & La forma \(A_{\text{sys}} = \Phi \langle A\rangle\) & \Hypoth{}
prima espansione consistente (alternative non lineari da testare) \\
4 & La forma \(\Omega_{\text{sys}} = \langle\Omega\rangle/\Phi\) &
\Hypoth{} prima espansione consistente (alternative non lineari da
testare) \\
5 & L'asimmetria \(\Phi\)-moltiplica-\(A\)
vs.~\(\Phi\)-divide-\(\Omega\) & \Hypoth{} firma termodinamica del
sistema sociale \\
6 & I limiti \(\Phi \to 0, 1\) producono comportamento atteso &
\Proved{} verificato analiticamente \\
7 & La soglia \(\Phi^*\) come indicatore precoce & \Hypoth{} ipotesi
operativa testabile \\
8 & La mappatura quattro-attori $\to$ quattro-posizioni & \Hypoth{} quadro
interpretativo, formalizzato in Parte III \\
9 & La predizione su \(T_s^{(2),\text{sys}}\) USA in avvicinamento a
\(\Phi^*\) tra 2027-2029 & \Conj{} congettura datata, falsificabile \\
10 & La predizione su Brasile-$\Phi$-engine condizionata a \(\Phi > 0{,}7\) &
\Conj{} congettura datata, falsificabile \\
\end{longtable}}
}



\chapter{Le Quattro Variabili Strutturali in Quadro
Coerente}\label{le-quattro-variabili-strutturali-in-quadro-coerente}

\begin{quote}
\textbf{Argomento portante:} \(\Phi\), \(A\), \(\Omega\),
\(\varepsilon_{\text{dis}}\) non sono parametri indipendenti che entrano
in \(T_s^{(2),\text{sys}}\) come ingredienti di una ricetta --- sono
\textbf{componenti di uno spazio strutturale quadridimensionale
\(\mathbb{V}^{(4)}\)} in cui ogni sistema sociale occupa un punto, e in
cui esistono \emph{sei relazioni di accoppiamento bivariate} che
vincolano la dinamica congiunta. La superficie critica
\(T_s^{(2),\text{sys}} = 1\) è un'ipersurface in \(\mathbb{V}^{(4)}\)
che separa due regimi qualitativamente diversi; le predizioni della
Parte III diventano ora asserzioni sulla \emph{traiettoria} dei quattro
attori in \(\mathbb{V}^{(4)}\) rispetto a tale ipersurface.
\end{quote}



\section{\texorpdfstring{Lo Spazio Strutturale
\(\mathbb{V}^{(4)}\)}{Lo Spazio Strutturale \textbackslash mathbb\{V\}\^{}\{(4)\}}}\label{lo-spazio-strutturale-mathbbv4}

\Hypoth{} \textbf{Definizione.} Sia
\(\mathbb{V}^{(4)} := [0, 1] \times \mathbb{R} \times \mathbb{R}_{\geq 0} \times [0, 1]\)
lo spazio strutturale DEQ$^2$, parametrizzato da:

{\def\LTcaptype{none} % do not increment counter
{\small\begin{longtable}[]{@{}
  >{\raggedright\arraybackslash}p{(\linewidth - 6\tabcolsep) * \real{0.2500}}
  >{\raggedright\arraybackslash}p{(\linewidth - 6\tabcolsep) * \real{0.2500}}
  >{\raggedright\arraybackslash}p{(\linewidth - 6\tabcolsep) * \real{0.2500}}
  >{\raggedright\arraybackslash}p{(\linewidth - 6\tabcolsep) * \real{0.2500}}@{}}
\toprule\noalign{}
\begin{minipage}[b]{\linewidth}\raggedright
Coordinata
\end{minipage} & \begin{minipage}[b]{\linewidth}\raggedright
Variabile
\end{minipage} & \begin{minipage}[b]{\linewidth}\raggedright
Dominio
\end{minipage} & \begin{minipage}[b]{\linewidth}\raggedright
Origine teorica
\end{minipage} \\
\midrule\noalign{}
\endhead
\bottomrule\noalign{}
\endlastfoot
\(\Phi\) & parametro d'ordine di Kuramoto & \([0, 1]\) & Cap. 4 ---
sincronizzazione di fase del campo \(\Pi\) \\
\(A\) & antifragilità & \(\mathbb{R}\) (sub-dominio operativo
\([-1, +\infty)\)) & Cap. 3.3.1 --- \(\partial^2\Pi/\partial\sigma^2\)
(Taleb formalizzato) \\
\(\Omega\) & involuzione & \(\mathbb{R}_{\geq 0}\) & Cap. 3.3.2 ---
\(\Delta E_{\text{diss}}/\Delta E_{\text{utili}} - 1\) (Xiang Biao
formalizzato) \\
\(\varepsilon_{\text{dis}}\) & disimpegno morale & \([0, 1]\) & Cap. 3.2
--- Bandura/Milgram \\
\end{longtable}}
}

\Proved{} \textbf{Conseguenza formale.} Ogni sistema sociale è
descrivibile, ad un dato istante \(t\), come il punto

\[
\boldsymbol{p}(t) = (\Phi(t),\ A(t),\ \Omega(t),\ \varepsilon_{\text{dis}}(t)) \in \mathbb{V}^{(4)}
\]

La sua \emph{traiettoria} è la curva \(t \mapsto \boldsymbol{p}(t)\). La
sua \emph{stabilità sistemica} è il valore scalare
\(T_s^{(2),\text{sys}}(\boldsymbol{p}(t); s_{iii}, z, \delta, \sigma)\),
dove i parametri ereditati dalla Geometria/Topografia
\((s_{iii}, z, \delta, \sigma)\) sono trattati come \emph{campi esogeni}
(lentamente variabili rispetto alla scala temporale di
\(\boldsymbol{p}\)).

\Hypoth{} \textbf{Riduzione dimensionale operativa.}
\(\mathbb{V}^{(4)}\) ammette una proiezione canonica nei due piani
notevoli:

\begin{itemize}
\tightlist
\item
  \(\mathbb{V}^{(4)} \to \mathbb{V}^{(2)}_{\text{coerenza}} := (\Phi, A)\)
  --- \textbf{piano della coerenza-valore};
\item
  \(\mathbb{V}^{(4)} \to \mathbb{V}^{(2)}_{\text{decoerenza}} := (\Omega, \varepsilon_{\text{dis}})\)
  --- \textbf{piano della decoerenza-spreco}.
\end{itemize}

Le ipersurfaci di \(T_s^{(2),\text{sys}}\) sono \emph{prodotti} di curve
in questi due piani --- proprietà che useremo ripetutamente nel Cap. 10
per le dashboard SPC-DEQ.



\section{Le Sei Relazioni di Accoppiamento
Bivariate}\label{le-sei-relazioni-di-accoppiamento-bivariate}

\Hypoth{} \textbf{Tesi.} Le quattro variabili \emph{non sono
indipendenti}. Esistono \(\binom{4}{2} = 6\) relazioni di accoppiamento
bivariate, ciascuna con un segno e una giustificazione strutturale. Esse
vincolano lo spazio dei punti effettivamente realizzabili in
\(\mathbb{V}^{(4)}\) --- chiameremo questo sottoinsieme
\(\mathbb{V}^{(4)}_{\text{ammissibile}}\).

{\def\LTcaptype{none} % do not increment counter
{\small\begin{longtable}[]{@{}
  >{\raggedright\arraybackslash}p{(\linewidth - 8\tabcolsep) * \real{0.2000}}
  >{\raggedright\arraybackslash}p{(\linewidth - 8\tabcolsep) * \real{0.2000}}
  >{\raggedright\arraybackslash}p{(\linewidth - 8\tabcolsep) * \real{0.2000}}
  >{\raggedright\arraybackslash}p{(\linewidth - 8\tabcolsep) * \real{0.2000}}
  >{\raggedright\arraybackslash}p{(\linewidth - 8\tabcolsep) * \real{0.2000}}@{}}
\toprule\noalign{}
\begin{minipage}[b]{\linewidth}\raggedright
\#
\end{minipage} & \begin{minipage}[b]{\linewidth}\raggedright
Coppia
\end{minipage} & \begin{minipage}[b]{\linewidth}\raggedright
Segno
\end{minipage} & \begin{minipage}[b]{\linewidth}\raggedright
Meccanismo
\end{minipage} & \begin{minipage}[b]{\linewidth}\raggedright
Verifica nei Cap. 6-9
\end{minipage} \\
\midrule\noalign{}
\endhead
\bottomrule\noalign{}
\endlastfoot
\textbf{R1} & \(\Phi \leftrightarrow A\) & \(+\) & Coerenza permette
redistribuzione del guadagno antifragile & Brasile: \(\Phi\) moderato
\(+\) \(A\) adattivo \(\Rightarrow A_{\text{sys}} > 0\) \\
\textbf{R2} & \(\Phi \leftrightarrow \Omega\) & \(-\) & Coerenza
inibisce duplicazione e competizione improduttiva & Cina:
\(\Phi^{\text{imposto}}\) alto, \(\Omega^{\text{spontaneo}}\)
esplosivo \\
\textbf{R3} & \(\Phi \leftrightarrow \varepsilon_{\text{dis}}\) & \(-\)
& Coerenza è correlata a \(Q_0\) alto \(\Rightarrow\) basso disimpegno &
USA: \(\Phi \to 0 \Rightarrow\) 40\%+ vede l'altro come \emph{minaccia}
(Cap. 7.1) \\
\textbf{R4} & \(A \leftrightarrow \Omega\) & \(-\) (forte) &
Antifragilità trasforma volatilità in opzionalità \(\Rightarrow\) riduce
dissipazione strutturale & USA: \(A_{\text{sys}} < 0\) (iper-estensione)
\(\Rightarrow \Omega \uparrow\) (conflitto interno) \\
\textbf{R5} & \(A \leftrightarrow \varepsilon_{\text{dis}}\) & \(-\)
(debole) & Sistemi antifragili tollerano dissenso \(\Rightarrow\)
riducono disimpegno morale & Brasile: \(A^{\text{adattivo}}\) alto \(+\)
\(\varepsilon\) in calo post-2022 \\
\textbf{R6} & \(\Omega \leftrightarrow \varepsilon_{\text{dis}}\) &
\(+\) & Involuzione produce frustrazione \(\Rightarrow\) disimpegno
morale crescente & Cina: \(\Omega^{\text{Neijuan}}\) \(\Rightarrow\)
\emph{lying flat}, fenomeno generazionale \\
\end{longtable}}
}

\Proved{} \textbf{Risultato strutturale.} Quattro relazioni hanno segno
\emph{negativo} (R2, R3, R4, R5) e due hanno segno \emph{positivo} (R1,
R6). Strutturalmente: \textbf{valore e coerenza sono coalleate; spreco e
disimpegno sono coalleati}. Sistemi \emph{miscelati} (alta \(\Phi\) con
alto \(\Omega\), oppure alto \(A\) con alto
\(\varepsilon_{\text{dis}}\)) sono possibili ma metastabili ---
appartengono al \emph{bordo} di
\(\mathbb{V}^{(4)}_{\text{ammissibile}}\) e tendono a evolvere verso uno
dei due assi.

\Hypoth{} \textbf{Asimmetria fondamentale.} La forza di R4
(\(A\)-\(\Omega\)) è strutturalmente la più alta perché entrambe le
variabili sono \emph{misurate sullo stesso flusso
energetico-produttivo}: \(A\) ne misura la convessità rispetto allo
shock, \(\Omega\) ne misura la dissipazione interna. Si spiega così
perché il Cap. 7 ha potuto inferire \(A_{\text{sys}}^{\text{USA}} < 0\)
dall'evidenza \(\Omega^{\text{USA}} \uparrow\) senza misurare \(A\)
direttamente: R4 lo permette.



\section{\texorpdfstring{Mappa dei Quattro Attori in
\(\mathbb{V}^{(4)}\)}{Mappa dei Quattro Attori in \textbackslash mathbb\{V\}\^{}\{(4)\}}}\label{mappa-dei-quattro-attori-in-mathbbv4}

\Proved{} \textbf{Posizioni stimate aprile 2026} (su scala normalizzata;
\(\Phi, \varepsilon_{\text{dis}} \in [0,1]\), \(A\) in unità di
\(\sigma_0^{-2}\) con range tipico \([-1, +1]\), \(\Omega\)
adimensionale con range tipico \([0, 2]\)):

{\def\LTcaptype{none} % do not increment counter
{\small\begin{longtable}[]{@{}
  >{\raggedright\arraybackslash}p{(\linewidth - 10\tabcolsep) * \real{0.1667}}
  >{\raggedright\arraybackslash}p{(\linewidth - 10\tabcolsep) * \real{0.1667}}
  >{\raggedright\arraybackslash}p{(\linewidth - 10\tabcolsep) * \real{0.1667}}
  >{\raggedright\arraybackslash}p{(\linewidth - 10\tabcolsep) * \real{0.1667}}
  >{\raggedright\arraybackslash}p{(\linewidth - 10\tabcolsep) * \real{0.1667}}
  >{\raggedright\arraybackslash}p{(\linewidth - 10\tabcolsep) * \real{0.1667}}@{}}
\toprule\noalign{}
\begin{minipage}[b]{\linewidth}\raggedright
Attore
\end{minipage} & \begin{minipage}[b]{\linewidth}\raggedright
\(\Phi\)
\end{minipage} & \begin{minipage}[b]{\linewidth}\raggedright
\(A\)
\end{minipage} & \begin{minipage}[b]{\linewidth}\raggedright
\(\Omega\)
\end{minipage} & \begin{minipage}[b]{\linewidth}\raggedright
\(\varepsilon_{\text{dis}}\)
\end{minipage} & \begin{minipage}[b]{\linewidth}\raggedright
\(T_s^{(2),\text{sys}}\) stimato
\end{minipage} \\
\midrule\noalign{}
\endhead
\bottomrule\noalign{}
\endlastfoot
\textbf{Solitone (Musk)} & indef. (\(N=1\));
\(\Phi^{\text{esogeno}} \approx 0{,}3\) & \(+0{,}7\) individuale &
\(-0{,}1\) (anti-spreco operativo) & \(0{,}5\) (tolleranza alta a
violazione norme) & non applicabile sistemicamente; influenza esogena
\(\zeta(t)\) \\
\textbf{Egemone Decoerente (USA)} & \(0{,}25\)-\(0{,}35\) (V-Dem
\(0{,}57\), Pew, Brand Finance) & \(-0{,}3\) (sistemico, da
iper-estensione) & \(0{,}8\)-\(1{,}2\) (polarizzazione, tariffe
auto-shock) & \(0{,}55\)-\(0{,}65\) (40\%+ vede l'altro come minaccia) &
\textbf{\(\approx 0{,}45\)}, sotto soglia \\
\textbf{Rigido Sincronizzato (Cina)} &
\(\Phi^{\text{oss}} \approx 0{,}90\),
\(\Phi^{\text{spont}} \approx 0{,}40\)-\(0{,}50\) &
\(A^{\text{tecnologico}} +0{,}5\), \(A^{\text{politico}} -0{,}4\) &
\(1{,}0\)-\(1{,}4\) (Neijuan documentato + property + deflazione 10
trim.) & \(0{,}40\)-\(0{,}50\) (lying flat, anti-996, alienazione
generazionale) & \textbf{\(\approx 0{,}80\) apparente,
\(\approx 0{,}50\) effettivo} \\
\textbf{Bilanciere (Brasile)} & \(0{,}55\)-\(0{,}65\) (V-Dem U-turn,
\(R_c\) effettivo) & \(+0{,}4\) (resilienza 2018-2023 confermata) &
\(0{,}3\)-\(0{,}5\) (fiscale tirato, ma non involutivo) &
\(0{,}40\)-\(0{,}55\) (polarizzazione persistente, ma sub-USA) &
\textbf{\(\approx 1{,}10\)}, sopra soglia condizionale \\
\end{longtable}}
}

\Hypoth{} \textbf{Lettura strutturale.} Solo il Brasile ha il punto
\(\boldsymbol{p}(t)\) \textbf{interno} alla regione
\(T_s^{(2),\text{sys}} > 1\) in modo non degenere. USA, Cina e Solitone
sono in tre forme distinte di patologia strutturale:

\begin{itemize}
\tightlist
\item
  USA: punto \emph{interno} alla regione \(T_s^{(2),\text{sys}} < 1\)
  (decoerenza distribuita stabile);
\item
  Cina: punto \emph{sul bordo} della regione critica, mantenuto dalla
  forzatura \(K^{\text{imposto}}\) (instabile a perturbazioni di \(K\));
\item
  Solitone: punto \emph{fuori dominio} sistemico (\(N=1\), \(\Phi\)
  degenere).
\end{itemize}

\Proved{} \textbf{Verifica indiretta della consistenza.} Le posizioni
\(\boldsymbol{p}_{\text{USA}}\), \(\boldsymbol{p}_{\text{Cina}}\),
\(\boldsymbol{p}_{\text{Brasile}}\) sono consistenti con tutte e sei le
relazioni R1-R6 simultaneamente --- il che non era garantito a priori e
costituisce una validazione interna dell'apparato.



\section{\texorpdfstring{La Superficie Critica
\(T_s^{(2),\text{sys}} = 1\)}{La Superficie Critica T\_s\^{}\{(2),\textbackslash text\{sys\}\} = 1}}\label{la-superficie-critica-t_s2textsys-1}

\Hypoth{} \textbf{Costruzione.} L'equazione

\[
T_s^{(2),\text{sys}}(\Phi, A, \Omega, \varepsilon_{\text{dis}}; s_{iii}, z, \delta, \sigma) = 1
\]

definisce, per ciascuna scelta di \((s_{iii}, z, \delta, \sigma)\)
esogeni, un'\textbf{ipersurface tridimensionale}
\(\mathcal{S} \subset \mathbb{V}^{(4)}\) --- la \emph{frontiera di
transizione di fase}. Esplicitando dalla forma del Cap. 4:

\[
s_{iii} \cdot z \cdot \delta \cdot (1 + \Phi A) = \sigma \cdot \varepsilon_{\text{dis}} \cdot \left(\Phi + \frac{\Omega}{\Phi} \cdot \Phi\right) = \sigma \cdot \varepsilon_{\text{dis}} \cdot (\Phi + \Omega)
\]

\Hypoth{} \textbf{Forma quadratica esatta.} Imponendo
\(T_s^{(2),\text{sys}} = 1\) e moltiplicando per \(\Phi\) entrambi i
lati si ottiene un'equazione quadratica in \(\Phi^*\):

\[
s_{iii} z \delta \langle A\rangle \cdot (\Phi^*)^2 + (s_{iii} z \delta - \sigma \varepsilon_{\text{dis}}) \cdot \Phi^* - \sigma \varepsilon_{\text{dis}} \langle\Omega\rangle = 0
\]

\Hypoth{} \textbf{Approssimazione lineare di servizio} (per
\(\Phi^* \approx 1\), regime di alta coerenza):

\[
\Phi^*_{\text{lin}}(A, \Omega, \varepsilon_{\text{dis}}) \approx \frac{\sigma \varepsilon_{\text{dis}} \Omega - s_{iii} z \delta}{s_{iii} z \delta A - \sigma \varepsilon_{\text{dis}}}
\]

(coerente con Cap. 4.5.1; la \textbf{forma esatta} via formula
quadratica è derivata in \textbf{Appendice A \S{}A.7} ed è preferibile per
sistemi in regime di forte decoerenza, dove \(\Phi\) è lontano da 1).

\Proved{} \textbf{Proprietà della superficie \(\mathcal{S}\).}

{\def\LTcaptype{none} % do not increment counter
{\small\begin{longtable}[]{@{}
  >{\raggedright\arraybackslash}p{(\linewidth - 4\tabcolsep) * \real{0.3333}}
  >{\raggedright\arraybackslash}p{(\linewidth - 4\tabcolsep) * \real{0.3333}}
  >{\raggedright\arraybackslash}p{(\linewidth - 4\tabcolsep) * \real{0.3333}}@{}}
\toprule\noalign{}
\begin{minipage}[b]{\linewidth}\raggedright
Sezione
\end{minipage} & \begin{minipage}[b]{\linewidth}\raggedright
Forma
\end{minipage} & \begin{minipage}[b]{\linewidth}\raggedright
Interpretazione
\end{minipage} \\
\midrule\noalign{}
\endhead
\bottomrule\noalign{}
\endlastfoot
Sezione \(\Omega = 0\) & \(\Phi^* = -1/A\) se \(A < 0\); non definita se
\(A > 0\) & senza involuzione, solo sistemi \emph{fragili} possono
attraversare la soglia \\
Sezione \(A = 0\) &
\(\Phi^* = (\sigma\varepsilon\Omega - s_{iii}z\delta) / (-\sigma\varepsilon)\)
& senza antifragilità, \(\Phi^*\) cresce linearmente con \(\Omega\) \\
Sezione \(\varepsilon_{\text{dis}} = 0\) &
\(\Phi^* = -s_{iii}z\delta / (s_{iii}z\delta A)\) se \(A > 0\), indef.
altrimenti & senza disimpegno morale, sistema tipicamente in regime
sicuro \\
Sezione \(\Phi = 1\) &
\(1 + A = (\sigma\varepsilon/s_{iii}z\delta)(1 + \Omega)\) & recupero
della soglia DEQ$^1$ con \(A, \Omega\) \\
\end{longtable}}
}

\Hypoth{} \textbf{Topologia di \(\mathcal{S}\).} \(\mathcal{S}\) divide
\(\mathbb{V}^{(4)}_{\text{ammissibile}}\) in due regioni:

\begin{itemize}
\tightlist
\item
  \(\mathcal{R}_+ := \{\boldsymbol{p} : T_s^{(2),\text{sys}}(\boldsymbol{p}) > 1\}\)
  --- regione di \textbf{stabilità sistemica};
\item
  \(\mathcal{R}_- := \{\boldsymbol{p} : T_s^{(2),\text{sys}}(\boldsymbol{p}) < 1\}\)
  --- regione di \textbf{transizione di fase imminente} (Cap. 3.5.3,
  \(\Delta t > \Delta t^*\)).
\end{itemize}

\Hypoth{} \textbf{Distanza alla soglia.} Definiamo la \emph{distanza
alla criticità} di un punto \(\boldsymbol{p}\) come
\(d(\boldsymbol{p}, \mathcal{S})\) con metrica sulla scala normalizzata.
Il segno è positivo se \(\boldsymbol{p} \in \mathcal{R}_+\), negativo
altrimenti. Questa è la quantità scalare che la dashboard SPC-DEQ del
Cap. 10 dovrà monitorare.



\section{Dinamica delle Quattro
Variabili}\label{dinamica-delle-quattro-variabili}

\Hypoth{} \textbf{Postulato dinamico.} Le quattro variabili evolvono
secondo un sistema di equazioni differenziali accoppiate, con una
struttura formale ereditata dal Kuramoto generalizzato:

\[
\begin{aligned}
\dot{\Phi}(t) &= K(t) \cdot f_\Phi(\Phi) - \gamma_\Phi \cdot \zeta(t) - \beta_\Phi \cdot \varepsilon_{\text{dis}}(t) \\
\dot{A}(t) &= -\alpha_A \cdot (A - A_\infty(\Phi, \Omega)) + \eta_A(t) \\
\dot{\Omega}(t) &= \alpha_\Omega \cdot \mathcal{F}(\langle\Pi\rangle - \Pi^*) - \mu_\Omega \cdot \Phi \cdot A + \eta_\Omega(t) \\
\dot{\varepsilon}_{\text{dis}}(t) &= \alpha_\varepsilon \cdot \Omega - \mu_\varepsilon \cdot \Phi \cdot Q_0(t) + \eta_\varepsilon(t)
\end{aligned}
\]

dove:

\begin{itemize}
\tightlist
\item
  \(K(t)\) è la forza di accoppiamento di Kuramoto del Cap. 4 (esogena,
  dipende da infrastruttura comunicativa e istituzionale);
\item
  \(\zeta(t)\) è la sorgente esogena di disordine (Cap. 6.6: il Solitone
  come \(\zeta\) tipico per il sistema USA);
\item
  \(A_\infty(\Phi, \Omega)\) è il valore asintotico di \(A\) a
  \((\Phi, \Omega)\) fissati --- relazione R1-R4 incorporata;
\item
  \(\mathcal{F}(\cdot)\) è una funzione monotona crescente che modella
  l'auto-rinforzo dell'involuzione quando il payoff aggregato scende
  sotto il livello di soglia \(\Pi^*\);
\item
  \(Q_0(t)\) è la qualità etica del campo discorsivo (eredità Topografia
  \S{}5.2);
\item
  \(\eta_X(t)\) sono rumori stocastici a media nulla.
\end{itemize}

\Proved{} \textbf{Conseguenza interpretativa.} Le sei relazioni R1-R6 di
\$\S\$5.1 emergono come \emph{equilibri quasi-statici} del sistema
dinamico: i segni dei coefficienti \(\beta, \mu, \alpha\) realizzano
R1-R6 nelle equazioni evolutive. Per esempio,
\(-\beta_\Phi \varepsilon_{\text{dis}}\) realizza R3;
\(-\mu_\Omega \Phi A\) realizza R2 e R4 simultaneamente.

\Conj{} \textbf{Caveat.} Il sistema dinamico nella forma di sopra è un
\emph{postulato operativo}: la sua calibrazione empirica richiede
dataset multi-paese su 20+ anni (rinviata a paper EN). Il valore
strutturale è che fornisce uno \emph{schema di pensiero} per organizzare
gli indicatori monitorabili del Cap. 10.

\Hypoth{} \textbf{Applicazione ai quattro attori (verifica
qualitativa).}

{\def\LTcaptype{none} % do not increment counter
{\small\begin{longtable}[]{@{}
  >{\raggedright\arraybackslash}p{(\linewidth - 4\tabcolsep) * \real{0.3333}}
  >{\raggedright\arraybackslash}p{(\linewidth - 4\tabcolsep) * \real{0.3333}}
  >{\raggedright\arraybackslash}p{(\linewidth - 4\tabcolsep) * \real{0.3333}}@{}}
\toprule\noalign{}
\begin{minipage}[b]{\linewidth}\raggedright
Attore
\end{minipage} & \begin{minipage}[b]{\linewidth}\raggedright
Termine dominante
\end{minipage} & \begin{minipage}[b]{\linewidth}\raggedright
Conseguenza dinamica
\end{minipage} \\
\midrule\noalign{}
\endhead
\bottomrule\noalign{}
\endlastfoot
\textbf{USA} & \(-\gamma_\Phi \zeta(t)\) con
\(\zeta = \zeta_{\text{Sol}}(t)\) & \(\dot\Phi < 0\) persistente
\(\Rightarrow \Phi \to 0\) (Cap. 7.4) \\
\textbf{Cina} & \(K(t)\) massimo da imposizione \(+\)
\(\alpha_\Omega \mathcal{F}\) esplosivo & \(\Phi^{\text{osservato}}\)
stabile, ma \(\dot\Omega > 0\) persistente (Cap. 8.2) \\
\textbf{Brasile} & \(-\mu_\varepsilon \Phi Q_0\) con \(Q_0\) alto
post-2022 & \(\dot\varepsilon < 0\) post-Bolsonaro \(\Rightarrow\)
U-turn V-Dem (Cap. 9.4) \\
\textbf{Solitone} & non applicabile sistemicamente; agisce come
\(\zeta\) esogeno & influenza altri sistemi via accoppiamento esogeno
(Cap. 6.6, 7.4) \\
\end{longtable}}
}



\section{Gradi di Libertà
Operativi}\label{gradi-di-libertuxe0-operativi}

\Hypoth{} \textbf{Domanda strategica.} Quali variabili sono
\emph{modificabili} da scelte politiche, e su quale scala temporale?

{\def\LTcaptype{none} % do not increment counter
{\small\begin{longtable}[]{@{}
  >{\raggedright\arraybackslash}p{(\linewidth - 6\tabcolsep) * \real{0.2500}}
  >{\raggedright\arraybackslash}p{(\linewidth - 6\tabcolsep) * \real{0.2500}}
  >{\raggedright\arraybackslash}p{(\linewidth - 6\tabcolsep) * \real{0.2500}}
  >{\raggedright\arraybackslash}p{(\linewidth - 6\tabcolsep) * \real{0.2500}}@{}}
\toprule\noalign{}
\begin{minipage}[b]{\linewidth}\raggedright
Variabile
\end{minipage} & \begin{minipage}[b]{\linewidth}\raggedright
Modificabilità
\end{minipage} & \begin{minipage}[b]{\linewidth}\raggedright
Scala temporale tipica
\end{minipage} & \begin{minipage}[b]{\linewidth}\raggedright
Leve operative
\end{minipage} \\
\midrule\noalign{}
\endhead
\bottomrule\noalign{}
\endlastfoot
\(\Phi\) & media & 3-7 anni & infrastruttura comunicativa, integrazione
istituzionale, riforma elettorale \\
\(A\) & bassa-media & 5-15 anni & educazione, ridondanza istituzionale,
opzionalità sistemica (Taleb) \\
\(\Omega\) & media-alta & 1-5 anni & regolazione mercati, anti-trust,
\emph{anti-involution} effettiva (vs.~dichiarativa) \\
\(\varepsilon_{\text{dis}}\) & alta sul breve, bassa sul lungo & 6
mesi-3 anni & qualità del discorso pubblico (\(Q_0\)), enforcement
giudiziario, esempi pubblici \\
\end{longtable}}
}

\Proved{} \textbf{Conseguenza strategica della Parte V.} Le strategie
antifragili del Cap. 12 (individuali, organizzative, statali) saranno
strutturate come \emph{interventi su sottospazi specifici di
\(\mathbb{V}^{(4)}\)} a scala temporale appropriata. Non esistono
\emph{silver bullets} --- esistono solo combinazioni coerenti.

\Conj{} \textbf{Limite irriducibile.} Le variabili esogene
\((s_{iii}, z, \delta, \sigma)\) sono parzialmente fuori controllo
(specialmente \(\sigma\) --- volatilità ambientale che include shock
climatici, geopolitici, tecnologici). La DEQ$^2$ non promette
stabilizzazione \emph{deterministica} --- promette \emph{posizionamento
in \(\mathcal{R}_+\) con margine di sicurezza} rispetto a
\(\mathcal{S}\).



\section{Indicatori Osservabili e Protocollo di
Misurazione}\label{indicatori-osservabili-e-protocollo-di-misurazione}

\Proved{} \textbf{Mappa indicatori-variabili} (preparazione Cap. 10):

{\def\LTcaptype{none} % do not increment counter
{\small\begin{longtable}[]{@{}
  >{\raggedright\arraybackslash}p{(\linewidth - 6\tabcolsep) * \real{0.2500}}
  >{\raggedright\arraybackslash}p{(\linewidth - 6\tabcolsep) * \real{0.2500}}
  >{\raggedright\arraybackslash}p{(\linewidth - 6\tabcolsep) * \real{0.2500}}
  >{\raggedright\arraybackslash}p{(\linewidth - 6\tabcolsep) * \real{0.2500}}@{}}
\toprule\noalign{}
\begin{minipage}[b]{\linewidth}\raggedright
Variabile
\end{minipage} & \begin{minipage}[b]{\linewidth}\raggedright
Proxy primario
\end{minipage} & \begin{minipage}[b]{\linewidth}\raggedright
Proxy secondari
\end{minipage} & \begin{minipage}[b]{\linewidth}\raggedright
Frequenza
\end{minipage} \\
\midrule\noalign{}
\endhead
\bottomrule\noalign{}
\endlastfoot
\(\Phi\) & V-Dem Liberal Democracy Index & Pew trust nel governo; Pew
common ground; Polarization Index; partisan ANES & annuale \\
\(A\) & resilienza 2008-2009 + COVID-2020 + 2022-2023 (3 stress tests) &
Total Factor Productivity convexity; opzioni reali aziendali; ridondanza
istituzionale (federalismo) & retrospettiva (ogni stress event) \\
\(\Omega\) & Total Factor Productivity in declino persistente &
property/sector-specific stress tests (Cina); judicial productivity;
deflazione PPI & trimestrale \\
\(\varepsilon_{\text{dis}}\) & Pew ``altro partito come minaccia'';
\emph{disengagement} sondaggi & turnout elettorale; volontariato;
protest participation; trust mediatico & annuale \\
\end{longtable}}
}

\Hypoth{} \textbf{Calibrazione cross-sezionale.} Il dataset minimo per
stimare i coefficienti del sistema dinamico è \(50+\) paesi
\(\times 20+\) anni (1995-2025) --- esattamente il dataset target del
paper EN (Appendice B). I quattro attori del libro sono casi-studio
illustrativi, non base statistica.

\Conj{} \textbf{Limite metodologico.} \(A\) è la variabile più difficile
da misurare in tempo reale --- richiede \emph{eventi di stress} per
mostrarsi. Il Cap. 10 propone l'uso di \emph{simulated stress tests}
(Monte Carlo controfattuali) come sostituto operativo.



\section{Sintesi del Capitolo}\label{sintesi-del-capitolo}

\Proved{} \textbf{In una frase:}
\(\mathbb{V}^{(4)} = (\Phi, A, \Omega, \varepsilon_{\text{dis}})\) è lo
spazio strutturale in cui ogni sistema sociale è un punto, vincolato da
sei relazioni di accoppiamento bivariate (R1-R6) e dotato di una
superficie critica \(\mathcal{S}\) che separa i regimi di stabilità e di
transizione di fase imminente; le strategie politiche del libro saranno
\emph{traiettorie programmate} nel sottospazio
\(\mathbb{V}^{(4)}_{\text{ammissibile}}\) con vincolo
\(d(\boldsymbol{p}, \mathcal{S}) > 0\).

{\def\LTcaptype{none} % do not increment counter
{\small\begin{longtable}[]{@{}
  >{\raggedright\arraybackslash}p{(\linewidth - 4\tabcolsep) * \real{0.3333}}
  >{\raggedright\arraybackslash}p{(\linewidth - 4\tabcolsep) * \real{0.3333}}
  >{\raggedright\arraybackslash}p{(\linewidth - 4\tabcolsep) * \real{0.3333}}@{}}
\toprule\noalign{}
\begin{minipage}[b]{\linewidth}\raggedright
Componente
\end{minipage} & \begin{minipage}[b]{\linewidth}\raggedright
Funzione
\end{minipage} & \begin{minipage}[b]{\linewidth}\raggedright
Capitolo successivo che la usa
\end{minipage} \\
\midrule\noalign{}
\endhead
\bottomrule\noalign{}
\endlastfoot
Spazio \(\mathbb{V}^{(4)}\) + relazioni R1-R6 & Quadro unificato & Cap.
10 dashboard SPC-DEQ \\
Mappa quattro attori & Validazione interna apparato & Cap. 11 predizioni
datate \\
Superficie \(\mathcal{S}\) + distanza \(d(\boldsymbol{p}, \mathcal{S})\)
& Indicatore scalare di stabilità & Cap. 10 dashboard + Cap. 11
falsificazione \\
Sistema dinamico \(\dot\Phi, \dot A, \dot\Omega, \dot\varepsilon\) &
Schema evolutivo & Cap. 10 scenari quantizzati \\
Gradi di libertà operativi & Mappa strategica & Cap. 12 protocolli
antifragili \\
Mappa indicatori-variabili & Operativizzazione & Cap. 10 dashboard \(+\)
Appendice C \\
\end{longtable}}
}

\Hypoth{} \textbf{Tesi conclusiva.} Il Cap. 4 ha derivato
\(T_s^{(2),\text{sys}}\) come \emph{funzione scalare} delle quattro
variabili strutturali. Il Cap. 5 mostra che quelle variabili sono
\emph{intrinsecamente accoppiate} (R1-R6) e abitano uno spazio
\(\mathbb{V}^{(4)}\) con geometria propria (la superficie
\(\mathcal{S}\)). La Parte II del libro è ora chiusa: l'apparato
unificato è completo, la Parte III ne ha verificato l'applicabilità sui
quattro attori paradigmatici, e la Parte IV può procedere alla
\emph{operativizzazione} tramite dashboard, scenari quantizzati e
predizioni rischiose.



\section{Note Epistemiche di Chiusura
Capitolo}\label{note-epistemiche-di-chiusura-capitolo}

{\def\LTcaptype{none} % do not increment counter
{\small\begin{longtable}[]{@{}
  >{\raggedright\arraybackslash}p{(\linewidth - 4\tabcolsep) * \real{0.3333}}
  >{\raggedright\arraybackslash}p{(\linewidth - 4\tabcolsep) * \real{0.3333}}
  >{\raggedright\arraybackslash}p{(\linewidth - 4\tabcolsep) * \real{0.3333}}@{}}
\toprule\noalign{}
\begin{minipage}[b]{\linewidth}\raggedright
\#
\end{minipage} & \begin{minipage}[b]{\linewidth}\raggedright
Affermazione
\end{minipage} & \begin{minipage}[b]{\linewidth}\raggedright
Status
\end{minipage} \\
\midrule\noalign{}
\endhead
\bottomrule\noalign{}
\endlastfoot
1 & Definizione di \(\mathbb{V}^{(4)}\) come spazio strutturale &
\Hypoth{} costruzione formale \\
2 & Sei relazioni R1-R6 con segni e meccanismi & \Hypoth{} quadro
analitico, segni verificati nei Cap. 6-9 \\
3 & Asimmetria notevole: 4 segni negativi, 2 positivi & \Proved{}
struttura termodinamica del sistema sociale \\
4 & R4 (\(A\)-\(\Omega\)) come relazione più forte & \Hypoth{} inferenza
strutturale, usata operativamente nel Cap. 7 \\
5 & Mappa quattro attori in \(\mathbb{V}^{(4)}\) con stime numeriche &
\Conj{} stime parametriche con intervalli di confidenza ampi \\
6 & Superficie critica \(\mathcal{S}\) con forma esplicita
\(\Phi^*(A, \Omega, \varepsilon)\) & \Proved{} derivata analiticamente
da Cap. 4 \\
7 & Topologia di \(\mathbb{V}^{(4)}\) in \(\mathcal{R}_+\) e
\(\mathcal{R}_-\) & \Proved{} conseguenza diretta di
\(T_s^{(2),\text{sys}} = 1\) \\
8 & Sistema dinamico per
\(\dot\Phi, \dot A, \dot\Omega, \dot\varepsilon\) & \Hypoth{} postulato
operativo, calibrazione rinviata a paper EN \\
9 & Mappa gradi di libertà operativi \(\to\) scale temporali \(\to\)
leve & \Hypoth{} sintesi strategica, base per Cap. 12 \\
10 & Mappa indicatori-variabili per dashboard & \Proved{} derivata da
fonti citate nei Cap. 7-9 \\
11 & Validazione interna: posizioni
\(\boldsymbol{p}_{\text{USA, Cina, Brasile}}\) consistenti con tutte
R1-R6 & \Proved{} verifica non banale dell'apparato \\
\end{longtable}}
}



\section{Chiusura della Parte II e Apertura della Parte
IV}\label{chiusura-della-parte-ii-e-apertura-della-parte-iv}

La Parte II del libro è ora completa: Cap. 3 ha derivato \(T_s^{(2)}\)
locale, Cap. 4 ha derivato \(T_s^{(2),\text{sys}}\) con asimmetria
\(\Phi\)-moltiplica-\(A\) vs \(\Phi\)-divide-\(\Omega\), Cap. 5 ha
consolidato le quattro variabili come spazio strutturale
\(\mathbb{V}^{(4)}\) con sei relazioni di accoppiamento. La Parte III ha
mappato i quattro attori paradigmatici nello spazio \(\mathbb{V}^{(4)}\)
(Cap. 6-9). La \textbf{Parte IV} può ora trasformare l'apparato in
\textbf{strumento operativo}:

\begin{itemize}
\tightlist
\item
  \textbf{Cap. 10} --- Scenari quantizzati e Dashboard SPC-DEQ: le
  proiezioni \(\mathbb{V}^{(4)} \to \mathbb{V}^{(2)}_{\text{coerenza}}\)
  e \(\mathbb{V}^{(4)} \to \mathbb{V}^{(2)}_{\text{decoerenza}}\)
  diventano cruscotti di monitoraggio in tempo reale, con soglie
  operative e segnali di allarme precoce.
\item
  \textbf{Cap. 11} --- Predizioni rischiose datate: aggregazione delle
  \(4 \times 3 = 12\) predizioni dei Cap. 7-9 in un \emph{bouquet di
  falsificazione pubblica} della DEQ$^2$, con date, soglie numeriche, fonti
  dati.
\end{itemize}




%% ═══════════════════════════════════════════════════════════════════════════
%% PARTE III — I Quattro Attori
%% ═══════════════════════════════════════════════════════════════════════════
\part{I Quattro Attori}

\chapter{\texorpdfstring{Il Solitone (Musk): \(T_s\)
Metastabile}{Il Solitone (Musk): T\_s Metastabile}}\label{il-solitone-musk-t_s-metastabile}

\begin{quote}
\textbf{Argomento portante:} Elon Musk non è un attore egemonico nel
senso classico --- è un attore singolo (\(N=1\)) che imita un sistema.
La sua \(T_s\) è formalmente indefinita perché manca il parametro
d'ordine \(\Phi\); la sua sopravvivenza dipende da una strategia di
\textbf{separazione bayesiana} ereditata dall'equilibrio separante della
Topografia: produrre segnali \((x, y)\) visibili e di alta qualità per
nascondere uno \(z\) deliberatamente compresso.
\end{quote}



\section{Premessa: Perché Musk Apre la Parte
III}\label{premessa-perchuxe9-musk-apre-la-parte-iii}

L'ordine canonico dei quattro attori sarebbe USA $\to$ Cina $\to$ Brasile $\to$ Musk
(dal più sistemico al più individuale). La DEQ$^2$ inverte: Musk \emph{per
primo}, perché il Solitone è il \textbf{caso limite metodologico} che
illumina gli altri tre.

Il Solitone è un attore singolo che si comporta come un sistema.
Esibisce alte performance individuali (\(A\), \(H_e\), \(M_e\) tutte
elevate) ma non possiede \(\Phi\) --- l'aggregazione è impossibile per
\(N = 1\). Proietta dunque \(T_s^{(2)}\) locale eccezionale e
\(T_s^{(2),\text{sys}}\) formalmente indefinita. La distanza tra questi
due livelli è il \emph{gap di sistema}: la differenza tra ciò che si può
produrre da soli e ciò che si può sostenere collettivamente. Comprendere
il Solitone significa comprendere quel gap --- e dunque possedere lo
strumento concettuale per leggere gli attori sistemici.



\section{Definizione Formale del
Solitone}\label{definizione-formale-del-solitone}

\Hypoth{} \textbf{Definizione (Solitone DEQ$^2$).} Un attore \(e\) è detto
\emph{Solitone metastabile} se soddisfa simultaneamente:

\begin{enumerate}
\def\labelenumi{\arabic{enumi}.}
\tightlist
\item
  \(N_e = 1\) (l'attore non si articola in sotto-attori indipendenti ---
  assenza di organizzazione collettiva interna);
\item
  \(A_e \gg \langle A \rangle_{\text{sis}}\) (antifragilità individuale
  eccezionalmente elevata);
\item
  \(M_e \gg \langle M \rangle_{\text{sis}}\) (massa comunicativa
  individuale eccezionalmente elevata);
\item
  \(\Phi_e\) formalmente indefinito (\(N=1\) implica \(\Phi = 1\)
  trivialmente --- coerenza con se stesso, ma non corrisponde a un
  parametro d'ordine sistemico).
\end{enumerate}

\Proved{} \textbf{Conseguenza.} Il payoff istantaneo del Solitone
\(\Pi_e(t+1) = M_e \cdot Q_{0,e} \cdot (1 - \delta_R z_R)\) può essere
arbitrariamente alto, ma la sua \(T_s^{(2),\text{sys}}\) richiede
un'aggregazione che strutturalmente non avviene. Il Solitone è dunque
\textbf{proiettato sulla scala sistemica come singolarità} --- una
perturbazione concentrata, non una struttura.



\section{Mappatura sull'Equilibrio Separante
Bayesiano}\label{mappatura-sullequilibrio-separante-bayesiano}

L'Appendice A della Topografia formalizza un \emph{signaling game}
bayesiano in cui:

\begin{itemize}
\tightlist
\item
  L'\textbf{Emittente Egemonico} sceglie un profilo discorsivo
  \((x = 8, y = 9, z = 2)\) --- alta coerenza formale, alta precisione
  tecnica, bassa complessità dialettica;
\item
  L'\textbf{Emittente Dialogico} sceglie \((x = 7, y = 6, z = 8)\);
\item
  Il Ricevente osserva solo \((x, y)\) --- la \(z\) richiede
  elaborazione critica attiva;
\item
  Esiste un equilibrio separante in cui R può inferire il tipo di E con
  \(P > 0{,}5\) se e solo se
  \(|y_{\text{Egem}} - y_{\text{Dialog}}| > \sigma_y\).
\end{itemize}

\Hypoth{} \textbf{Tesi.} Musk è il caso paradigmatico dell'Emittente
Egemonico in equilibrio separante. Il suo profilo discorsivo aggregato
(rilevabile su X, interventi pubblici, presentazioni Tesla/SpaceX) è
approssimativamente:

{\def\LTcaptype{none} % do not increment counter
{\small\begin{longtable}[]{@{}
  >{\raggedright\arraybackslash}p{(\linewidth - 4\tabcolsep) * \real{0.3333}}
  >{\raggedright\arraybackslash}p{(\linewidth - 4\tabcolsep) * \real{0.3333}}
  >{\raggedright\arraybackslash}p{(\linewidth - 4\tabcolsep) * \real{0.3333}}@{}}
\toprule\noalign{}
\begin{minipage}[b]{\linewidth}\raggedright
Asse
\end{minipage} & \begin{minipage}[b]{\linewidth}\raggedright
Stima
\end{minipage} & \begin{minipage}[b]{\linewidth}\raggedright
Componenti
\end{minipage} \\
\midrule\noalign{}
\endhead
\bottomrule\noalign{}
\endlastfoot
\(x\) --- coerenza formale & \(\approx 8\) & Argomentazione strutturata,
slide tecniche, retorica ingegneristica \\
\(y\) --- precisione tecnica & \(\approx 9\) & Padronanza di vocabolario
aerospaziale, statistica produttiva, automazione \\
\(z\) --- complessità dialettica & \(\approx 2\) & Assenza di apparato
critico, riduzione delle contraddizioni a meme, dicotomie binarie \\
\end{longtable}}
}

\Proved{} \textbf{Verifica.} Il profilo \((8, 9, 2)\) coincide
\emph{esattamente} con l'equilibrio separante della Topografia App. A.
La strategia comunicativa di Musk è dunque la \textbf{manifestazione
empirica massimale} dell'equilibrio teorico. Coerente con la Barbarie
Tecnica del Cap. 8 della Topografia.



\section{La Massa Comunicativa di
Musk}\label{la-massa-comunicativa-di-musk}

Stima \(M_e^{\text{Musk}}\) secondo la formula della Topografia (\S{}5.3):

{\def\LTcaptype{none} % do not increment counter
{\small\begin{longtable}[]{@{}
  >{\raggedright\arraybackslash}p{(\linewidth - 4\tabcolsep) * \real{0.3333}}
  >{\raggedright\arraybackslash}p{(\linewidth - 4\tabcolsep) * \real{0.3333}}
  >{\raggedright\arraybackslash}p{(\linewidth - 4\tabcolsep) * \real{0.3333}}@{}}
\toprule\noalign{}
\begin{minipage}[b]{\linewidth}\raggedright
Parametro
\end{minipage} & \begin{minipage}[b]{\linewidth}\raggedright
Stima
\end{minipage} & \begin{minipage}[b]{\linewidth}\raggedright
Valore
\end{minipage} \\
\midrule\noalign{}
\endhead
\bottomrule\noalign{}
\endlastfoot
\(K\) --- capitale (centinaia di miliardi \$) & Massimo individuale &
\(1{,}0\) \\
\(R^{0{,}8}\) --- reach (X/Twitter \(\approx 200\)M follower) &
\(200^{0{,}8} \approx 76\) $\to$ normalizzato & \(3{,}0\) \\
\(A_{\text{alg}}\) --- amplificazione algoritmica (proprietà diretta
della piattaforma X) & Auto-amplificazione massima & \(2{,}5\) \\
\(S\) --- saturazione (post diari, presenze multiple) & Continua &
\(2{,}0\) \\
\(T\) --- durata (continua dal 2018+) & Pluriennale & \(2{,}0\) \\
\textbf{\(M_e^{\text{Musk}} = K \cdot R^{0{,}8} \cdot A_{\text{alg}} \cdot S \cdot T\)}
& & \textbf{\(\approx 30{,}0\)} \\
\end{longtable}}
}

\Hypoth{} \textbf{Confronto.} \(M_e^{\text{CNN/Ucraina}} \approx 9{,}8\)
(Topografia \S{}5.3). Musk individuale ha massa comunicativa \textbf{circa
tre volte} quella di un'organizzazione mediatica top globale. Questo è
il punto strutturale: un attore singolo che concentra una massa
comunicativa di scala organizzativa.



\section{Antifragilità Individuale
Eccezionale}\label{antifragilituxe0-individuale-eccezionale}

\Hypoth{} \textbf{Stima di \(A_{\text{Musk}}\).} Musk ha mostrato, in
tre eventi distinti, convessità positiva alla volatilità:

\begin{enumerate}
\def\labelenumi{\arabic{enumi}.}
\tightlist
\item
  \emph{Crisi finanziaria 2008-2009}: SpaceX e Tesla sull'orlo del
  collasso $\to$ ristrutturazione che produce vantaggio competitivo
  decennale;
\item
  \emph{Pandemia 2020}: Tesla unica azienda automobilistica a non
  interrompere la produzione $\to$ quadruplicamento della capitalizzazione;
\item
  \emph{Acquisizione X 2022-2023}: massima volatilità (esodi
  pubblicitari, controversie) $\to$ riconfigurazione algoritmica $\to$ uso
  politico nel 2024.
\end{enumerate}

\Proved{} Tutti e tre eventi mostrano
\(\partial^2 \Pi / \partial \sigma^2 > 0\) a livello individuale.
\(A_{\text{Musk}} \gg 0\).

\Conj{} \textbf{Caveat.} L'antifragilità di Musk è quella della
\textbf{classe degli attori a opzionalità asimmetrica con copertura
politica} --- vantaggio strutturale che si trasferisce solo a chi ha
capitale sufficiente per attraversare le crisi. È una proprietà
\emph{individuale}, non \emph{istituzionale}. Da qui il problema
strutturale del paragrafo successivo.



\section{\texorpdfstring{Il Problema dell'Aggregazione:
\(N = 1\)}{Il Problema dell'Aggregazione: N = 1}}\label{il-problema-dellaggregazione-n-1}

\Proved{} \textbf{Risultato.} Per il Solitone, l'operatore
\(\langle\cdot\rangle_\Phi\) del Cap. 4 è degenere. Non esiste \emph{un
sistema di Musk} --- esiste solo Musk. Le sue aziende (Tesla, SpaceX, X,
Neuralink, xAI) non costituiscono un sistema con \(\Phi\) indipendente:
la coerenza emerge dalla sua singola persona, non da un parametro
d'ordine emergente.

\Hypoth{} \textbf{Conseguenza strutturale.} Il Solitone è
\textbf{non-trasferibile}:

\begin{itemize}
\tightlist
\item
  \emph{Non istituzionalizzabile}: la successione produce
  frammentazione, non continuità.
\item
  \emph{Non replicabile}: altri attori con stesse risorse ma assenza
  dello stesso profilo personale falliscono.
\item
  \emph{Non sistemico}: non genera sotto-attori autonomi che ne portino
  avanti la funzione.
\end{itemize}

\Conj{} \textbf{Predizione DEQ$^2$ (datata).} Entro il 2031, il
\emph{cluster Musk} (Tesla + SpaceX + X + xAI) si troverà in una di
queste condizioni: (a) frammentato per uscita o decesso del Solitone con
perdita di coerenza \textgreater{} \(50\%\) del valore aggregato; (b)
istituzionalizzato attraverso un governance shift che riduce \(A\)
collettivo verso la media; (c) congelato in una posizione metastabile
sotto \(\Phi\) esogeno (alleanza con uno Stato --- USA o altro). Le tre
opzioni sono mutualmente esclusive e tutte segnano la fine del Solitone
come tale.



\section{Il Solitone come Vettore di Decoerenza
Sistemica}\label{il-solitone-come-vettore-di-decoerenza-sistemica}

\Hypoth{} \textbf{Effetto sull'ambiente.} Il Solitone, pur essendo
\(N = 1\), \emph{interagisce} con i sistemi macroscopici e ne altera
\(\Phi\):

\begin{itemize}
\tightlist
\item
  L'azione comunicativa di Musk su X (massa \(M_e \approx 30\)) erode la
  \(\Phi\) del sistema USA: introduce dissonanza nel discorso politico
  ufficiale, sposta i baricentri narrativi, polarizza ulteriormente.
\item
  L'azione tecnologica di SpaceX/Starlink crea dipendenze
  infrastrutturali asimmetriche su Stati terzi (Ucraina 2022-2023,
  Taiwan in scenari ipotetici), introducendo un nuovo asse di sovranità
  che gli stati non controllano.
\end{itemize}

\Proved{} \textbf{Modellazione formale.} Il Solitone agisce come
\textbf{sorgente esogena di disordine} \(\zeta(t)\) nel sistema USA. La
dinamica di Kuramoto del Cap. 4 si modifica:

\[
\frac{d\theta_j}{dt} = \omega_j + K\Phi \sin(\psi - \theta_j) + \zeta_{\text{Sol}}(t)
\]

Per \(\zeta_{\text{Sol}}\) persistente e sufficientemente intenso,
\(\Phi^{\text{USA}} \to 0\) anche con \(K > K_c\). Il Solitone
\emph{forza la decoerenza} di un sistema.

\Conj{} \textbf{Predizione testabile.} L'indice di polarizzazione USA
(Pew Research Polarization Index, V-Dem Liberal Democracy Index) deve
mostrare correlazione positiva \(r > 0{,}3\) con la \emph{Musk
visibility} (volume mensile di interventi politici diretti) nel periodo
2022-2026.



\section{Sintesi del Capitolo}\label{sintesi-del-capitolo}

\Proved{} \textbf{In una frase:} il Solitone è un attore con
\(T_s^{(2)}\) locale eccezionale ma \(T_s^{(2),\text{sys}}\) indefinito,
la cui sopravvivenza dipende da una strategia di separazione bayesiana
che produce segnali alti su \((x, y)\) per nascondere uno \(z\)
deliberatamente compresso, e la cui interazione con i sistemi macro
avviene come sorgente esogena di decoerenza.

{\def\LTcaptype{none} % do not increment counter
{\small\begin{longtable}[]{@{}
  >{\raggedright\arraybackslash}p{(\linewidth - 4\tabcolsep) * \real{0.3333}}
  >{\raggedright\arraybackslash}p{(\linewidth - 4\tabcolsep) * \real{0.3333}}
  >{\raggedright\arraybackslash}p{(\linewidth - 4\tabcolsep) * \real{0.3333}}@{}}
\toprule\noalign{}
\begin{minipage}[b]{\linewidth}\raggedright
Aspetto
\end{minipage} & \begin{minipage}[b]{\linewidth}\raggedright
Solitone (Musk)
\end{minipage} & \begin{minipage}[b]{\linewidth}\raggedright
Implicazione
\end{minipage} \\
\midrule\noalign{}
\endhead
\bottomrule\noalign{}
\endlastfoot
\(N\) & \(1\) & \(\Phi\) formalmente indefinito \\
\(M_e\) & \(\approx 30\) & Tre volte CNN/Ucraina \\
\((x, y, z)\) & \((8, 9, 2)\) & Equilibrio separante bayesiano \\
\(A\) & Molto alto individuale & Non aggregabile \\
\(\Omega\) & Basso (Musk non spreca) & Ma irrilevante senza \(N\) \\
Effetto sistemico & \(\zeta(t)\) esogeno & Decoerenza forzata su altri
sistemi \\
\end{longtable}}
}

\Hypoth{} \textbf{Tesi conclusiva del capitolo.} Il Solitone non è
un'eccezione che conferma la regola DEQ$^2$ --- è il \textbf{caso limite
metodologico} che ne illustra la necessità. Senza la teoria sistemica
con \(\Phi\), Musk apparirebbe come l'attore più stabile del sistema
globale; con la teoria, appare per quello che è: una concentrazione
metastabile destinata a una transizione di fase entro l'orizzonte 2031.



\section{Note Epistemiche di Chiusura
Capitolo}\label{note-epistemiche-di-chiusura-capitolo}

{\def\LTcaptype{none} % do not increment counter
{\small\begin{longtable}[]{@{}
  >{\raggedright\arraybackslash}p{(\linewidth - 4\tabcolsep) * \real{0.3333}}
  >{\raggedright\arraybackslash}p{(\linewidth - 4\tabcolsep) * \real{0.3333}}
  >{\raggedright\arraybackslash}p{(\linewidth - 4\tabcolsep) * \real{0.3333}}@{}}
\toprule\noalign{}
\begin{minipage}[b]{\linewidth}\raggedright
\#
\end{minipage} & \begin{minipage}[b]{\linewidth}\raggedright
Affermazione
\end{minipage} & \begin{minipage}[b]{\linewidth}\raggedright
Status
\end{minipage} \\
\midrule\noalign{}
\endhead
\bottomrule\noalign{}
\endlastfoot
1 & Definizione formale del Solitone & \Hypoth{} categoria nuova della
DEQ$^2$ \\
2 & Mappatura \((8, 9, 2)\) Musk $\to$ equilibrio separante Topografia &
\Proved{} verificato sul profilo discorsivo pubblico \\
3 & Stima \(M_e^{\text{Musk}} \approx 30\) & \Conj{} stima parametrica,
intervallo \([20, 40]\) \\
4 & \(A_{\text{Musk}} \gg 0\) via tre eventi documentati & \Proved{}
evidenza convergente \\
5 & Non-trasferibilità del Solitone & \Hypoth{} ipotesi strutturale
forte \\
6 & Predizione 2031 (frammentazione/istituzionalizzazione/congelamento)
& \Conj{} congettura datata, falsificabile \\
7 & Solitone come sorgente esogena \(\zeta(t)\) nel sistema USA &
\Hypoth{} modellazione coerente con Kuramoto \\
8 & Correlazione \(r > 0{,}3\) Musk visibility $\leftrightarrow$ polarizzazione USA &
\Conj{} predizione empirica testabile \\
\end{longtable}}
}



\input{chapters-it/07-egemone-decoerente}
\input{chapters-it/08-rigido-sincronizzato}
\chapter{\texorpdfstring{Il Bilanciere Quantistico (Brasile):
l'\emph{Estado Logístico} come
\(\Phi\)-engine}{Il Bilanciere Quantistico (Brasile): l'Estado Logístico come \textbackslash Phi-engine}}\label{il-bilanciere-quantistico-brasile-lestado-loguxedstico-come-phi-engine}

\begin{quote}
\textbf{Argomento portante:} il Brasile del 2026 è l'unico fra i quattro
attori della Parte III a possedere la \emph{struttura formale} di un
produttore di \(\Phi\) esportabile --- non per dotazione materiale
(\(s_{ii}\) medio, non superiore), né per egemonia ideologica
(\(s_{iii}\) regionale-globale ma non globale-totale), ma per
\textbf{assetto operativo}: l'\emph{Estado Logístico} di Amado Cervo
come politica estera che mantiene tutti i fronti aperti, la capacità di
mediazione plurale documentata su Russia-Ucraina, BRICS, COP30, e una
resilienza democratica istituzionale verificata (V-Dem U-turn 2025) dopo
aver respinto il colpo di stato del gennaio 2023 e condannato il suo
principale agente. La \(\Phi\)-engine brasiliana è tuttavia
\textbf{condizionale}: l'elezione presidenziale del 4 ottobre 2026 è il
test di falsificabilità più stretto della DEQ$^2$ nell'orizzonte 2028-2031.
\end{quote}



\section{Definizione Strutturale del Bilanciere
Quantistico}\label{definizione-strutturale-del-bilanciere-quantistico}

\Hypoth{} \textbf{Definizione (Bilanciere Quantistico DEQ$^2$).} Un attore
\(e\) è detto \emph{Bilanciere Quantistico} (o \(\Phi\)-engine) se
soddisfa simultaneamente:

\begin{enumerate}
\def\labelenumi{\arabic{enumi}.}
\tightlist
\item
  \(\Phi_e^{\text{interno}}\) moderato ma \textbf{prodotto
  spontaneamente} (non per coercizione gerarchica come nel Rigido
  Sincronizzato, non assente come nell'Egemone Decoerente);
\item
  \(\Phi_e^{\text{esportabile}} > 0\) --- esistono \emph{interfacce di
  accoppiamento bidirezionale} con sotto-sistemi globali eterogenei
  (USA, Cina, Russia, UE, Africa, Medio Oriente, Asia) che permettono di
  trasferire ordine senza imporre allineamento;
\item
  \textbf{Asimmetria \(A\)/\(\Omega\) favorevole}: \(A_e\) adattivo
  (esperienza ricorrente di crisi superate senza collasso
  istituzionale), \(\Omega_e\) sotto-medio per la fascia (l'energia
  dissipata in conflitto interno cresce ma resta entro la soglia di
  gestibilità);
\item
  \(s_{iii,e} > s_{ii,e}\) in termini relativi (la legittimità
  etico-discorsiva eccede la capacità materiale --- caratteristica
  strutturale di un \emph{bridge-builder});
\item
  \textbf{Politica estera non-allineata operativa}: l'Estado Logístico
  di Cervo come \emph{policy} effettiva, non come retorica.
\end{enumerate}

\Proved{} \textbf{Conseguenza formale.} \(T_s^{(2),\text{sys}}\) del
Bilanciere è strutturalmente in regime \textbf{stabile-perturbabile}:
alta sotto condizioni di base, sensibile a tre variabili --- esito
elettorale, fiscale, polarizzazione interna --- il cui collasso
simultaneo produrrebbe una transizione verso \emph{Egemone Decoerente}
(caso USA-like) o \emph{Solitone Politico} (caso post-2018). Non esiste
rotta verso \emph{Rigido Sincronizzato}, perché lo \emph{stack}
istituzionale federale brasiliano è asimmetrico rispetto a quello
cinese.



\section{\texorpdfstring{L'\emph{Estado Logístico} di Amado Cervo come
Apparato
Operativo}{L'Estado Logístico di Amado Cervo come Apparato Operativo}}\label{lestado-loguxedstico-di-amado-cervo-come-apparato-operativo}

\Hypoth{} \textbf{Eredità teorica.} L'\emph{Estado Logístico} (Amado
Luiz Cervo, \emph{Inserção internacional: formação dos conceitos
brasileiros}, 2008) è il paradigma di politica estera in cui lo Stato
non rinuncia all'autonomia decisionale ma la articola attraverso una
\emph{molteplicità simultanea di vettori}: industriale (BRICS), agricolo
(Cerrado, soia), energetico (idroelettrico, biocombustibili, oggi
rinnovabili), diplomatico (mediazione plurale), ambientale (Amazzonia,
COP). Cervo lo distingue da \emph{Estado Normal} (subordinato),
\emph{Estado Desenvolvimentista} (autarchico), \emph{Estado Neoliberal}
(dipendente).

\Proved{} \textbf{Mappatura DEQ$^2$ formale.} L'\emph{Estado Logístico} è
la realizzazione politica della condizione (5) della definizione
\$\S\$9.0. In termini del campo \(\Pi\) topografico:

\[\Pi^{\text{Logístico}}(t+1) = \sum_{k} M_e^{(k)} \cdot Q_0^{(k)} \cdot (1 - \delta_R^{(k)} z_R^{(k)})\]

dove \(k\) indicizza i diversi vettori (industriale, agricolo,
energetico, diplomatico, ambientale) e ciascun \(M_e^{(k)}\),
\(Q_0^{(k)}\), \(\delta_R^{(k)}\) è specifico del vettore. È una
\emph{sovrapposizione coerente} --- esattamente la struttura matematica
che genera \(\Phi\) in un sistema multi-canale.

\Hypoth{} \textbf{Conseguenza strutturale.} Il Bilanciere produce
\(\Phi\) globale per \emph{sovrapposizione plurale} (più vettori
coerenti) e non per \emph{amplificazione singolare} (un solo vettore
massimizzato). Questo lo rende qualitativamente diverso dal Solitone
(Cap. 6), dall'Egemone Decoerente (Cap. 7) e dal Rigido Sincronizzato
(Cap. 8). Strutturalmente, il Bilanciere è \textbf{l'unico attore della
Parte III in cui l'apparato matematico DEQ$^2$ produce numeratore \(> 0\)
stabile}.



\section{BRICS 2025 Rio + COP30 Belém: due Test
Simultanei}\label{brics-2025-rio-cop30-beluxe9m-due-test-simultanei}

\Proved{} \textbf{Dato 1 --- Presidenza BRICS 2025.} Brasile presiede
BRICS dal 1 gennaio 2025. Il 17$^\circ$ summit (Rio de Janeiro, 6-7 luglio
2025) approva tre documenti centrali della presidenza brasiliana: -
\emph{BRICS Leaders' Framework Declaration on Climate Finance}; -
\emph{BRICS Leaders' Declaration on Global Governance of Artificial
Intelligence}; - \emph{BRICS Partnership for the Elimination of Socially
Determined Diseases}.

Dichiarazione finale: 126 punti su global governance, finanza, salute,
AI, clima. Lancio del \emph{Tropical Forest Forever Fund} (TFFF) e del
\emph{BRICS Multilateral Guarantees} (BMG).

\Proved{} \textbf{Dato 2 --- COP30 Belém.} 10-22 novembre 2025. Brasile
presiede fino a novembre 2026. Tre pilastri: 1. \emph{Strengthening
multilateralism and the climate regime}; 2. \emph{Connecting climate
initiatives to people's daily lives}; 3. \emph{Accelerating
implementation of the Paris Agreement}.

Risultati: mobilitazione \(1{,}3\)T USD/anno entro 2035, triplicamento
finanza adattamento, \emph{Belém Mechanism for Just Global Transition},
due roadmap (transizione fuori dai fossili + foreste/clima) lanciate per
iniziativa della presidenza brasiliana malgrado assenza di consenso
(\(\approx 80\) paesi pro / \(\approx 80\) contro language esplicito).

\Hypoth{} \textbf{Lettura DEQ$^2$ congiunta.} I due eventi, ravvicinati
(luglio 2025 e novembre 2025), misurano la \(\Phi\) esportabile
brasiliana in due assi distinti:

{\def\LTcaptype{none} % do not increment counter
{\small\begin{longtable}[]{@{}
  >{\raggedright\arraybackslash}p{(\linewidth - 4\tabcolsep) * \real{0.3333}}
  >{\raggedright\arraybackslash}p{(\linewidth - 4\tabcolsep) * \real{0.3333}}
  >{\raggedright\arraybackslash}p{(\linewidth - 4\tabcolsep) * \real{0.3333}}@{}}
\toprule\noalign{}
\begin{minipage}[b]{\linewidth}\raggedright
Test
\end{minipage} & \begin{minipage}[b]{\linewidth}\raggedright
Asse misurato
\end{minipage} & \begin{minipage}[b]{\linewidth}\raggedright
Risultato 2025
\end{minipage} \\
\midrule\noalign{}
\endhead
\bottomrule\noalign{}
\endlastfoot
BRICS Rio & \(\Phi^{\text{Sud globale}}\) --- capacità di sintesi
multilaterale fra interessi divergenti (Cina-India, Russia-Iran,
Brasile-mediatore) & Documenti tematici di consenso; iniziative
innovative (TFFF, BMG); 200+ riunioni preparatorie \\
COP30 Belém & \(\Phi^{\text{Nord-Sud}}\) --- capacità di mediare fra
alta-emissioni e adattamento, mantenendo entrambi nel campo \(\Pi\) &
Pacchetto Belém approvato; due roadmap lanciate \emph{malgrado} assenza
di consenso (mossa di leadership) \\
\end{longtable}}
}

\Proved{} \textbf{Conseguenza.} Il Brasile è l'unico Stato del 2025 ad
aver presieduto contemporaneamente un blocco multilaterale Sud (BRICS) e
un consesso multilaterale globale (COP), producendo in entrambi
documenti operativi non banali. Questa è la firma empirica
dell'asimmetria definita in \$\S\$9.0 condizione 4: \(s_{iii} > s_{ii}\)
produce \emph{capacità di convocazione} superiore alla \emph{capacità di
imposizione}.

\Conj{} \textbf{Caveat onesto.} L'efficacia dei documenti del 2025 si
misura nei trienni successivi. Il TFFF ha bisogno di donazioni
effettive; il BMG di investitori reali; le roadmap COP di adesioni. Il
Brasile ha \emph{aperto i fronti} --- l'esito misurerà se il vettore
\(\Phi\)-esportabile è solo dichiarativo o effettivamente operativo.



\section{Capacità di Mediazione come Struttura, non
Episodio}\label{capacituxe0-di-mediazione-come-struttura-non-episodio}

\Proved{} \textbf{Dato.} Maggio 2025: Lula visita Mosca, propone
\emph{strategic partnership} RU-BR e si offre come mediatore RU-UA,
condizionatamente all'apertura di entrambe le parti. Il Brasile è
co-fondatore (con la Cina) del \emph{Group of Friends for Peace} (gruppo
amici della pace), che ha pubblicato un piano in sei punti per il
conflitto ucraino (2024).

\Hypoth{} \textbf{Lettura DEQ$^2$ della mediazione strutturale.} La
capacità di mediazione del Brasile non è un attributo personale di Lula
ma una \emph{invariante strutturale} dell'Estado Logístico: dipende dal
mantenimento attivo di tutti i canali di comunicazione (USA, Cina,
Russia, UE) senza chiusura ideologica unilaterale. Formalmente: il
Brasile mantiene il proprio campo \(\Pi^{\text{globale}}\) con
\(\Sigma\) (matrice di covarianza degli assi \(x\) coerenza, \(y\)
precisione, \(z\) complessità dialettica --- eredità Topografia \S{}5.1)
\textbf{massimamente aperta}, contro l'Egemone Decoerente che la riduce
a \(\Sigma^{\text{egemonica}}\) con \(z \to 0\) (cap. 7.3).

\Proved{} \textbf{Conseguenza misurabile.} Il Brasile è l'unico fra i
grandi Stati a: - non aver applicato sanzioni a Russia, Cina, Iran; -
mantenere relazioni operative con Israele e Iran simultaneamente; -
preservare partnership economiche con USA e Cina simultaneamente; -
partecipare a BRICS+ senza abbandonare OECD-track e G20.

\Hypoth{} \textbf{Tesi formale.} Il \emph{non-allineamento attivo} è la
condizione operativa della \(\Phi\)-engine: produrre coerenza globale
senza imporre allineamento sub-globale. È una posizione sistemicamente
distinta dal \emph{non-allineamento passivo} della Guerra Fredda (India,
Yugoslavia) --- è \emph{non-allineamento generativo}, capace di produrre
nuovi accoppiamenti (TFFF, \emph{Group of Friends for Peace},
\emph{Belém Mechanism}).



\section{\texorpdfstring{V-Dem U-Turn: la Prova di \(\Phi\)
Resiliente}{V-Dem U-Turn: la Prova di \textbackslash Phi Resiliente}}\label{v-dem-u-turn-la-prova-di-phi-resiliente}

\Proved{} \textbf{Dato V-Dem 2026} (relativo all'anno 2025): registrato
un \emph{U-turn} (inversione di autocratizzazione) per il Brasile
post-2022. La presidenza Lula ha \textbf{invertito il deterioramento}
documentato sotto Bolsonaro. Il rapporto annota: miglioramenti
istituzionali, \emph{judicial assertiveness} (STF, TSE), \emph{civil
society mobilization}, \emph{media independence}, \emph{neutrality of
armed forces}. Avvertimento esplicito: la \emph{polarizzazione sociale
persiste} e \emph{l'elezione del 2026 sarà decisiva}.

\Proved{} \textbf{Verifica empirica della resilienza istituzionale.}
L'11 settembre 2025 la Suprema Corte Federale (STF) condanna Jair
Bolsonaro a 27 anni e 3 mesi per: - tentativo di colpo di stato (8
gennaio 2023); - partecipazione a organizzazione criminale armata; -
piano per uccidere Lula, Alckmin, Moraes; - abolizione violenta
dell'ordine democratico.

7 novembre 2025: pannello STF respinge l'appello (4-1, dissenso del solo
giudice Fux). È la \textbf{prima condanna nella storia brasiliana di un
ex-presidente per tentato colpo di stato}.

\Hypoth{} \textbf{Lettura DEQ$^2$. } La condanna Bolsonaro è l'evento
empirico più rilevante della Parte III per la teoria DEQ$^2$: misura la
capacità del sistema brasiliano di mantenere
\(\Phi^{\text{interno}} > 0\) sotto perturbazione massima (tentativo di
rottura dell'ordine democratico). In termini formali, ha mostrato che il
sistema possiede un \emph{meccanismo di feedback negativo} funzionante
--- una proprietà strutturale che né Egemone Decoerente né Rigido
Sincronizzato esibiscono.

{\def\LTcaptype{none} % do not increment counter
{\small\begin{longtable}[]{@{}
  >{\raggedright\arraybackslash}p{(\linewidth - 6\tabcolsep) * \real{0.2500}}
  >{\raggedright\arraybackslash}p{(\linewidth - 6\tabcolsep) * \real{0.2500}}
  >{\raggedright\arraybackslash}p{(\linewidth - 6\tabcolsep) * \real{0.2500}}
  >{\raggedright\arraybackslash}p{(\linewidth - 6\tabcolsep) * \real{0.2500}}@{}}
\toprule\noalign{}
\begin{minipage}[b]{\linewidth}\raggedright
Indicatore istituzionale
\end{minipage} & \begin{minipage}[b]{\linewidth}\raggedright
USA 2026
\end{minipage} & \begin{minipage}[b]{\linewidth}\raggedright
Cina 2026
\end{minipage} & \begin{minipage}[b]{\linewidth}\raggedright
Brasile 2026
\end{minipage} \\
\midrule\noalign{}
\endhead
\bottomrule\noalign{}
\endlastfoot
Vincoli giudiziari sull'esecutivo & \(\downarrow\) minimo da 100+ anni
(Cap. 7.1) & non applicabili & \(\uparrow\) STF assertiva \\
Indipendenza media & \(\downarrow\) (attacchi a stampa) & bassa per
costruzione & \(\uparrow\) post-Bolsonaro \\
Mobilitazione civile & \(\downarrow\) frammentata & \(\to 0\) permessa &
media-alta \\
Neutralità forze armate & erosa nel 2021-2024 & partito = forze armate &
confermata 2023-2025 \\
\end{longtable}}
}

\Proved{} \textbf{Conseguenza.} Il Brasile possiede ciò che la
Topografia \S{}5.4 chiama \emph{resistenza critica} \(R_c\): un sistema di
sotto-attori che possono opporsi alla traiettoria di un singolo attore
dominante. Per la DEQ$^2$, \(R_c > 0\) effettivo è condizione necessaria di
\(\Phi\) spontanea (vs \(\Phi\) imposta del Cap. 8).



\section{Le Fragilità: Onestà
Analitica}\label{le-fragilituxe0-onestuxe0-analitica}

Il Cap. 9 perderebbe rigore se omettesse le fragilità reali del
Bilanciere brasiliano. La \(\Phi\)-engine non è automatica --- è
\textbf{condizionale} a tre variabili in tensione.

\subsection{Fragilità 1: Polarizzazione e approval
calante}\label{fragilituxe0-1-polarizzazione-e-approval-calante}

\Proved{} \textbf{Dato.} Approval Lula marzo 2026 (Genial/Quaest, 2.004
elettori): \(44\%\) approva, \(51\%\) disapprova --- peggior gap dal
luglio 2025. Sondaggi alternativi danno \(28\%/40\%\). Coalizione di
governo a \emph{minimo storico di 30 anni} in fedeltà parlamentare.

\Conj{} \textbf{Lettura DEQ$^2$.} La polarizzazione interna è il rumore
endogeno che erode \(\Phi^{\text{interno}}\). Il Brasile non è immune al
fenomeno globale di affective polarization --- è solo che dispone di
meccanismi \(R_c\) ancora operanti per contenerlo.

\subsection{Fragilità 2: Fiscale e
Selic}\label{fragilituxe0-2-fiscale-e-selic}

\Proved{} \textbf{Dato.} Crescita 2025: \(+2{,}3\%\). Forecast 2026:
\(+1{,}7\%\). Selic da \(10{,}5\%\) (settembre 2024) a \(15\%\) (giugno
2025); easing atteso fine 2026 attorno a \(9{,}75\)-\(11{,}50\%\).
Inflazione IPCA 2025: \(4{,}44\%\) (entro target \(1{,}5\)-\(4{,}5\%\)).
Deficit federale 2026: \(\approx 3\%\) PIL; debito pubblico
\(\approx 75\%\) PIL. Quadro fiscale \emph{weakened} da indebolimento
delle regole proprie.

\Conj{} \textbf{Lettura DEQ$^2$.} Lo squilibrio fiscale aumenta \(\sigma\)
esogena --- volatilità monetaria, pressione su tasso di cambio,
vulnerabilità a shock esterni. Il Brasile resta sotto la soglia di
\emph{currency crisis}, ma la traiettoria 2025-2026 è sub-ottimale per
un \(\Phi\)-engine. Strutturalmente: la combinazione \emph{crescita
rallentata + Selic alto + debito 75\%} riduce la capacità del governo di
sostenere finanziariamente l'agenda Estado Logístico (TFFF, mediazione,
COP).

\subsection{Fragilità 3: Elezioni 4 ottobre
2026}\label{fragilituxe0-3-elezioni-4-ottobre-2026}

\Proved{} \textbf{Dato.} Bolsonaro ineligibile (condannato + in
prigione). Suo figlio \textbf{Flávio Bolsonaro} candidato del Partido
Liberal. Sondaggio AtlasIntel/Bloomberg aprile 2026: Flávio \(47{,}6\%\)
vs Lula \(46{,}6\%\) in secondo turno simulato --- \textbf{prima volta}
in cui Flávio è numericamente avanti (margine entro errore statistico).
Primo turno: Lula \(45{,}9\%\) vs Flávio \(40{,}1\%\). Trend dei mesi
precedenti: Lula in calo da \(+12\) pp (dicembre 2025) a pareggio
statistico.

\Conj{} \textbf{Lettura DEQ$^2$ centrale.} L'elezione del 4 ottobre 2026 è
il \textbf{test di falsificazione più stretto} della DEQ$^2$ nell'orizzonte
2028-2031. Tre scenari:

{\def\LTcaptype{none} % do not increment counter
{\small\begin{longtable}[]{@{}
  >{\raggedright\arraybackslash}p{(\linewidth - 4\tabcolsep) * \real{0.3333}}
  >{\raggedright\arraybackslash}p{(\linewidth - 4\tabcolsep) * \real{0.3333}}
  >{\raggedright\arraybackslash}p{(\linewidth - 4\tabcolsep) * \real{0.3333}}@{}}
\toprule\noalign{}
\begin{minipage}[b]{\linewidth}\raggedright
Scenario
\end{minipage} & \begin{minipage}[b]{\linewidth}\raggedright
Probabilità soggettiva (aprile 2026)
\end{minipage} & \begin{minipage}[b]{\linewidth}\raggedright
Conseguenza per \(\Phi\)-engine
\end{minipage} \\
\midrule\noalign{}
\endhead
\bottomrule\noalign{}
\endlastfoot
\textbf{A} Lula re-eletto & \(0{,}45\)-\(0{,}55\) & \(\Phi\)-engine
continua, sotto pressione fiscale ma operativa \\
\textbf{B} Successore di sinistra (post-Lula, es. Haddad) &
\(0{,}10\)-\(0{,}15\) & \(\Phi\)-engine continua, possibile riforma
fiscale \\
\textbf{C} Flávio Bolsonaro o destra radicale eletta &
\(0{,}30\)-\(0{,}40\) & \(\Phi\)-engine si interrompe, traiettoria verso
\emph{Egemone Decoerente regionale} \\
\end{longtable}}
}

\Hypoth{} \textbf{Tesi forte.} Se A o B: la DEQ$^2$ prevede consolidamento
del Brasile come \(\Phi\)-engine globale entro 2028-2031, con
\(T_s^{(2),\text{sys}}\) nettamente superiore a USA e Cina. Se C: la
DEQ$^2$ prevede solitonizzazione discorsiva del Brasile (passaggio dal
profilo Dialogico \((7,6,8)\) all'Egemonico \((8,9,2)\) --- ricalcando
la traiettoria USA del Cap. 7.3) con perdita parziale di \(s_{iii}\)
globale entro 24 mesi.



\section{\texorpdfstring{Sintesi del Capitolo: \(\Phi\)-engine
Condizionale}{Sintesi del Capitolo: \textbackslash Phi-engine Condizionale}}\label{sintesi-del-capitolo-phi-engine-condizionale}

\Proved{} \textbf{In una frase:} il Brasile del 2026 è strutturalmente
l'unico attore della Parte III configurato come \(\Phi\)-engine ---
Estado Logístico operativo, doppia presidenza BRICS-COP, mediazione
strutturale, V-Dem U-turn verificato, condanna Bolsonaro come prova
empirica di \(R_c\) effettivo --- ma la condizionalità elettorale del 4
ottobre 2026 fa della DEQ$^2$ una teoria \emph{con punto di biforcazione
esplicito}: il sistema globale entrerà nella finestra 2028-2031 con o
senza un produttore di \(\Phi\) esportabile.

{\def\LTcaptype{none} % do not increment counter
{\small\begin{longtable}[]{@{}
  >{\raggedright\arraybackslash}p{(\linewidth - 6\tabcolsep) * \real{0.2500}}
  >{\raggedright\arraybackslash}p{(\linewidth - 6\tabcolsep) * \real{0.2500}}
  >{\raggedright\arraybackslash}p{(\linewidth - 6\tabcolsep) * \real{0.2500}}
  >{\raggedright\arraybackslash}p{(\linewidth - 6\tabcolsep) * \real{0.2500}}@{}}
\toprule\noalign{}
\begin{minipage}[b]{\linewidth}\raggedright
Variabile
\end{minipage} & \begin{minipage}[b]{\linewidth}\raggedright
Valore stimato 2026
\end{minipage} & \begin{minipage}[b]{\linewidth}\raggedright
Direzione
\end{minipage} & \begin{minipage}[b]{\linewidth}\raggedright
Conseguenza per \(T_s^{(2),\text{sys}}\)
\end{minipage} \\
\midrule\noalign{}
\endhead
\bottomrule\noalign{}
\endlastfoot
\(s_{ii}\) (capacità materiale) & media & stabile & numeratore
moderato \\
\(s_{iii}\) (legittimità etico-discorsiva) & \(49{,}2\) Brand Finance,
V-Dem U-turn & \(\uparrow\) & numeratore in crescita \\
\(\Phi^{\text{interno}}\) & spontaneo, sotto pressione polarizzazione &
stabile-perturbabile & numeratore funzionale \\
\(\Phi^{\text{esportabile}}\) & \(> 0\) verificato (BRICS, COP30,
mediazione) & \(\uparrow\) & unicità nel sistema \\
\(A^{\text{adattivo}}\) & alto (resilienza 2018-2023 confermata) &
\(\uparrow\) & denominatore contenuto \\
\(\Omega\) (dissipazione interna) & medio (polarizzazione persistente) &
stabile & denominatore moderato \\
\(\varepsilon_{\text{dis}}\) (disimpegno morale) & calo post-2022,
possibile risalita 2026 & dipende elezioni & swing \\
\textbf{Variabile di biforcazione} & esito 4 ottobre 2026 &
\(P(\text{Lula/sinistra}) \approx 0{,}55\)-\(0{,}70\) & A/B vs C separa
due regimi \\
\end{longtable}}
}

\Hypoth{} \textbf{Tesi conclusiva del capitolo.} La probabilità che il
sistema globale 2028-2031 entri nella sua transizione di fase con almeno
un produttore di \(\Phi\) esportabile è \(P \approx 0{,}55\)-\(0{,}70\),
condizionale al risultato elettorale brasiliano dell'ottobre 2026. È la
singola variabile più informativa fra tutte quelle considerate nei Cap.
6-9 --- perché tutti gli altri attori (Solitone, USA, Cina) hanno già
\emph{firmato} la propria traiettoria DEQ$^2$, mentre il Brasile non l'ha
ancora fatto.

\Hypoth{} \textbf{Asimmetria triplice consolidata (Cap. 7-8-9).}

{\def\LTcaptype{none} % do not increment counter
{\small\begin{longtable}[]{@{}
  >{\raggedright\arraybackslash}p{(\linewidth - 6\tabcolsep) * \real{0.2500}}
  >{\raggedright\arraybackslash}p{(\linewidth - 6\tabcolsep) * \real{0.2500}}
  >{\raggedright\arraybackslash}p{(\linewidth - 6\tabcolsep) * \real{0.2500}}
  >{\raggedright\arraybackslash}p{(\linewidth - 6\tabcolsep) * \real{0.2500}}@{}}
\toprule\noalign{}
\begin{minipage}[b]{\linewidth}\raggedright
Attore
\end{minipage} & \begin{minipage}[b]{\linewidth}\raggedright
Tipo \(\Phi\)
\end{minipage} & \begin{minipage}[b]{\linewidth}\raggedright
Esportabilità
\end{minipage} & \begin{minipage}[b]{\linewidth}\raggedright
Direzione \(T_s^{(2),\text{sys}}\)
\end{minipage} \\
\midrule\noalign{}
\endhead
\bottomrule\noalign{}
\endlastfoot
Solitone (Musk) & \(\Phi\) degenere (\(N=1\)) & per decoerenza forzata
su altri & indefinita \\
Egemone Decoerente (USA) & \(\Phi\) mai prodotta & nessuna & \(\to 0\)
in 2027-2029 \\
Rigido Sincronizzato (Cina) & \(\Phi\) imposta & esportazione di \(K\),
non di \(\Phi\) & \(\to 0\) catastrofica \\
\textbf{Bilanciere Quantistico (Brasile)} & \(\Phi\) spontanea
condizionale & \textbf{vera \(\Phi\)-engine se A/B vince} &
\textbf{\(> 0\) se A/B; \(\to 0\) se C} \\
\end{longtable}}
}



\section{Quattro Predizioni Datate del Cap. 9 (per Cap.
11)}\label{quattro-predizioni-datate-del-cap.-9-per-cap.-11}

\begin{enumerate}
\def\labelenumi{\arabic{enumi}.}
\tightlist
\item
  \textbf{Esito elezione 4 ottobre 2026 condiziona la traiettoria.}
  Predizione pubblica:
  \(P(\text{vittoria Lula o successore di sinistra}) = 0{,}55\)-\(0{,}70\)
  aprile 2026, da rivedere al 30 giugno 2026 e al 30 settembre 2026
  sulla base di sondaggi aggregati (Genial/Quaest, AtlasIntel,
  Datafolha).
\item
  \textbf{Se Lula/sinistra vince}: TFFF raccoglie \(> 5\)B USD impegnati
  entro fine 2027 e BRICS Multilateral Guarantees parte operativamente
  entro Q2 2027. Falsificabile su comunicati BRICS e World Bank.
\item
  \textbf{V-Dem Brasile in salita}: indice Liberal Democracy oltre
  \(0{,}65\) entro report 2028 (su anno 2027), condizionale alla
  continuità Lula/successore. Falsificabile su \url{https://v-dem.net}.
\item
  \textbf{Se vince Flávio Bolsonaro o destra radicale}: deterioramento
  V-Dem \(> 0{,}05\) punti in 24 mesi; Brasile esce da \emph{Group of
  Friends for Peace}; partecipazione attiva BRICS si trasforma in
  adesione formale ma passiva. Falsificabile su corso degli eventi
  2027-2028.
\end{enumerate}



\section{Note Epistemiche di Chiusura
Capitolo}\label{note-epistemiche-di-chiusura-capitolo}

{\def\LTcaptype{none} % do not increment counter
{\small\begin{longtable}[]{@{}
  >{\raggedright\arraybackslash}p{(\linewidth - 4\tabcolsep) * \real{0.3333}}
  >{\raggedright\arraybackslash}p{(\linewidth - 4\tabcolsep) * \real{0.3333}}
  >{\raggedright\arraybackslash}p{(\linewidth - 4\tabcolsep) * \real{0.3333}}@{}}
\toprule\noalign{}
\begin{minipage}[b]{\linewidth}\raggedright
\#
\end{minipage} & \begin{minipage}[b]{\linewidth}\raggedright
Affermazione
\end{minipage} & \begin{minipage}[b]{\linewidth}\raggedright
Status
\end{minipage} \\
\midrule\noalign{}
\endhead
\bottomrule\noalign{}
\endlastfoot
1 & Definizione formale Bilanciere Quantistico & \Hypoth{} categoria
DEQ$^2$ nuova \\
2 & Estado Logístico (Cervo) come \(\Phi^{\text{esportabile}}\)
operativo & \Hypoth{} mappatura concettuale forte \\
3 & Doppia presidenza BRICS Rio + COP30 Belém 2025 come test simultaneo
& \Proved{} dato pubblicato \\
4 & Mediazione RU-UA via Group of Friends for Peace & \Proved{} dato
pubblicato \\
5 & V-Dem 2026 U-turn per Brasile & \Proved{} dato pubblicato \\
6 & Condanna Bolsonaro 27 anni come prova di \(R_c\) & \Proved{} dato
pubblicato 11/09/2025 \\
7 & Approval Lula in calo, fragilità reali & \Proved{} dati
Genial/Quaest, Bloomberg \\
8 & Selic \(15\%\), deficit \(\approx 3\%\) PIL 2026, debito
\(\approx 75\%\) & \Proved{} dati BBVA, OECD \\
9 & Sondaggio AtlasIntel: Flávio \(47{,}6\%\) vs Lula \(46{,}6\%\)
secondo turno & \Proved{} dato aprile 2026 \\
10 & Tre scenari elettorali con probabilità
\(A \approx 0{,}45\)-\(0{,}55\), \(B \approx 0{,}10\)-\(0{,}15\),
\(C \approx 0{,}30\)-\(0{,}40\) & \Conj{} stima soggettiva da sondaggi
aggregati \\
11 & \(\Phi\)-engine condizionale all'esito elettorale & \Hypoth{}
risultato strutturale forte \\
12 & Predizioni 1-4 (TFFF, BRICS BMG, V-Dem 2027) & \Conj{}
falsificabili \\
\end{longtable}}
}



\section{Fonti dei dati 2025-2026
utilizzate}\label{fonti-dei-dati-2025-2026-utilizzate}

\begin{itemize}
\tightlist
\item
  Cervo, A. L., \emph{Inserção internacional: formação dos conceitos
  brasileiros} (2008) --- base teorica Estado Logístico
\item
  BRICS Brasil --- 17$^\circ$ summit Rio 6-7/07/2025, Dichiarazione 126 punti,
  TFFF, BMG
\item
  COP30 Belém 10-22/11/2025 --- UN, UNFCCC, Carbon Brief, World
  Resources Institute, House of Commons Library
\item
  V-Dem 2026 --- U-turn Brasile post-2022, \emph{judicial assertiveness}
\item
  STF Brasile --- condanna Bolsonaro 11/09/2025, conferma appello
  7/11/2025 (Al Jazeera, PBS, Time, CNN)
\item
  Genial/Quaest, AtlasIntel/Bloomberg, Datafolha --- sondaggi 2026
\item
  BBVA Research, OECD, Deloitte --- Brazil Economic Outlook 2026
\item
  gov.br/planalto --- Lula visita Mosca 05/2025, Group of Friends for
  Peace
\item
  Brand Finance Soft Power Index 2026 --- Brasile \(49{,}2\)
\item
  Munich Security Report 2025 --- Cap. 8 \emph{Lula Land}
\end{itemize}



\section{Chiusura della Parte III e Apertura della Parte
IV}\label{chiusura-della-parte-iii-e-apertura-della-parte-iv}

I quattro capitoli applicativi della Parte III (Solitone, Egemone
Decoerente, Rigido Sincronizzato, Bilanciere) hanno mappato il sistema
globale 2026 sull'apparato DEQ$^2$. Tre dei quattro attori hanno
traiettorie già definite --- il Solitone verso una metastabilità
esauribile entro 2031, l'Egemone verso \(\Phi \to 0\) distribuito entro
2027-2029, il Rigido verso \(\Phi\) catastrofica entro l'orizzonte
demografico 2046 ma con punti di rottura possibili nel triennio
2027-2029. Il quarto --- il Bilanciere brasiliano --- è in
\emph{biforcazione esplicita} con data e meccanismo noti: 4 ottobre
2026.

La \textbf{Parte IV} del libro tratterà la transizione di fase 2028-2031
considerando le quattro traiettorie come funzioni del tempo accoppiate,
con dashboard SPC-DEQ (Cap. 10) e predizioni datate falsificabili (Cap.
11).




%% ═══════════════════════════════════════════════════════════════════════════
%% PARTE IV — La Transizione di Fase 2028--2031
%% ═══════════════════════════════════════════════════════════════════════════
\part[La Transizione di Fase 2028--2031]{La Transizione di Fase\\2028--2031}

\chapter{Scenari Quantizzati e Dashboard
SPC-DEQ}\label{scenari-quantizzati-e-dashboard-spc-deq}

\begin{quote}
\textbf{Argomento portante:} la DEQ$^2$ è una teoria \emph{operativa}
perché ammette una traduzione diretta in \textbf{Statistical Process
Control adattato ai sistemi sociali}. La Dashboard SPC-DEQ è la
proiezione monitorabile di \(\boldsymbol{p}(t) \in \mathbb{V}^{(4)}\) in
\emph{sei pannelli bivariati} (uno per ciascuna relazione R1-R6) più un
\emph{pannello aggregato} della distanza scalare
\(d(\boldsymbol{p}(t), \mathcal{S})\). Le soglie operative non sono
regole arbitrarie ma derivano dalla geometria della superficie critica
\(\mathcal{S}\) e dalla varianza naturale degli indicatori. I tre attori
principali (USA, Cina, Brasile) sono mappati in scenari quantizzati
2026-2031, con uno scenario aggregato e due contro-fattuali brasiliani
che dipendono dall'esito elettorale del 4 ottobre 2026.
\end{quote}



\section{Perché SPC}\label{perchuxe9-spc}

\Hypoth{} \textbf{Tradizione tecnica.} Lo \emph{Statistical Process
Control} (Walter Shewhart, 1924; Edward Deming, 1950s; Six Sigma 1980s)
è la disciplina che monitora processi industriali tramite \emph{carte di
controllo} con limiti \(UCL\) (upper) e \(LCL\) (lower) calibrati sulla
varianza naturale del processo. Le \emph{regole di Western Electric}
(1956) identificano pattern di derive sistematiche prima che il processo
esca dai limiti.

\Hypoth{} \textbf{Adattamento alla DEQ$^2$.} Il sistema sociale è
formalmente analogo a un processo industriale con due differenze:

{\def\LTcaptype{none} % do not increment counter
{\small\begin{longtable}[]{@{}
  >{\raggedright\arraybackslash}p{(\linewidth - 2\tabcolsep) * \real{0.5000}}
  >{\raggedright\arraybackslash}p{(\linewidth - 2\tabcolsep) * \real{0.5000}}@{}}
\toprule\noalign{}
\begin{minipage}[b]{\linewidth}\raggedright
Differenza SPC industriale $\to$ SPC-DEQ$^2$
\end{minipage} & \begin{minipage}[b]{\linewidth}\raggedright
Adattamento richiesto
\end{minipage} \\
\midrule\noalign{}
\endhead
\bottomrule\noalign{}
\endlastfoot
Variabili continue univariate $\to$ variabili strutturali
\(\mathbb{V}^{(4)}\) multivariate & Sei pannelli bivariati (R1-R6) +
pannello aggregato \(d(\boldsymbol{p}, \mathcal{S})\) \\
Limiti \(UCL/LCL\) simmetrici $\to$ superficie critica \(\mathcal{S}\)
asimmetrica & \(\mathcal{S}\) ricavata analiticamente dal Cap. 5.3, non
dalla varianza \\
Causa speciale = fuori-limite $\to$ transizione di fase =
\(T_s^{(2),\text{sys}} < 1\) persistente & Soglia di persistenza
\(\Delta t > \Delta t^*\) (Cap. 3.5.3) \\
Capability index \(C_p = (USL-LSL)/6\sigma\) $\to$ indice DEQ$^2$
\(C_{\text{DEQ}}^2 := d(\boldsymbol{p}, \mathcal{S}) / \sigma_d\) &
Quantità scalare confrontabile fra sistemi \\
\end{longtable}}
}

\Proved{} \textbf{Conseguenza operativa.} La Dashboard SPC-DEQ è
\emph{azione politica formalizzata}: ogni leva del Cap. 5.5 (gradi di
libertà operativi) corrisponde a un intervento mirato su uno o più
pannelli. Non è strumento di previsione passiva --- è \emph{interfaccia
di pilotaggio} del punto \(\boldsymbol{p}(t)\) in
\(\mathbb{V}^{(4)}_{\text{ammissibile}}\).



\section{I Sei Pannelli Bivariati}\label{i-sei-pannelli-bivariati}

\Hypoth{} \textbf{Costruzione.} Per ciascuna relazione \(R_i\) del Cap.
5.1, costruiamo un \emph{pannello bidimensionale} sulle due coordinate
accoppiate, in cui:

\begin{itemize}
\tightlist
\item
  l'asse \(x\) è la prima variabile della coppia;
\item
  l'asse \(y\) è la seconda;
\item
  una \emph{iso-curva} \(\mathcal{S}_i\) rappresenta la sezione di
  \(\mathcal{S}\) a \((s_{iii}, z, \delta, \sigma)\) esogeni fissi e
  altre due variabili al loro valore mediano sistemico;
\item
  il punto attuale \(\boldsymbol{p}(t)\) è proiettato e tracciato come
  \emph{trail temporale} (ultimi 12-24 mesi);
\item
  una \emph{zona di near-miss} (banda colorata tra \(\mathcal{S}_i\) e
  \(0{,}9 \cdot \mathcal{S}_i\) in distanza euclidea normalizzata)
  segnala l'avvicinamento alla criticità.
\end{itemize}

\Proved{} \textbf{I sei pannelli con interpretazione operativa.}

{\def\LTcaptype{none} % do not increment counter
{\small\begin{longtable}[]{@{}
  >{\raggedright\arraybackslash}p{(\linewidth - 8\tabcolsep) * \real{0.2000}}
  >{\raggedright\arraybackslash}p{(\linewidth - 8\tabcolsep) * \real{0.2000}}
  >{\raggedright\arraybackslash}p{(\linewidth - 8\tabcolsep) * \real{0.2000}}
  >{\raggedright\arraybackslash}p{(\linewidth - 8\tabcolsep) * \real{0.2000}}
  >{\raggedright\arraybackslash}p{(\linewidth - 8\tabcolsep) * \real{0.2000}}@{}}
\toprule\noalign{}
\begin{minipage}[b]{\linewidth}\raggedright
Pannello
\end{minipage} & \begin{minipage}[b]{\linewidth}\raggedright
Coppia
\end{minipage} & \begin{minipage}[b]{\linewidth}\raggedright
Iso-curva
\end{minipage} & \begin{minipage}[b]{\linewidth}\raggedright
Direzione di sicurezza
\end{minipage} & \begin{minipage}[b]{\linewidth}\raggedright
Allarme tipico
\end{minipage} \\
\midrule\noalign{}
\endhead
\bottomrule\noalign{}
\endlastfoot
\textbf{P1} & \(\Phi\)-\(A\) &
\(A = (\Phi^{-1} - 1)/\langle A\rangle_{\text{ref}}\) & NE (alto
\(\Phi\), alto \(A\)) & \(\boldsymbol{p}\) scende verso SW (decoerenza
fragilizzante) \\
\textbf{P2} & \(\Phi\)-\(\Omega\) & \(\Phi = \Omega/\Omega_c\) con
\(\Omega_c\) critica & NW (alto \(\Phi\), basso \(\Omega\)) &
\(\boldsymbol{p}\) migra verso SE (decoerenza con dissipazione) \\
\textbf{P3} & \(\Phi\)-\(\varepsilon_{\text{dis}}\) &
\(\Phi = 1 - \varepsilon_{\text{dis}}/\varepsilon_c\) & NW (alto
\(\Phi\), basso \(\varepsilon\)) & \(\boldsymbol{p}\) verso SE
(decoerenza con disimpegno) \\
\textbf{P4} & \(A\)-\(\Omega\) & \(A = -\Omega \cdot k_4\) (\(k_4 > 0\))
& NW (alto \(A\), basso \(\Omega\)) & \(\boldsymbol{p}\) verso SE
(fragilità con dissipazione --- pattern \emph{cinese}) \\
\textbf{P5} & \(A\)-\(\varepsilon_{\text{dis}}\) &
\(A = (1 - \varepsilon)/k_5\) & NW (alto \(A\), basso \(\varepsilon\)) &
\(\boldsymbol{p}\) verso SE (fragilità con disimpegno --- pattern
\emph{USA}) \\
\textbf{P6} & \(\Omega\)-\(\varepsilon_{\text{dis}}\) &
\(\varepsilon_{\text{dis}} = k_6 \cdot \Omega\) & SW (basso \(\Omega\),
basso \(\varepsilon\)) & \(\boldsymbol{p}\) verso NE (dissipazione che
alimenta disimpegno --- pattern \emph{involutivo}) \\
\end{longtable}}
}

\Hypoth{} \textbf{Pannello aggregato \(P_0\).} La distanza scalare
\(d(\boldsymbol{p}(t), \mathcal{S})\) è tracciata su un'unica carta di
controllo con:

\begin{itemize}
\tightlist
\item
  \emph{zona verde}: \(d > 0{,}30\) (margine di sicurezza ampio);
\item
  \emph{zona gialla}: \(0 < d \leq 0{,}30\) (vigilanza);
\item
  \emph{zona rossa}: \(-0{,}20 < d \leq 0\) (sotto soglia ma
  reversibile);
\item
  \emph{zona nera}: \(d \leq -0{,}20\) (transizione di fase consolidata,
  \(\Delta t > \Delta t^*\) richiesto).
\end{itemize}



\section{Soglie Operative e Segnali di Allarme
Precoce}\label{soglie-operative-e-segnali-di-allarme-precoce}

\Hypoth{} \textbf{Adattamento delle regole di Western Electric a
\(\mathbb{V}^{(4)}\).} La logica WE classica si traduce in cinque
pattern DEQ$^2$:

{\def\LTcaptype{none} % do not increment counter
{\small\begin{longtable}[]{@{}
  >{\raggedright\arraybackslash}p{(\linewidth - 4\tabcolsep) * \real{0.3333}}
  >{\raggedright\arraybackslash}p{(\linewidth - 4\tabcolsep) * \real{0.3333}}
  >{\raggedright\arraybackslash}p{(\linewidth - 4\tabcolsep) * \real{0.3333}}@{}}
\toprule\noalign{}
\begin{minipage}[b]{\linewidth}\raggedright
Pattern WE-DEQ$^2$
\end{minipage} & \begin{minipage}[b]{\linewidth}\raggedright
Definizione
\end{minipage} & \begin{minipage}[b]{\linewidth}\raggedright
Interpretazione
\end{minipage} \\
\midrule\noalign{}
\endhead
\bottomrule\noalign{}
\endlastfoot
\textbf{WE-1} & Un punto in zona rossa o nera & Allarme immediato ---
entrata in \(\mathcal{R}_-\) \\
\textbf{WE-2} & Due di tre punti consecutivi in zona gialla o peggio &
Deriva confermata --- intervenire entro un trimestre \\
\textbf{WE-3} & Quattro di cinque punti consecutivi in zona gialla &
Trend strutturale --- riallineare leve operative del Cap. 5.5 \\
\textbf{WE-4} & Otto punti consecutivi in derivata \(\dot d < 0\) &
Erosione lenta --- diagnosi delle relazioni R1-R6 più stressate \\
\textbf{WE-5} & Inversione di derivata seconda \(\ddot d < 0\) con
\(d > 0\) & Pre-allarme --- sistema ancora sicuro ma sta accelerando il
declino \\
\end{longtable}}
}

\Proved{} \textbf{Test sui dati 2024-2026 dei tre attori} (verifica
retrospettiva qualitativa):

{\def\LTcaptype{none} % do not increment counter
{\small\begin{longtable}[]{@{}
  >{\raggedright\arraybackslash}p{(\linewidth - 4\tabcolsep) * \real{0.3333}}
  >{\raggedright\arraybackslash}p{(\linewidth - 4\tabcolsep) * \real{0.3333}}
  >{\raggedright\arraybackslash}p{(\linewidth - 4\tabcolsep) * \real{0.3333}}@{}}
\toprule\noalign{}
\begin{minipage}[b]{\linewidth}\raggedright
Attore
\end{minipage} & \begin{minipage}[b]{\linewidth}\raggedright
Pattern attivati 2024-2026
\end{minipage} & \begin{minipage}[b]{\linewidth}\raggedright
Coerenza con Cap. 7-9
\end{minipage} \\
\midrule\noalign{}
\endhead
\bottomrule\noalign{}
\endlastfoot
USA & WE-1 (zona rossa, V-Dem \(0{,}57\)), WE-3 (Pew trust in declino 4
anni), WE-4 (erosione lenta su tutti i pannelli) & \Proved{} Cap. 7 \\
Cina & WE-4 su P2 (deflazione PPI 10 trim.), WE-2 su P4 (involution +
property) & \Proved{} Cap. 8 \\
Brasile & WE-5 su pannello aggregato (approval calante, ma \(d > 0\)) &
\Proved{} Cap. 9.5 \\
\end{longtable}}
}

\Hypoth{} \textbf{Conseguenza metodologica.} I tre attori esibiscono
pattern WE-DEQ$^2$ \emph{qualitativamente diversi} --- il che valida la
pertinenza diagnostica del cruscotto. La diagnosi non sostituisce
l'interpretazione strutturale dei Cap. 7-9 ma la \emph{quantifica} in
segnali operativi standard.



\section{Scenari Quantizzati per Attore:
2026-2031}\label{scenari-quantizzati-per-attore-2026-2031}

\Hypoth{} \textbf{Discretizzazione per scenari.} Anziché simulare il
sistema dinamico continuo del Cap. 5.4, adottiamo \emph{scenari
quantizzati}: tre traiettorie per attore (best/median/worst case)
parametrizzate dal valore di una variabile critica esogena (per gli USA:
\(\zeta_{\text{Sol}}(t)\); per la Cina: \(K^{\text{imposto}}\); per il
Brasile: esito elettorale 4/10/2026). Le tre traiettorie partono dal
punto stimato \(\boldsymbol{p}(2026)\) del Cap. 5.2 e si sviluppano in
sei tappe annuali.

\subsection{USA --- Scenari 2026-2031}\label{usa-scenari-2026-2031}

{\def\LTcaptype{none} % do not increment counter
{\small\begin{longtable}[]{@{}
  >{\raggedright\arraybackslash}p{(\linewidth - 6\tabcolsep) * \real{0.2500}}
  >{\raggedright\arraybackslash}p{(\linewidth - 6\tabcolsep) * \real{0.2500}}
  >{\raggedright\arraybackslash}p{(\linewidth - 6\tabcolsep) * \real{0.2500}}
  >{\raggedright\arraybackslash}p{(\linewidth - 6\tabcolsep) * \real{0.2500}}@{}}
\toprule\noalign{}
\begin{minipage}[b]{\linewidth}\raggedright
Anno
\end{minipage} & \begin{minipage}[b]{\linewidth}\raggedright
Scenario \emph{median} (probabilità \(\approx 0{,}50\))
\end{minipage} & \begin{minipage}[b]{\linewidth}\raggedright
Scenario \emph{worst} (\(\approx 0{,}30\))
\end{minipage} & \begin{minipage}[b]{\linewidth}\raggedright
Scenario \emph{best} (\(\approx 0{,}20\))
\end{minipage} \\
\midrule\noalign{}
\endhead
\bottomrule\noalign{}
\endlastfoot
2026 & \(d \approx -0{,}10\) (zona rossa) & \(d \approx -0{,}10\) &
\(d \approx -0{,}10\) \\
2027 & \(d \approx -0{,}15\) ; midterm split, \(\Omega \uparrow\) &
\(d \approx -0{,}25\) ; escalation costituzionale &
\(d \approx -0{,}05\) ; SCOTUS riequilibra \\
2028 & \(d \approx -0{,}18\) ; presidenziali, \(\Phi\) scende
ulteriormente & \(d \approx -0{,}35\) ; transizione di fase consolidata
& \(d \approx 0\) ; nuovo equilibrio post-elettorale \\
2029 & \(d \approx -0{,}20\) ; stato stazionario decoerente &
\(d \approx -0{,}40\) ; instabilità acuta & \(d \approx +0{,}05\) ;
recupero parziale \(s_{iii}\) \\
2030 & \(d \approx -0{,}22\) & \(d \approx -0{,}45\) ; crisi monetaria
possibile & \(d \approx +0{,}08\) \\
2031 & \(d \approx -0{,}25\) ; \emph{potenza materiale senza fase} &
\(d \approx -0{,}50\) ; rottura strutturale & \(d \approx +0{,}10\) ;
consolidamento del recupero \\
\end{longtable}}
}

\Conj{} \textbf{Predizione DEQ$^2$ mediana}:
\(T_s^{(2),\text{sys},\text{USA}} < 1\) in tutti gli scenari per
l'intero orizzonte 2026-2031. Recupero solo nello scenario best,
condizionale a SCOTUS attivo + nuova amministrazione \(\neq\) Trump 2.0.

\subsection{Cina --- Scenari 2026-2031}\label{cina-scenari-2026-2031}

{\def\LTcaptype{none} % do not increment counter
{\small\begin{longtable}[]{@{}
  >{\raggedright\arraybackslash}p{(\linewidth - 6\tabcolsep) * \real{0.2500}}
  >{\raggedright\arraybackslash}p{(\linewidth - 6\tabcolsep) * \real{0.2500}}
  >{\raggedright\arraybackslash}p{(\linewidth - 6\tabcolsep) * \real{0.2500}}
  >{\raggedright\arraybackslash}p{(\linewidth - 6\tabcolsep) * \real{0.2500}}@{}}
\toprule\noalign{}
\begin{minipage}[b]{\linewidth}\raggedright
Anno
\end{minipage} & \begin{minipage}[b]{\linewidth}\raggedright
Scenario \emph{median} (\(\approx 0{,}45\))
\end{minipage} & \begin{minipage}[b]{\linewidth}\raggedright
Scenario \emph{worst} (\(\approx 0{,}30\))
\end{minipage} & \begin{minipage}[b]{\linewidth}\raggedright
Scenario \emph{best} (\(\approx 0{,}25\))
\end{minipage} \\
\midrule\noalign{}
\endhead
\bottomrule\noalign{}
\endlastfoot
2026 & \(d^{\text{oss}} \approx +0{,}25\) ;
\(d^{\text{spont}} \approx -0{,}05\) & come median & come median \\
2027 & \(d^{\text{oss}}\) stabile ; \(d^{\text{spont}} \approx -0{,}10\)
; deflazione \(> 12\) trim. & \(d^{\text{oss}}\) scende a \(+0{,}10\) ;
tensione Taiwan acuta & \(d^{\text{oss}} \approx +0{,}28\) ;
anti-involution efficace, riforma proprietà \\
2028 & divergenza oss/spont si amplia & possibile evento \emph{bloqueo}
Taiwan & recupero TFP modesto \\
2029 & \(d^{\text{oss}}\) inizia a calare per \emph{fatica del K
imposto} & crisi politica interna & \(d^{\text{oss}}\) stabile,
\(d^{\text{spont}}\) recupera \\
2030 & \(d^{\text{oss}} \approx +0{,}10\) ;
\(d^{\text{spont}} \approx -0{,}20\) & rottura del cluster sincronizzato
& \(d^{\text{oss}} \approx +0{,}30\) \\
2031 & candidato a \emph{transizione catastrofica} nel triennio
successivo & transizione catastrofica avvenuta & metà-soft landing \\
\end{longtable}}
}

\Conj{} \textbf{Predizione DEQ$^2$ mediana}: divario crescente
\(d^{\text{oss}} - d^{\text{spont}}\), con \emph{fatica del \(K\)
imposto} visibile a partire dal 2029. La transizione di fase cinese ha
probabilità \(\approx 0{,}40\)-\(0{,}55\) di attivarsi nel triennio
2030-2032 --- \emph{fuori} dall'orizzonte centrale del libro ma dentro
alla coda di rilevanza.

\subsection{Brasile --- Scenari 2026-2031 con Biforcazione 4 ottobre
2026}\label{brasile-scenari-2026-2031-con-biforcazione-4-ottobre-2026}

\Hypoth{} \textbf{Schema doppio.} Per il Brasile gli scenari sono
\emph{condizionali} all'esito elettorale. Dividiamo in due rami
principali:

\textbf{Ramo A/B --- Vittoria Lula o successore di sinistra}
(probabilità aprile 2026 \(\approx 0{,}55\)-\(0{,}70\)):

{\def\LTcaptype{none} % do not increment counter
{\small\begin{longtable}[]{@{}
  >{\raggedright\arraybackslash}p{(\linewidth - 6\tabcolsep) * \real{0.2500}}
  >{\raggedright\arraybackslash}p{(\linewidth - 6\tabcolsep) * \real{0.2500}}
  >{\raggedright\arraybackslash}p{(\linewidth - 6\tabcolsep) * \real{0.2500}}
  >{\raggedright\arraybackslash}p{(\linewidth - 6\tabcolsep) * \real{0.2500}}@{}}
\toprule\noalign{}
\begin{minipage}[b]{\linewidth}\raggedright
Anno
\end{minipage} & \begin{minipage}[b]{\linewidth}\raggedright
Scenario \emph{median}
\end{minipage} & \begin{minipage}[b]{\linewidth}\raggedright
\emph{Worst-A/B}
\end{minipage} & \begin{minipage}[b]{\linewidth}\raggedright
\emph{Best-A/B}
\end{minipage} \\
\midrule\noalign{}
\endhead
\bottomrule\noalign{}
\endlastfoot
2027 & \(d \approx +0{,}12\) ; consolidamento istituzionale &
\(d \approx +0{,}05\) ; fiscale teso & \(d \approx +0{,}20\) ; TFFF
parte forte \\
2028 & \(d \approx +0{,}15\) ; BRICS BMG operativo &
\(d \approx +0{,}08\) & \(d \approx +0{,}25\) ; \emph{$\Phi$-engine}
riconosciuto \\
2029 & \(d \approx +0{,}18\) & \(d \approx +0{,}10\) &
\(d \approx +0{,}28\) \\
2030 & \(d \approx +0{,}20\) ; presidenziali fisiologiche &
\(d \approx +0{,}05\) ; alternanza stretta & \(d \approx +0{,}30\) \\
2031 & \(d \approx +0{,}22\) ; \emph{$\Phi$-engine} consolidato &
\(d \approx +0{,}10\) ; condizionale alla nuova amministrazione &
\(d \approx +0{,}32\) ; leadership globale Sud \\
\end{longtable}}
}

\textbf{Ramo C --- Vittoria Flávio Bolsonaro o destra radicale}
(probabilità \(\approx 0{,}30\)-\(0{,}40\)):

{\def\LTcaptype{none} % do not increment counter
{\small\begin{longtable}[]{@{}
  >{\raggedright\arraybackslash}p{(\linewidth - 6\tabcolsep) * \real{0.2500}}
  >{\raggedright\arraybackslash}p{(\linewidth - 6\tabcolsep) * \real{0.2500}}
  >{\raggedright\arraybackslash}p{(\linewidth - 6\tabcolsep) * \real{0.2500}}
  >{\raggedright\arraybackslash}p{(\linewidth - 6\tabcolsep) * \real{0.2500}}@{}}
\toprule\noalign{}
\begin{minipage}[b]{\linewidth}\raggedright
Anno
\end{minipage} & \begin{minipage}[b]{\linewidth}\raggedright
Scenario \emph{median-C}
\end{minipage} & \begin{minipage}[b]{\linewidth}\raggedright
\emph{Worst-C}
\end{minipage} & \begin{minipage}[b]{\linewidth}\raggedright
\emph{Best-C}
\end{minipage} \\
\midrule\noalign{}
\endhead
\bottomrule\noalign{}
\endlastfoot
2027 & \(d \approx +0{,}05\) ; primi 100 giorni transizione &
\(d \approx -0{,}05\) ; smantellamento STF tentato &
\(d \approx +0{,}10\) ; vincoli istituzionali tengono \\
2028 & \(d \approx -0{,}05\) ; uscita Group of Friends; V-Dem cala
\(0{,}05\) & \(d \approx -0{,}15\) ; rottura BRICS &
\(d \approx +0{,}08\) ; resilienza istituzionale \\
2029 & \(d \approx -0{,}10\) ; \emph{solitonizzazione discorsiva} &
\(d \approx -0{,}25\) ; pattern USA-like & \(d \approx +0{,}05\) \\
2030 & \(d \approx -0{,}12\) & \(d \approx -0{,}30\) & \(d \approx 0\) ;
pareggio precario \\
2031 & \(d \approx -0{,}15\) ; \emph{$\Phi$-engine globale assente} &
\(d \approx -0{,}35\) ; transizione di fase regionale &
\(d \approx -0{,}05\) ; mantenimento basale \\
\end{longtable}}
}

\Hypoth{} \textbf{Tesi forte.} Il \textbf{divario fra Ramo A/B e Ramo C}
nel 2031 è \(\Delta d \approx 0{,}30\)-\(0{,}40\) --- il più ampio fra
qualunque variabile considerata in tutto il libro. È la \textbf{misura
quantitativa dell'importanza dell'elezione brasiliana} del 4 ottobre
2026 per l'intero sistema globale.



\section{Lo Scenario Aggregato 2031: Quattro Mondi
Possibili}\label{lo-scenario-aggregato-2031-quattro-mondi-possibili}

\Hypoth{} \textbf{Costruzione.} Combinando gli scenari mediani per
attore con i due rami brasiliani, otteniamo quattro \emph{mondi
possibili} al 2031, ciascuno caratterizzato da una distribuzione
\(\{d_{\text{USA}}, d_{\text{Cina}}, d_{\text{Brasile}}\}\):

{\def\LTcaptype{none} % do not increment counter
{\small\begin{longtable}[]{@{}
  >{\raggedright\arraybackslash}p{(\linewidth - 10\tabcolsep) * \real{0.1667}}
  >{\raggedright\arraybackslash}p{(\linewidth - 10\tabcolsep) * \real{0.1667}}
  >{\raggedright\arraybackslash}p{(\linewidth - 10\tabcolsep) * \real{0.1667}}
  >{\raggedright\arraybackslash}p{(\linewidth - 10\tabcolsep) * \real{0.1667}}
  >{\raggedright\arraybackslash}p{(\linewidth - 10\tabcolsep) * \real{0.1667}}
  >{\raggedright\arraybackslash}p{(\linewidth - 10\tabcolsep) * \real{0.1667}}@{}}
\toprule\noalign{}
\begin{minipage}[b]{\linewidth}\raggedright
Mondo
\end{minipage} & \begin{minipage}[b]{\linewidth}\raggedright
USA
\end{minipage} & \begin{minipage}[b]{\linewidth}\raggedright
Cina
\end{minipage} & \begin{minipage}[b]{\linewidth}\raggedright
Brasile
\end{minipage} & \begin{minipage}[b]{\linewidth}\raggedright
Probabilità soggettiva
\end{minipage} & \begin{minipage}[b]{\linewidth}\raggedright
Caratteristica strutturale
\end{minipage} \\
\midrule\noalign{}
\endhead
\bottomrule\noalign{}
\endlastfoot
\textbf{M1: Decoerenza globale con \(\Phi\)-engine attivo} & \(-0{,}25\)
& \(+0{,}10/-0{,}20\) & \(+0{,}22\) (A/B) &
\(\approx 0{,}35\)-\(0{,}40\) & Brasile come polo di compensazione del
declino USA-Cina \\
\textbf{M2: Decoerenza globale senza \(\Phi\)-engine} & \(-0{,}25\) &
\(+0{,}10/-0{,}20\) & \(-0{,}15\) (C) & \(\approx 0{,}20\)-\(0{,}25\) &
Multipolarità degenere --- frammentazione senza centro coerente \\
\textbf{M3: Rottura cinese precoce + \(\Phi\)-engine attivo} &
\(-0{,}25\) & \(-0{,}30\) (worst) & \(+0{,}22\) (A/B) &
\(\approx 0{,}12\)-\(0{,}18\) & Riassetto sistemico veloce, Brasile
riferimento \\
\textbf{M4: Recupero USA + Cina soft landing + \(\Phi\)-engine attivo} &
\(+0{,}10\) (best) & \(+0{,}30\) (best) & \(+0{,}22\) (A/B) &
\(\approx 0{,}05\)-\(0{,}10\) & Stabilizzazione plurale --- scenario
meno probabile ma operativamente più stabile \\
\end{longtable}}
}

\Hypoth{} \textbf{Lettura.} I mondi M1 e M2 differiscono \emph{solo} per
il ramo brasiliano. La loro probabilità combinata
(\(\approx 0{,}55\)-\(0{,}65\)) misura la \emph{probabilità che il
sistema globale 2031 sia in regime di decoerenza generalizzata}; la
presenza o assenza di un \(\Phi\)-engine ne determina la
\emph{governabilità}. La Parte V del libro (Cap. 12-13) costruisce
protocolli antifragili condizionati alla distinzione M1/M2.

\Proved{} \textbf{Risultato sintetico.}
\(P(\text{sistema globale 2031 in regime di decoerenza}) \approx 0{,}80\)-\(0{,}85\).
\(P(\text{sistema globale 2031 con almeno un $\Phi$-engine}) \approx 0{,}55\)-\(0{,}70\).
La differenza misura l'\textbf{ampiezza dello spazio di manovra residuo
per il Sud globale e per le strategie del Cap. 12}.



\section{\texorpdfstring{Indice \(C_{\text{DEQ}}^2\) e Confronto
Cross-Attore}{Indice C\_\{\textbackslash text\{DEQ\}\}\^{}2 e Confronto Cross-Attore}}\label{indice-c_textdeq2-e-confronto-cross-attore}

\Hypoth{} \textbf{Definizione operativa.}

\[
C_{\text{DEQ}}^2(\boldsymbol{p}, t) := \frac{d(\boldsymbol{p}(t), \mathcal{S})}{\sigma_d(t)}
\]

dove \(\sigma_d(t)\) è la deviazione standard naturale della distanza
calcolata su una finestra mobile di 36 mesi precedenti \(t\).

\Proved{} \textbf{Interpretazione.} \(C_{\text{DEQ}}^2 \geq 1\): il
sistema è a oltre una deviazione standard dalla soglia critica ---
\emph{capability} sistemica accettabile. \(C_{\text{DEQ}}^2 \leq 0\): il
sistema è dentro \(\mathcal{R}_-\). Il valore tipico desiderato per
sistemi industriali è \(C_p \geq 1{,}33\); per sistemi sociali, anche
per la maggiore stochasticity intrinseca, una soglia operativa
ragionevole è \(C_{\text{DEQ}}^2 \geq 1{,}0\).

{\def\LTcaptype{none} % do not increment counter
{\small\begin{longtable}[]{@{}
  >{\raggedright\arraybackslash}p{(\linewidth - 4\tabcolsep) * \real{0.3333}}
  >{\raggedright\arraybackslash}p{(\linewidth - 4\tabcolsep) * \real{0.3333}}
  >{\raggedright\arraybackslash}p{(\linewidth - 4\tabcolsep) * \real{0.3333}}@{}}
\toprule\noalign{}
\begin{minipage}[b]{\linewidth}\raggedright
Attore
\end{minipage} & \begin{minipage}[b]{\linewidth}\raggedright
\(C_{\text{DEQ}}^2\) stimato 2026
\end{minipage} & \begin{minipage}[b]{\linewidth}\raggedright
Lettura
\end{minipage} \\
\midrule\noalign{}
\endhead
\bottomrule\noalign{}
\endlastfoot
USA & \(\approx -0{,}5\) & sotto soglia, dentro \(\mathcal{R}_-\) \\
Cina & \(\approx +1{,}1\) apparente / \(\approx +0{,}3\) effettivo &
apparenza stabile, struttura fragile \\
Brasile & \(\approx +1{,}3\) & unica configurazione
\emph{industrial-grade} \\
\end{longtable}}
}

\Hypoth{} \textbf{Conseguenza.} Il Brasile è l'unico attore della Parte
III a soddisfare il criterio di \emph{capability industriale} applicato
ai sistemi sociali --- coerente con la definizione di \(\Phi\)-engine
del Cap. 9.



\section{Sintesi del Capitolo}\label{sintesi-del-capitolo}

\Proved{} \textbf{In una frase:} la Dashboard SPC-DEQ traduce
\(\boldsymbol{p}(t) \in \mathbb{V}^{(4)}\) in sei pannelli bivariati
monitorabili, un pannello aggregato \(d(\boldsymbol{p}, \mathcal{S})\),
cinque regole di pattern detection adattate da Western Electric, e un
indice scalare \(C_{\text{DEQ}}^2\) --- fornendo l'interfaccia operativa
per le strategie del Cap. 12 e il quadro evidenziale per le predizioni
del Cap. 11. I quattro mondi possibili al 2031 (M1-M4) consolidano
l'analisi della Parte III in \emph{scenari probabilizzati} con
probabilità soggettive esplicite.

{\def\LTcaptype{none} % do not increment counter
{\small\begin{longtable}[]{@{}
  >{\raggedright\arraybackslash}p{(\linewidth - 4\tabcolsep) * \real{0.3333}}
  >{\raggedright\arraybackslash}p{(\linewidth - 4\tabcolsep) * \real{0.3333}}
  >{\raggedright\arraybackslash}p{(\linewidth - 4\tabcolsep) * \real{0.3333}}@{}}
\toprule\noalign{}
\begin{minipage}[b]{\linewidth}\raggedright
Strumento
\end{minipage} & \begin{minipage}[b]{\linewidth}\raggedright
Funzione
\end{minipage} & \begin{minipage}[b]{\linewidth}\raggedright
Capitolo che lo usa
\end{minipage} \\
\midrule\noalign{}
\endhead
\bottomrule\noalign{}
\endlastfoot
Sei pannelli bivariati P1-P6 & Diagnosi delle relazioni R1-R6 più
stressate & Cap. 12 (interventi mirati) \\
Pannello aggregato \(P_0 = d(\boldsymbol{p}, \mathcal{S})\) & Indicatore
scalare di stabilità & Cap. 11 (predizioni) \\
Cinque regole WE-DEQ$^2$ & Pattern detection precoce & Cap. 12 (allarme
operativo) \\
Indice \(C_{\text{DEQ}}^2\) & Confronto cross-attore & Cap. 11 (sintesi
quantitativa) \\
Scenari quantizzati per attore & Proiezioni 2026-2031 & Cap. 11
(predizioni datate) \\
Quattro mondi possibili 2031 & Scenari aggregati & Cap. 12-13 (strategie
condizionali) \\
\end{longtable}}
}

\Hypoth{} \textbf{Tesi conclusiva.} La DEQ$^2$ è ora \emph{operativa}. La
Parte IV chiude la fase teorica del libro e apre la fase normativa: dato
lo spazio degli scenari, \emph{quali strategie sono antifragili in
ciascuno di essi?} Questo è l'oggetto della Parte V.



\section{Note Epistemiche di Chiusura
Capitolo}\label{note-epistemiche-di-chiusura-capitolo}

{\def\LTcaptype{none} % do not increment counter
{\small\begin{longtable}[]{@{}
  >{\raggedright\arraybackslash}p{(\linewidth - 4\tabcolsep) * \real{0.3333}}
  >{\raggedright\arraybackslash}p{(\linewidth - 4\tabcolsep) * \real{0.3333}}
  >{\raggedright\arraybackslash}p{(\linewidth - 4\tabcolsep) * \real{0.3333}}@{}}
\toprule\noalign{}
\begin{minipage}[b]{\linewidth}\raggedright
\#
\end{minipage} & \begin{minipage}[b]{\linewidth}\raggedright
Affermazione
\end{minipage} & \begin{minipage}[b]{\linewidth}\raggedright
Status
\end{minipage} \\
\midrule\noalign{}
\endhead
\bottomrule\noalign{}
\endlastfoot
1 & Adattamento SPC industriale $\to$ DEQ$^2$ & \Hypoth{} traduzione
strutturale rigorosa \\
2 & Sei pannelli bivariati con iso-curve esplicite & \Hypoth{}
costruzione operativa \\
3 & Cinque regole WE-DEQ$^2$ & \Hypoth{} adattamento di Western Electric
(1956) \\
4 & Test retrospettivo 2024-2026 sui tre attori & \Proved{} pattern
coerenti con Cap. 7-9 \\
5 & Scenari quantizzati USA con \(d \to -0{,}25\) entro 2031 (median) &
\Conj{} congettura datata, derivata da Cap. 7 \\
6 & Scenari Cina con divergenza \(d^{\text{oss}}/d^{\text{spont}}\)
crescente & \Conj{} congettura, derivata da Cap. 8 \\
7 & Doppio ramo Brasile (A/B vs C) con
\(\Delta d \approx 0{,}30\)-\(0{,}40\) al 2031 & \Conj{} quantificazione
della biforcazione del Cap. 9.5 \\
8 & Quattro mondi possibili al 2031 con probabilità soggettive & \Conj{}
stime soggettive su evidenza Cap. 7-9 \\
9 & \(P(\text{decoerenza globale 2031}) \approx 0{,}80\)-\(0{,}85\) &
\Conj{} sintesi probabilistica \\
10 & Indice \(C_{\text{DEQ}}^2\) + soglia operativa \(\geq 1{,}0\) &
\Hypoth{} adattamento di \(C_p\) industriale \\
11 & Solo Brasile soddisfa \(C_{\text{DEQ}}^2 \geq 1{,}0\) & \Hypoth{}
conseguenza derivata coerente con Cap. 9 \\
\end{longtable}}
}



\chapter{Predizioni Rischiose Datate: Falsificabilità
Pubblica}\label{predizioni-rischiose-datate-falsificabilituxe0-pubblica}

\begin{quote}
\textbf{Argomento portante:} una teoria sociale che non si espone alla
falsificazione \emph{datata} è indistinguibile da una narrazione. Questo
capitolo aggrega le sedici predizioni rischiose della DEQ$^2$ --- tredici
derivate dai Cap. 7-9 sui tre attori sistemici (cinque USA, quattro
Cina, quattro Brasile), tre derivate dal Cap. 10 sull'aggregato globale
--- in un \emph{pre-registered prediction set} con soglie numeriche,
date di scadenza, fonti dati pubbliche e protocollo di valutazione. Per
ciascuna predizione è inoltre specificato \emph{cosa accadrebbe alla
teoria se fallisse}: la falsificazione di alcune predizioni è gestibile
come correzione locale, quella di altre forzerebbe la riformulazione
strutturale dell'apparato.
\end{quote}



\section{Criteri DEQ$^2$ per una Predizione
Rischiosa}\label{criteri-dequxb2-per-una-predizione-rischiosa}

\Hypoth{} \textbf{Eredità popperiana operativizzata.} Una predizione
\(\mathcal{P}\) è \emph{rischiosa} nel senso DEQ$^2$ se soddisfa
simultaneamente:

\begin{enumerate}
\def\labelenumi{\arabic{enumi}.}
\tightlist
\item
  \textbf{Datazione esplicita}: esiste \(t^*\) entro cui \(\mathcal{P}\)
  è valutabile;
\item
  \textbf{Soglia numerica}: esiste un valore di confronto \(v^*\) tale
  che \(\mathcal{P}\) asserisce \(X(t^*) \lessgtr v^*\);
\item
  \textbf{Fonte dati pubblica}: esiste una sorgente \(\mathcal{D}\)
  verificabile da terzi indipendenti;
\item
  \textbf{Probabilità soggettiva esplicita}:
  \(\Pr(\mathcal{P} \text{ vera}) = p\) è dichiarata, \(p \in (0, 1)\);
\item
  \textbf{Conseguenza teorica}: è dichiarato cosa accade alla DEQ$^2$ se
  \(\mathcal{P}\) fallisce --- \emph{correzione locale},
  \emph{aggiustamento parametrico}, oppure \emph{riformulazione
  strutturale}.
\end{enumerate}

\Proved{} \textbf{Risultato metodologico.} Una predizione che soddisfa
(1)-(5) è \emph{Bayes-aggiornabile}: la sua valutazione produce un
aggiornamento esplicito delle probabilità posteriori sui parametri della
teoria. La DEQ$^2$ si presenta come \emph{modello probabilistico vivente},
non come dottrina chiusa.



\section{Le Tredici Predizioni dei Cap.
7-9}\label{le-tredici-predizioni-dei-cap.-7-9}

\subsection{USA --- Egemone Decoerente (Cap.
7)}\label{usa-egemone-decoerente-cap.-7}

{\def\LTcaptype{none} % do not increment counter
{\small\begin{longtable}[]{@{}
  >{\raggedright\arraybackslash}p{(\linewidth - 12\tabcolsep) * \real{0.1429}}
  >{\raggedright\arraybackslash}p{(\linewidth - 12\tabcolsep) * \real{0.1429}}
  >{\raggedright\arraybackslash}p{(\linewidth - 12\tabcolsep) * \real{0.1429}}
  >{\raggedright\arraybackslash}p{(\linewidth - 12\tabcolsep) * \real{0.1429}}
  >{\raggedright\arraybackslash}p{(\linewidth - 12\tabcolsep) * \real{0.1429}}
  >{\raggedright\arraybackslash}p{(\linewidth - 12\tabcolsep) * \real{0.1429}}
  >{\raggedright\arraybackslash}p{(\linewidth - 12\tabcolsep) * \real{0.1429}}@{}}
\toprule\noalign{}
\begin{minipage}[b]{\linewidth}\raggedright
\#
\end{minipage} & \begin{minipage}[b]{\linewidth}\raggedright
Predizione
\end{minipage} & \begin{minipage}[b]{\linewidth}\raggedright
\(t^*\)
\end{minipage} & \begin{minipage}[b]{\linewidth}\raggedright
\(v^*\)
\end{minipage} & \begin{minipage}[b]{\linewidth}\raggedright
Fonte \(\mathcal{D}\)
\end{minipage} & \begin{minipage}[b]{\linewidth}\raggedright
\(p\)
\end{minipage} & \begin{minipage}[b]{\linewidth}\raggedright
Falsificazione
\end{minipage} \\
\midrule\noalign{}
\endhead
\bottomrule\noalign{}
\endlastfoot
\textbf{U1} & V-Dem Liberal Democracy Index USA \(< 0{,}55\) & report
aprile 2027 (anno 2026) & \(0{,}55\) & \url{https://v-dem.net} &
\(0{,}70\) & locale (parametro \(\Phi\)) \\
\textbf{U2} & Quota dollaro nelle riserve IMF COFER \(< 55\%\) & Q4 2027
& \(0{,}55\) & IMF COFER trimestrale & \(0{,}55\) & locale (parametro
\(\sigma\) esogena) \\
\textbf{U3} & Quota dollaro \(< 52\%\) & Q4 2030 & \(0{,}52\) & IMF
COFER trimestrale & \(0{,}50\) & locale (cumulativa con U2) \\
\textbf{U4} & Vetoes presidenziali \(> 15\) + emergency docket SCOTUS
\(> 30\) in 12 mesi & gennaio-luglio 2027, primo anno Congresso 119$^\circ$ &
\(15\) vetoes, \(30\) emergency & Federal Register + SCOTUSblog &
\(0{,}55\) & parametrica (\(\Omega\) acuta vs cronica) \\
\textbf{U5} & Correlazione Musk visibility--polarizzazione USA
\(r > 0{,}3\) & calibrazione 2022-2026 & \(r = 0{,}3\) & dataset
compilato (Cap. 7) & \(0{,}65\) & locale (modello
\(\zeta_{\text{Sol}}\)) \\
\end{longtable}}
}

\subsection{Cina --- Rigido Sincronizzato (Cap.
8)}\label{cina-rigido-sincronizzato-cap.-8}

{\def\LTcaptype{none} % do not increment counter
{\small\begin{longtable}[]{@{}
  >{\raggedright\arraybackslash}p{(\linewidth - 12\tabcolsep) * \real{0.1429}}
  >{\raggedright\arraybackslash}p{(\linewidth - 12\tabcolsep) * \real{0.1429}}
  >{\raggedright\arraybackslash}p{(\linewidth - 12\tabcolsep) * \real{0.1429}}
  >{\raggedright\arraybackslash}p{(\linewidth - 12\tabcolsep) * \real{0.1429}}
  >{\raggedright\arraybackslash}p{(\linewidth - 12\tabcolsep) * \real{0.1429}}
  >{\raggedright\arraybackslash}p{(\linewidth - 12\tabcolsep) * \real{0.1429}}
  >{\raggedright\arraybackslash}p{(\linewidth - 12\tabcolsep) * \real{0.1429}}@{}}
\toprule\noalign{}
\begin{minipage}[b]{\linewidth}\raggedright
\#
\end{minipage} & \begin{minipage}[b]{\linewidth}\raggedright
Predizione
\end{minipage} & \begin{minipage}[b]{\linewidth}\raggedright
\(t^*\)
\end{minipage} & \begin{minipage}[b]{\linewidth}\raggedright
\(v^*\)
\end{minipage} & \begin{minipage}[b]{\linewidth}\raggedright
Fonte \(\mathcal{D}\)
\end{minipage} & \begin{minipage}[b]{\linewidth}\raggedright
\(p\)
\end{minipage} & \begin{minipage}[b]{\linewidth}\raggedright
Falsificazione
\end{minipage} \\
\midrule\noalign{}
\endhead
\bottomrule\noalign{}
\endlastfoot
\textbf{C1} & Deflazione cinese \(> 12\) trimestri consecutivi & Q4 2026
& \(12\) trim. & NBS / PIIE / WB & \(0{,}65\) & locale (\(\Omega\)
misurato) \\
\textbf{C2} & \(P(\text{evento *bloqueo* Taiwan})\) in 2027-2029,
condizionale a (deflazione \(\geq 12\) trim.) \(\wedge\) (GDP \(< 4\%\))
\(\wedge\) (nascite \(< 5/1000\)) & 2027-2029 &
\(P \in [0{,}25, 0{,}40]\) & corso eventi + INSS / AEI / Crisis Group &
\(0{,}50\) entro intervallo & strutturale se \(P\) fuori intervallo \\
\textbf{C3} & Yuan in IMF COFER \(< 4\%\) nel 2027 e \(< 5\%\) nel 2030
& Q4 2027, Q4 2030 & \(4\%\), \(5\%\) & IMF COFER trimestrale &
\(0{,}60\) & locale (modello dedollarizzazione) \\
\textbf{C4} & PPI cinese in territorio negativo \(\geq 5\) anni cumulati
& Q4 2027 & \(5\) anni & NBS PPI mensile & \(0{,}70\) & locale
(anti-involution efficacia) \\
\end{longtable}}
}

\subsection{Brasile --- Bilanciere Quantistico (Cap.
9)}\label{brasile-bilanciere-quantistico-cap.-9}

{\def\LTcaptype{none} % do not increment counter
{\small\begin{longtable}[]{@{}
  >{\raggedright\arraybackslash}p{(\linewidth - 12\tabcolsep) * \real{0.1429}}
  >{\raggedright\arraybackslash}p{(\linewidth - 12\tabcolsep) * \real{0.1429}}
  >{\raggedright\arraybackslash}p{(\linewidth - 12\tabcolsep) * \real{0.1429}}
  >{\raggedright\arraybackslash}p{(\linewidth - 12\tabcolsep) * \real{0.1429}}
  >{\raggedright\arraybackslash}p{(\linewidth - 12\tabcolsep) * \real{0.1429}}
  >{\raggedright\arraybackslash}p{(\linewidth - 12\tabcolsep) * \real{0.1429}}
  >{\raggedright\arraybackslash}p{(\linewidth - 12\tabcolsep) * \real{0.1429}}@{}}
\toprule\noalign{}
\begin{minipage}[b]{\linewidth}\raggedright
\#
\end{minipage} & \begin{minipage}[b]{\linewidth}\raggedright
Predizione
\end{minipage} & \begin{minipage}[b]{\linewidth}\raggedright
\(t^*\)
\end{minipage} & \begin{minipage}[b]{\linewidth}\raggedright
\(v^*\)
\end{minipage} & \begin{minipage}[b]{\linewidth}\raggedright
Fonte \(\mathcal{D}\)
\end{minipage} & \begin{minipage}[b]{\linewidth}\raggedright
\(p\)
\end{minipage} & \begin{minipage}[b]{\linewidth}\raggedright
Falsificazione
\end{minipage} \\
\midrule\noalign{}
\endhead
\bottomrule\noalign{}
\endlastfoot
\textbf{B1} & \(P(\text{Lula o successore di sinistra vince elezioni})\)
aprile 2026 & aprile, giugno, settembre 2026 (revisione mobile) &
\(p \in [0{,}55, 0{,}70]\) & Genial/Quaest, AtlasIntel, Datafolha &
meta-stima soggettiva & aggiustamento parametrico al cambiare
evidenze \\
\textbf{B2} & Se A/B vince: TFFF impegni \(> 5\)B USD entro fine 2027 e
BRICS BMG operativo entro Q2 2027 & Q4 2027, Q2 2027 & \(5\)B,
\emph{operativo} & comunicati BRICS, World Bank & \(0{,}55\)
condizionale & locale (operatività iniziativa) \\
\textbf{B3} & Se A/B vince: V-Dem Brasile Liberal Democracy Index
\(> 0{,}65\) & report 2028 (anno 2027) & \(0{,}65\) &
\url{https://v-dem.net} & \(0{,}65\) condizionale & locale (\(\Phi\)
Brasile) \\
\textbf{B4} & Se C vince: deterioramento V-Dem \(> 0{,}05\) in 24 mesi +
uscita da Group of Friends for Peace + BRICS partecipazione passiva & 24
mesi post-insediamento (gen 2027 $\to$ gen 2029) & \(\Delta \geq 0{,}05\)
punti, evento ufficiale, evento osservabile & V-Dem + comunicati
ufficiali & \(0{,}60\) condizionale & strutturale (la
\emph{solitonizzazione} discorsiva è il meccanismo principale del Cap.
7.3) \\
\end{longtable}}
}

\Proved{} \textbf{Sintesi delle 13 predizioni Cap. 7-9.} Probabilità
soggettiva mediana: \(p \approx 0{,}60\). Distribuzione: 2 strutturali
(C2, B4), 2 di meta-stima (B1, U5), 9 locali. Il fallimento di una
predizione locale aggiusterebbe parametri specifici; il fallimento di
una predizione strutturale costringerebbe a ripensare l'apparato
corrispondente.



\section{Tre Predizioni di Sistema Globale (Cap.
10)}\label{tre-predizioni-di-sistema-globale-cap.-10}

{\def\LTcaptype{none} % do not increment counter
{\small\begin{longtable}[]{@{}
  >{\raggedright\arraybackslash}p{(\linewidth - 12\tabcolsep) * \real{0.1429}}
  >{\raggedright\arraybackslash}p{(\linewidth - 12\tabcolsep) * \real{0.1429}}
  >{\raggedright\arraybackslash}p{(\linewidth - 12\tabcolsep) * \real{0.1429}}
  >{\raggedright\arraybackslash}p{(\linewidth - 12\tabcolsep) * \real{0.1429}}
  >{\raggedright\arraybackslash}p{(\linewidth - 12\tabcolsep) * \real{0.1429}}
  >{\raggedright\arraybackslash}p{(\linewidth - 12\tabcolsep) * \real{0.1429}}
  >{\raggedright\arraybackslash}p{(\linewidth - 12\tabcolsep) * \real{0.1429}}@{}}
\toprule\noalign{}
\begin{minipage}[b]{\linewidth}\raggedright
\#
\end{minipage} & \begin{minipage}[b]{\linewidth}\raggedright
Predizione
\end{minipage} & \begin{minipage}[b]{\linewidth}\raggedright
\(t^*\)
\end{minipage} & \begin{minipage}[b]{\linewidth}\raggedright
\(v^*\)
\end{minipage} & \begin{minipage}[b]{\linewidth}\raggedright
Fonte \(\mathcal{D}\)
\end{minipage} & \begin{minipage}[b]{\linewidth}\raggedright
\(p\)
\end{minipage} & \begin{minipage}[b]{\linewidth}\raggedright
Falsificazione
\end{minipage} \\
\midrule\noalign{}
\endhead
\bottomrule\noalign{}
\endlastfoot
\textbf{G1} &
\(P(\text{sistema globale 2031 in regime di decoerenza generalizzata})\)
--- operazionalizzato come \emph{almeno due fra USA, Cina, UE in
\(\mathcal{R}_-\)} & dicembre 2031 & \(\geq 2\) attori in
\(\mathcal{R}_-\) & aggregazione V-Dem, Pew, Brand Finance, IMF, NBS al
2031 & \(0{,}80\) & strutturale (forma di \(T_s^{(2),\text{sys}}\)) \\
\textbf{G2} &
\(P(\text{sistema globale 2031 con almeno un $\Phi$-engine globale operativo})\)
--- operazionalizzato come Brasile \(C_{\text{DEQ}}^2 \geq 1\)
\emph{oppure} nuovo attore equivalente & dicembre 2031 &
\(C_{\text{DEQ}}^2 \geq 1\) per almeno un attore & calcolo dashboard
SPC-DEQ & \([0{,}55, 0{,}70]\), condizionale a B1 & locale (dipende da
B1) \\
\textbf{G3} & Cina: divergenza \(d^{\text{oss}} - d^{\text{spont}}\)
crescente, con \emph{fatica del \(K\) imposto} visibile & finestra di
osservazione 2029-2031 &
\(\Delta(d^{\text{oss}} - d^{\text{spont}}) > 0{,}10\) in 24 mesi &
dashboard SPC-DEQ + V-Dem + NBS & \(0{,}55\) & strutturale (definizione
di \emph{Rigido Sincronizzato} del Cap. 8.0) \\
\end{longtable}}
}

\Hypoth{} \textbf{Lettura.} G1 e G2 sono \emph{complementari}: la teoria
afferma simultaneamente che il sistema globale 2031 sarà in decoerenza
\emph{e} che esiste con probabilità \(\approx 0{,}55\)-\(0{,}70\) un
\(\Phi\)-engine operativo. La conferma di entrambe
(\(\Pr \approx 0{,}50\)) consoliderebbe il quadro M1 del Cap. 10.4. La
conferma di G1 con falsificazione di G2 (\(\Pr \approx 0{,}25\))
realizzerebbe M2 --- multipolarità degenere senza centro coerente.



\section{La Predizione Meta: l'Elezione Brasiliana del 4 Ottobre
2026}\label{la-predizione-meta-lelezione-brasiliana-del-4-ottobre-2026}

\Hypoth{} \textbf{Tesi.} B1 non è una predizione \emph{fra le altre} ---
è la \textbf{predizione che condiziona tutte quelle brasiliane (B2, B3,
B4) e una predizione globale (G2)}. La sua valutazione, prevista per il
4 ottobre 2026 (primo turno) e l'ottobre/novembre 2026 (eventuale
secondo turno), è il \emph{singolo evento più informativo} dell'intero
apparato DEQ$^2$ fra aprile 2026 e dicembre 2031.

\Proved{} \textbf{Protocollo di valutazione di B1.}

{\def\LTcaptype{none} % do not increment counter
{\small\begin{longtable}[]{@{}
  >{\raggedright\arraybackslash}p{(\linewidth - 4\tabcolsep) * \real{0.3333}}
  >{\raggedright\arraybackslash}p{(\linewidth - 4\tabcolsep) * \real{0.3333}}
  >{\raggedright\arraybackslash}p{(\linewidth - 4\tabcolsep) * \real{0.3333}}@{}}
\toprule\noalign{}
\begin{minipage}[b]{\linewidth}\raggedright
Data
\end{minipage} & \begin{minipage}[b]{\linewidth}\raggedright
Azione
\end{minipage} & \begin{minipage}[b]{\linewidth}\raggedright
Conseguenza
\end{minipage} \\
\midrule\noalign{}
\endhead
\bottomrule\noalign{}
\endlastfoot
30 giugno 2026 & Re-stima \(p_{B1}\) su evidenze sondaggi aggregati &
aggiornamento bayesiano della probabilità del Mondo M1/M2 \\
30 settembre 2026 & Re-stima \(p_{B1}\) alla vigilia del primo turno &
aggiornamento finale pre-evento \\
4 ottobre 2026, sera & Risultato primo turno: \(\geq 50\%\) Lula? Se sì
$\to$ B1 vera & falsificazione binaria \\
Ottobre/novembre 2026 (eventuale 2$^\circ$ turno) & Risultato secondo turno:
vince A/B o C? & falsificazione binaria definitiva \\
Entro 30 novembre 2026 & Pubblicazione del \emph{post-mortem} B1 &
aggiornamento di tutta la cascata B2-B3-B4-G2 \\
\end{longtable}}
}

\Hypoth{} \textbf{Conseguenza per la teoria.} Se B1 = \emph{vera}
(vittoria A/B): la cascata B2, B3, G2 conserva la sua probabilità
condizionale. Se B1 = \emph{falsa} (vittoria C): la cascata B2, B3, G2
si riorienta verso B4 e M2. In entrambi i casi, la DEQ$^2$ resta valida ---
ma il \emph{contenuto empirico futuro} del libro cambia
significativamente.



\section{Protocollo di Pubblicazione e
Falsificazione}\label{protocollo-di-pubblicazione-e-falsificazione}

\Hypoth{} \textbf{Pre-registrazione.} Le sedici predizioni di questo
capitolo sono \emph{pre-registrate} alla data di pubblicazione del libro
(target 2026 Q3 --- KDP IT/EN/ES/PT). La pre-registrazione consiste
nella pubblicazione contestuale di:

\begin{enumerate}
\def\labelenumi{\arabic{enumi}.}
\tightlist
\item
  \textbf{Tabella delle predizioni} (sezioni 11.1, 11.2) con tutti i
  campi (1)-(5);
\item
  \textbf{Repository pubblico} (GitHub
  \texttt{bilanciere-quantistico/predizioni-deq2}) contenente:

  \begin{itemize}
  \tightlist
  \item
    file CSV con tutte le predizioni e i loro stati (pendente / vera /
    falsa);
  \item
    script Python di calcolo automatico delle predizioni quantitative
    dai dataset pubblici;
  \item
    log dei \emph{checkpoint} trimestrali con re-stima probabilità
    soggettive;
  \end{itemize}
\item
  \textbf{Protocollo di aggiornamento}: ogni 31 marzo dell'anno
  successivo, pubblicazione di un \emph{update report} con:

  \begin{itemize}
  \tightlist
  \item
    stato di ciascuna predizione;
  \item
    aggiornamento bayesiano delle probabilità;
  \item
    eventuali correzioni alla teoria (locali/parametriche/strutturali).
  \end{itemize}
\end{enumerate}

\Proved{} \textbf{Conseguenza editoriale.} Il libro non è
\emph{autorevole per autorevolezza} ma \emph{autorevole per esposizione
al rischio}. Ogni copia venduta riceve l'URL del repository e ha accesso
al ciclo di update. Questo è coerente con la posizione epistemologica
della Topografia App. A: produrre segnali \((x, y, z)\) alti su
\emph{tutti} gli assi --- incluso \(z\) (complessità dialettica) tramite
l'esposizione esplicita alla falsificazione.



\section{Cosa Accadrebbe se le Predizioni
Fallissero}\label{cosa-accadrebbe-se-le-predizioni-fallissero}

\Hypoth{} \textbf{Mappa di vulnerabilità della teoria.}

{\def\LTcaptype{none} % do not increment counter
{\small\begin{longtable}[]{@{}
  >{\raggedright\arraybackslash}p{(\linewidth - 2\tabcolsep) * \real{0.5000}}
  >{\raggedright\arraybackslash}p{(\linewidth - 2\tabcolsep) * \real{0.5000}}@{}}
\toprule\noalign{}
\begin{minipage}[b]{\linewidth}\raggedright
Predizione
\end{minipage} & \begin{minipage}[b]{\linewidth}\raggedright
Fallimento $\to$ Conseguenza per DEQ$^2$
\end{minipage} \\
\midrule\noalign{}
\endhead
\bottomrule\noalign{}
\endlastfoot
U1 & V-Dem USA migliora rapidamente: aggiusta \(\Phi^{\text{USA}}\) e
\(\zeta_{\text{Sol}}\) effettivo. Locale. \\
U2-U3 & Dollaro tiene: rivedere stima \(\sigma\) esogena globale.
Locale. \\
U4 & Bassa attività veto/SCOTUS: rivedere stima \(\Omega^{\text{USA}}\).
Locale. \\
U5 & Correlazione \(r < 0{,}3\): il modello \(\zeta_{\text{Sol}}(t)\)
del Cap. 6.6 va riformulato. Locale-parametrica. \\
C1 & Deflazione finisce prima: anti-involution funziona. Aggiustamento
\(\Omega^{\text{Cina}}\). Locale. \\
C2 & Nessun bloqueo malgrado condizioni soddisfatte: ipotesi
\emph{Taiwan come valvola} va rivista. \textbf{Strutturale}. \\
C3 & Yuan oltre \(5\%\) entro 2030: la dedollarizzazione converge sullo
yuan, contro la tesi Cap. 8.6. Locale-parametrica. \\
C4 & PPI positivo prima del 2027: anti-involution efficace.
Aggiustamento \(\Omega^{\text{Cina}}\). Locale. \\
B1 & Vittoria di Flávio o destra radicale: B2-B3 cadono, attiva ramo B4.
Cascata controfattuale prevista. \\
B2-B3 & (condizionali a B1 vera) fallimento: rivedere capacità operativa
BRICS+ e \(s_{iii}^{\text{Brasile}}\). Locale. \\
B4 & (condizionale a B1 falsa) fallimento --- Brasile resiste sotto
destra radicale: la \emph{solitonizzazione} del Cap. 7.3 va ripensata
come meccanismo non automatico. \textbf{Strutturale}. \\
G1 & Solo uno o nessun grande attore in \(\mathcal{R}_-\) entro 2031: la
forma stessa di \(T_s^{(2),\text{sys}}\) va ricalibrata.
\textbf{Strutturale}. \\
G2 & Brasile fallisce \(C_{\text{DEQ}}^2 \geq 1\) malgrado A/B: la
definizione di \(\Phi\)-engine è insufficiente. \textbf{Strutturale}. \\
G3 & Cina mantiene \(d^{\text{oss}} \approx d^{\text{spont}}\): la
categoria \emph{Rigido Sincronizzato} non descrive la Cina come
ipotizzato. \textbf{Strutturale}. \\
\end{longtable}}
}

\Proved{} \textbf{Sintesi.} Su 16 predizioni: \textbf{5 strutturali}
(C2, B4, G1, G2, G3), \textbf{9 locali}, \textbf{2 di meta-stima} (B1,
U5). Il fallimento di una predizione strutturale impone
\emph{riformulazione}; il fallimento di tutte le strutturali
simultaneamente \emph{invaliderebbe la teoria nel suo insieme}.
Probabilità di fallimento simultaneo (assumendo indipendenza
condizionale debole): \(\Pr \approx (1-0{,}55)^5 \approx 0{,}018\) ---
basso ma non trascurabile.



\section{Sintesi del Capitolo}\label{sintesi-del-capitolo}

\Proved{} \textbf{In una frase:} la DEQ$^2$ si espone deliberatamente a
sedici predizioni datate, con probabilità soggettive esplicite e fonti
dati pubbliche; cinque di queste sono strutturali e potrebbero
invalidare l'apparato; il libro è dunque un \emph{pre-registered
prediction set} il cui valore epistemico cresce o diminuisce in
proporzione alla performance verificabile della teoria.

{\def\LTcaptype{none} % do not increment counter
{\small\begin{longtable}[]{@{}
  >{\raggedright\arraybackslash}p{(\linewidth - 6\tabcolsep) * \real{0.2500}}
  >{\raggedright\arraybackslash}p{(\linewidth - 6\tabcolsep) * \real{0.2500}}
  >{\raggedright\arraybackslash}p{(\linewidth - 6\tabcolsep) * \real{0.2500}}
  >{\raggedright\arraybackslash}p{(\linewidth - 6\tabcolsep) * \real{0.2500}}@{}}
\toprule\noalign{}
\begin{minipage}[b]{\linewidth}\raggedright
Categoria
\end{minipage} & \begin{minipage}[b]{\linewidth}\raggedright
Numero
\end{minipage} & \begin{minipage}[b]{\linewidth}\raggedright
Probabilità mediana \(p\)
\end{minipage} & \begin{minipage}[b]{\linewidth}\raggedright
Tipo di falsificazione
\end{minipage} \\
\midrule\noalign{}
\endhead
\bottomrule\noalign{}
\endlastfoot
Cap. 7 USA & 5 & \(0{,}60\) & 4 locali, 1 parametrica/meta (U5) \\
Cap. 8 Cina & 4 & \(0{,}60\) & 3 locali, 1 strutturale (C2) \\
Cap. 9 Brasile & 4 & meta + condizionali & 2 condizionali, 1 strutturale
(B4), 1 meta (B1) \\
Cap. 10 globali & 3 & \(0{,}65\) & 2 strutturali (G1, G2 e G3 entrambe),
1 condizionale \\
\textbf{Totale} & \textbf{16} & \(\approx 0{,}60\) & \textbf{5
strutturali, 9 locali/condizionali, 2 meta} \\
\end{longtable}}
}

\Hypoth{} \textbf{Tesi conclusiva.} Le predizioni di questo capitolo
sono il \emph{contenuto empirico futuro} della DEQ$^2$. La Parte V del
libro (Cap. 12-13) costruirà invece il \emph{contenuto normativo}: dato
lo spazio degli scenari probabilizzati, quali sono le strategie
antifragili --- individuali, organizzative, statali --- coerenti con
ciascuno dei quattro mondi possibili al 2031?



\section{Note Epistemiche di Chiusura
Capitolo}\label{note-epistemiche-di-chiusura-capitolo}

{\def\LTcaptype{none} % do not increment counter
{\small\begin{longtable}[]{@{}
  >{\raggedright\arraybackslash}p{(\linewidth - 4\tabcolsep) * \real{0.3333}}
  >{\raggedright\arraybackslash}p{(\linewidth - 4\tabcolsep) * \real{0.3333}}
  >{\raggedright\arraybackslash}p{(\linewidth - 4\tabcolsep) * \real{0.3333}}@{}}
\toprule\noalign{}
\begin{minipage}[b]{\linewidth}\raggedright
\#
\end{minipage} & \begin{minipage}[b]{\linewidth}\raggedright
Affermazione
\end{minipage} & \begin{minipage}[b]{\linewidth}\raggedright
Status
\end{minipage} \\
\midrule\noalign{}
\endhead
\bottomrule\noalign{}
\endlastfoot
1 & Cinque criteri DEQ$^2$ per predizione rischiosa & \Hypoth{}
metodologia \\
2 & Tredici predizioni Cap. 7-9 aggregate con campi (1)-(5) & \Proved{}
aggregazione completa \\
3 & Tre predizioni di sistema globale dal Cap. 10 & \Hypoth{}
derivazione coerente \\
4 & B1 come predizione meta che condiziona la cascata & \Hypoth{}
architettura logica \\
5 & Protocollo di pre-registrazione e repository GitHub & \Hypoth{}
commitment editoriale \\
6 & Mappa di vulnerabilità delle 16 predizioni & \Hypoth{} analisi
strutturale \\
7 & \(\Pr(\text{fallimento simultaneo strutturale}) \approx 0{,}018\) &
\Conj{} stima sotto indipendenza condizionale \\
8 & \(5\) strutturali + \(9\) locali/condizionali + \(2\) meta &
\Proved{} classificazione completa \\
9 & Update report annuale al 31 marzo & \Hypoth{} protocollo
operativo \\
\end{longtable}}
}



\section{Chiusura della Parte IV e Apertura della Parte
V}\label{chiusura-della-parte-iv-e-apertura-della-parte-v}

La \textbf{Parte IV} è ora completa. Cap. 10 ha tradotto l'apparato in
dashboard operativa (sei pannelli + aggregato + cinque regole WE-DEQ$^2$ +
indice \(C_{\text{DEQ}}^2\) + scenari + quattro mondi possibili 2031).
Cap. 11 ha aggregato le 16 predizioni rischiose datate con criteri
popperiani operativi e protocollo di falsificazione pubblica.

La \textbf{Parte V} può ora costruire il contenuto normativo:

\begin{itemize}
\tightlist
\item
  \textbf{Cap. 12} --- \emph{Cosa fare}: strategie individuali (livello
  \(\boldsymbol{p}_{\text{personale}}\)), organizzative (livello
  \(\boldsymbol{p}_{\text{azienda/comunità}}\)), statali (livello
  \(\boldsymbol{p}_{\text{nazionale}}\)), condizionate ai quattro mondi
  possibili M1-M4 del Cap. 10.4. La distinzione \emph{protocollo per
  ramo A/B} vs \emph{protocollo per ramo C} è esplicita.
\item
  \textbf{Cap. 13} --- Conclusioni e ricerca futura: posizione del libro
  nella tradizione (Arrighi, Turchin, Wallerstein, Cervo,
  Marcuse-Adorno), agenda di estensione (paper EN, articolo PT, dataset
  50+ paesi), inviti a falsificazione.
\end{itemize}




%% ═══════════════════════════════════════════════════════════════════════════
%% PARTE V — Protocolli Antifragili
%% ═══════════════════════════════════════════════════════════════════════════
\part{Protocolli Antifragili}

\input{chapters-it/12-strategie-antifragili}
\chapter{Conclusioni e Ricerca
Futura}\label{conclusioni-e-ricerca-futura}

\begin{quote}
\textbf{Argomento portante:} la DEQ$^2$ è una teoria \emph{aperta} ---
esposta alla falsificazione (Cap. 11), strutturata per estensione (paper
EN, articolo PT, dataset multi-paese), inserita esplicitamente in una
tradizione (Arrighi, Turchin, Wallerstein, Cervo, Marcuse-Adorno, Taleb)
di cui rivendica l'eredità senza pretendere di esaurirla. Questo
capitolo conclude il libro come \emph{atto epistemico e politico}:
chiude la trattazione, dichiara cosa resta da fare, invita
esplicitamente al dissenso informato.
\end{quote}



\section{Riepilogo della Tesi
Centrale}\label{riepilogo-della-tesi-centrale}

\Proved{} \textbf{Tesi.} Il destino del sistema globale nel 2031 sarà
deciso dalla distribuzione della \textbf{coerenza di fase \(\Phi\)} tra
quattro attori strutturalmente diversi:

\begin{enumerate}
\def\labelenumi{\arabic{enumi}.}
\tightlist
\item
  \textbf{Solitone (Musk)} --- attore singolo con \(\Phi\) degenere ma
  sorgente esogena \(\zeta_{\text{Sol}}(t)\) di decoerenza per i sistemi
  macro;
\item
  \textbf{Egemone Decoerente (USA)} --- capacità materiale residua ma
  \(\Phi \to 0\) per polarizzazione, iper-estensione, perdita di
  legittimità esterna;
\item
  \textbf{Rigido Sincronizzato (Cina)} --- \(\Phi^{\text{osservato}}\)
  massima per imposizione gerarchica, \(\Phi^{\text{spontaneo}}\) in
  collasso strutturale (Neijuan, demografia, decoupling);
\item
  \textbf{Bilanciere Quantistico (Brasile)} --- unico attore con
  condizioni strutturali per produrre \(\Phi\) esportabile,
  \emph{condizionalmente} all'esito elettorale del 4 ottobre 2026.
\end{enumerate}

\Hypoth{} \textbf{Risultato strutturale principale.} Solo il Bilanciere
produce \emph{\(\Phi\) esportabile vera}. Solitone, Egemone e Rigido
producono rispettivamente \(\zeta\) esogeno, decoerenza distribuita, e
\(K\) imposto --- nessuno dei tre è \(\Phi\). La differenza è formale,
non retorica: emerge dalla struttura matematica di
\(T_s^{(2),\text{sys}}\) derivata nel Cap. 4 e consolidata nel Cap. 5.



\section{I Risultati Strutturali
Acquisiti}\label{i-risultati-strutturali-acquisiti}

\Proved{} \textbf{Risultati formali della Parte II (Cap. 3-5).}

{\def\LTcaptype{none} % do not increment counter
{\small\begin{longtable}[]{@{}
  >{\raggedright\arraybackslash}p{(\linewidth - 4\tabcolsep) * \real{0.3333}}
  >{\raggedright\arraybackslash}p{(\linewidth - 4\tabcolsep) * \real{0.3333}}
  >{\raggedright\arraybackslash}p{(\linewidth - 4\tabcolsep) * \real{0.3333}}@{}}
\toprule\noalign{}
\begin{minipage}[b]{\linewidth}\raggedright
\#
\end{minipage} & \begin{minipage}[b]{\linewidth}\raggedright
Risultato
\end{minipage} & \begin{minipage}[b]{\linewidth}\raggedright
Capitolo
\end{minipage} \\
\midrule\noalign{}
\endhead
\bottomrule\noalign{}
\endlastfoot
R-I & Identità \(\delta^* = 0{,}50\) tra Folk Theorem Discorsivo
(Topografia) e Folk Theorem Geometrico (Geometria) --- invarianza di
scala dimostrata & Cap. 3.1.3 \\
R-II & Operatore di trasposizione
\(\langle X \rangle_\Phi := \Phi \cdot N^{-1}\sum_j X_j\) via Kuramoto &
Cap. 3.1bis \\
R-III & Forma chiusa
\(T_s^{(2),\text{sys}} = s_{iii} z \delta (1+\Phi\langle A\rangle) / [\sigma \varepsilon_{\text{dis}} (1 + \langle\Omega\rangle/\Phi)]\)
& Cap. 4.4 \\
R-IV & Asimmetria \(\Phi\)-moltiplica-\(A\) vs
\(\Phi\)-divide-\(\Omega\) --- firma della seconda legge termodinamica
nei sistemi sociali & Cap. 4.3.2 \\
R-V & Spazio strutturale \(\mathbb{V}^{(4)}\) con sei relazioni di
accoppiamento bivariate R1-R6 & Cap. 5.1 \\
R-VI & Forma esplicita \(\Phi^*(A, \Omega, \varepsilon_{\text{dis}})\)
della superficie critica \(\mathcal{S}\) & Cap. 5.3 \\
R-VII & Sistema dinamico postulato
\(\dot\Phi, \dot A, \dot\Omega, \dot\varepsilon\) con coefficienti
realizzanti R1-R6 & Cap. 5.4 \\
\end{longtable}}
}

\Proved{} \textbf{Risultati applicativi della Parte III (Cap. 6-9).}

{\def\LTcaptype{none} % do not increment counter
{\small\begin{longtable}[]{@{}
  >{\raggedright\arraybackslash}p{(\linewidth - 4\tabcolsep) * \real{0.3333}}
  >{\raggedright\arraybackslash}p{(\linewidth - 4\tabcolsep) * \real{0.3333}}
  >{\raggedright\arraybackslash}p{(\linewidth - 4\tabcolsep) * \real{0.3333}}@{}}
\toprule\noalign{}
\begin{minipage}[b]{\linewidth}\raggedright
\#
\end{minipage} & \begin{minipage}[b]{\linewidth}\raggedright
Risultato
\end{minipage} & \begin{minipage}[b]{\linewidth}\raggedright
Capitolo
\end{minipage} \\
\midrule\noalign{}
\endhead
\bottomrule\noalign{}
\endlastfoot
R-VIII & Profilo discorsivo Musk \((8, 9, 2)\) = equilibrio separante
bayesiano della Topografia App. A & Cap. 6.2 \\
R-IX & Solitone come sorgente \(\zeta(t)\) esogena nel sistema USA,
verificata empiricamente 2025-2026 & Cap. 6.6 + Cap. 7.4 \\
R-X & Solitonizzazione discorsiva degli USA: \((7,6,8) \to (8,9,2)\),
firma su Brand Finance Soft Power 2026 & Cap. 7.3 \\
R-XI & Mappatura Neijuan (Xiang Biao) \(\to \Omega\) DEQ$^2$;
anti-involution campaign cinese come ammissione istituzionale & Cap.
8.2 \\
R-XII & Estado Logístico (Cervo) come \(\Phi^{\text{esportabile}}\)
operativa: sovrapposizione coerente del campo \(\Pi\) & Cap. 9.1 \\
R-XIII & Asimmetria triplice USA/Cina/Brasile: \(\Phi \to 0\) per
percorsi opposti, solo Brasile produce \(\Phi\) vera & Cap. 9.7 \\
\end{longtable}}
}

\Proved{} \textbf{Risultati operativi della Parte IV (Cap. 10-11).}

{\def\LTcaptype{none} % do not increment counter
{\small\begin{longtable}[]{@{}
  >{\raggedright\arraybackslash}p{(\linewidth - 4\tabcolsep) * \real{0.3333}}
  >{\raggedright\arraybackslash}p{(\linewidth - 4\tabcolsep) * \real{0.3333}}
  >{\raggedright\arraybackslash}p{(\linewidth - 4\tabcolsep) * \real{0.3333}}@{}}
\toprule\noalign{}
\begin{minipage}[b]{\linewidth}\raggedright
\#
\end{minipage} & \begin{minipage}[b]{\linewidth}\raggedright
Risultato
\end{minipage} & \begin{minipage}[b]{\linewidth}\raggedright
Capitolo
\end{minipage} \\
\midrule\noalign{}
\endhead
\bottomrule\noalign{}
\endlastfoot
R-XIV & Dashboard SPC-DEQ con sei pannelli + aggregato + cinque regole
WE-DEQ$^2$ + indice \(C_{\text{DEQ}}^2\) & Cap. 10 \\
R-XV & Quattro mondi possibili al 2031 (M1-M4) con probabilità
soggettive esplicite & Cap. 10.4 \\
R-XVI &
\(\Pr(\text{decoerenza globale 2031}) \approx 0{,}80\)-\(0{,}85\),
\(\Pr(\Phi\text{-engine}) \approx 0{,}55\)-\(0{,}70\) & Cap. 10.4 \\
R-XVII & 16 predizioni rischiose datate con campi popperiani (1)-(5) &
Cap. 11 \\
R-XVIII & Classificazione: 5 strutturali, 8 locali, 2 meta.
\(\Pr(\text{fall. simult. strutt.}) \approx 0{,}018\) & Cap. 11.5 \\
\end{longtable}}
}

\Proved{} \textbf{Risultati normativi della Parte V (Cap. 12).}

{\def\LTcaptype{none} % do not increment counter
{\small\begin{longtable}[]{@{}
  >{\raggedright\arraybackslash}p{(\linewidth - 4\tabcolsep) * \real{0.3333}}
  >{\raggedright\arraybackslash}p{(\linewidth - 4\tabcolsep) * \real{0.3333}}
  >{\raggedright\arraybackslash}p{(\linewidth - 4\tabcolsep) * \real{0.3333}}@{}}
\toprule\noalign{}
\begin{minipage}[b]{\linewidth}\raggedright
\#
\end{minipage} & \begin{minipage}[b]{\linewidth}\raggedright
Risultato
\end{minipage} & \begin{minipage}[b]{\linewidth}\raggedright
Capitolo
\end{minipage} \\
\midrule\noalign{}
\endhead
\bottomrule\noalign{}
\endlastfoot
R-XIX & Matrice 3$\times$3 di protocolli operativi (livelli $\times$ orizzonti) con
condizione antifragile esplicita & Cap. 12.1-12.4 \\
R-XX & Strategia \(S_{\mathcal{R}_-}\) per cittadini in sistemi sotto
soglia & Cap. 12.6 \\
R-XXI & Distinzione politica antifragile vs realista/utopica come
superamento concettuale & Cap. 12.7 \\
\end{longtable}}
}

\Hypoth{} \textbf{Sintesi.} Ventuno risultati strutturali acquisiti,
organizzati in quattro parti, con tracciabilità formale a Topografia,
Geometria, MQNU. La DEQ$^2$ è ora un \emph{corpus teorico chiuso} ---
pronto per estensione (paper EN, articolo PT) e per falsificazione (Cap.
11 protocollo).



\section{Posizione del Libro nella
Tradizione}\label{posizione-del-libro-nella-tradizione}

\Hypoth{} \textbf{Mappa esplicita.}

{\def\LTcaptype{none} % do not increment counter
{\small\begin{longtable}[]{@{}
  >{\raggedright\arraybackslash}p{(\linewidth - 6\tabcolsep) * \real{0.2500}}
  >{\raggedright\arraybackslash}p{(\linewidth - 6\tabcolsep) * \real{0.2500}}
  >{\raggedright\arraybackslash}p{(\linewidth - 6\tabcolsep) * \real{0.2500}}
  >{\raggedright\arraybackslash}p{(\linewidth - 6\tabcolsep) * \real{0.2500}}@{}}
\toprule\noalign{}
\begin{minipage}[b]{\linewidth}\raggedright
Tradizione
\end{minipage} & \begin{minipage}[b]{\linewidth}\raggedright
Eredità accolta
\end{minipage} & \begin{minipage}[b]{\linewidth}\raggedright
Mossa DEQ$^2$
\end{minipage} & \begin{minipage}[b]{\linewidth}\raggedright
Distanza dichiarata
\end{minipage} \\
\midrule\noalign{}
\endhead
\bottomrule\noalign{}
\endlastfoot
\textbf{Arrighi} --- cicli egemonici & finanziarizzazione terminale =
segnale di transizione & \(T_s^{(2),\text{sys}}\) formalizza la
transizione & la transizione 2028-2031 è quella di Arrighi
\emph{misurata} \\
\textbf{Turchin} --- structural-demographic & elite overproduction
\(\subset \Omega\); PSI \(\subset \varepsilon_{\text{dis}}\) & aggiunta
di \(\Phi\) come parametro mancante & Turchin cattura il denominatore,
manca il numeratore \\
\textbf{Wallerstein} --- world-system & Brasile non
\emph{semi-periferia} ma \(\Phi\)-engine & riposizionamento brasiliano
nello schema centro-periferia & rifiuto della linearità
centro-periferia \\
\textbf{Cervo} --- Estado Logístico & base operativa del Cap. 9 &
formalizzazione matematica via \(\Pi^{\text{Logístico}}\)
multi-vettoriale & radicamento esplicito \\
\textbf{Taleb} --- antifragilità & \(A\) come
\(\partial^2\Pi/\partial\sigma^2\) & inserimento in apparato sistemico,
asimmetria \(\Phi\)-\(A\)/\(\Omega\) & da euristica a struttura \\
\textbf{Gramsci/Cox} --- egemonia critica & egemonia come operatore di
proiezione (eredità MQNU) & inserimento in DEQ$^2$ come \(s_{iii}\)
proiettata sull'asse \(z\) & continuità filosofica \\
\textbf{Luttwak} --- grande strategia & modello di registro: formalismo
+ compattezza & integrazione: registro tecnico + apparato matematico &
continuità di stile \\
\textbf{Scuola di Francoforte} (Marcuse, Adorno, Habermas) --- via
Topografia & microfondazione comunicativa di \(\Pi\) & \(T_s^{(2)}\)
emerge dal campo \(\Pi\) topografico & la critica francofortese è ora
\emph{misurabile} \\
\textbf{Xiang Biao} --- Neijuan etnografato & \(\Omega\) DEQ$^2$ come
formalizzazione del Neijuan & da osservazione antropologica a variabile
sistemica & formalizzazione \\
\textbf{Kuramoto} --- sincronizzazione fisica & \(\Phi\) parametro
d'ordine & trasposizione di scala micro$\to$macro per accoppiamenti
emittente-ricevente & applicazione metodologica \\
\end{longtable}}
}

\Proved{} \textbf{Conseguenza posizionale.} La DEQ$^2$ \emph{non
sostituisce} nessuna delle tradizioni elencate --- le \emph{integra} in
un apparato in cui ciascuna riceve una collocazione formale. Il rifiuto
del manicheismo (DEQ$^2$ \emph{contro} X) è esplicito: ogni tradizione
cattura una parte vera e parziale dello stesso fenomeno.



\section{Agenda di Estensione}\label{agenda-di-estensione}

\Hypoth{} \textbf{Tre vettori già aperti.}

\subsection{\texorpdfstring{Paper EN --- \emph{From Discursive
Topography to Hegemonic Phase Coherence} (target:
\emph{Complexity})}{Paper EN --- From Discursive Topography to Hegemonic Phase Coherence (target: Complexity)}}\label{paper-en-from-discursive-topography-to-hegemonic-phase-coherence-target-complexity}

Stato attuale: scaffold completo (PAPER-EN-Topography-to-Hegemony.md,
2026-04-25). Lavoro residuo:

\begin{itemize}
\tightlist
\item
  Dataset 50+ paesi $\times$ 20+ anni (1995-2025) per calibrazione
  \(\Sigma^{\text{geo}} \leftarrow \Sigma^{\text{comm}}\) (Appendice B);
\item
  Bootstrap \(n \geq 5000\) per intervalli di confidenza dei
  coefficienti del sistema dinamico Cap. 5.4;
\item
  Repository GitHub \texttt{bilanciere-quantistico/predizioni-deq2}
  operativo (Cap. 11.4);
\item
  Sottomissione \emph{Complexity} o, in alternativa, \emph{Physica A},
  \emph{EPJ B}, \emph{Journal of Conflict Resolution}.
\end{itemize}

\subsection{\texorpdfstring{Articolo PT --- versione brasiliana (target:
\emph{Revista Brasileira de Política Internacional} o \emph{Contexto
Internacional})}{Articolo PT --- versione brasiliana (target: Revista Brasileira de Política Internacional o Contexto Internacional)}}\label{articolo-pt-versione-brasiliana-target-revista-brasileira-de-poluxedtica-internacional-o-contexto-internacional}

Focus: il Brasile come \(\Phi\)-engine condizionale, l'\emph{Estado
Logístico} come politica formalizzabile, le predizioni B1-B4 con
dettaglio sui contesti regionali (Mercosul, BRICS, COP30 follow-up). Da
scrivere in portoghese, registro accademico brasiliano. Lunghezza target
\(20\)-\(25\) pagine. Tempistiche: post-elettorali (novembre 2026) per
integrare l'esito B1.

\subsection{Saggio IT/EN/ES/PT --- distribuzione
KDP}\label{saggio-itenespt-distribuzione-kdp}

Versione attuale: bozza completa di tutti i 13 capitoli + Appendice B.
Lavoro residuo:

\begin{itemize}
\tightlist
\item
  Appendici A (derivazione formale completa), C (protocolli di
  misurazione --- preparata da \S{}5.6 e Cap. 10), D (dialogo testuale con
  tradizioni);
\item
  Revisione editoriale, illustrazioni concettuali per i sei pannelli
  SPC-DEQ;
\item
  Traduzioni EN/ES/PT;
\item
  Upload KDP via pipeline esistente.
\end{itemize}



\section{Tre Vettori Ancora da
Aprire}\label{tre-vettori-ancora-da-aprire}

\Hypoth{} \textbf{R\&D futuro non ancora avviato.}

\begin{enumerate}
\def\labelenumi{\arabic{enumi}.}
\tightlist
\item
  \textbf{Calibrazione real-time della Dashboard SPC-DEQ} come \emph{web
  service pubblico}: aggregare V-Dem, IMF, NBS, Pew, Brand Finance in
  cruscotto live aggiornato trimestralmente. Costo stimato:
  \(40\)-\(80\)k\$ + 1-2 anni di sviluppo. Possibile partnership con
  think tank (Brookings, Chatham House, Stimson, Itamaraty).
\item
  \textbf{Estensione del modello a sotto-attori non-statali}: applicare
  \(\boldsymbol{p}(t) \in \mathbb{V}^{(4)}\) a aziende, città, ONG
  transnazionali, comunità religiose. Necessita ricalibrazione delle sei
  relazioni R1-R6 sulle scale appropriate (la R4 \(A\)-\(\Omega\) è
  probabilmente meno forte a scala aziendale, dove la concorrenza spinge
  \(A\) e \(\Omega\) verso correlazione meno negativa).
\item
  \textbf{Integrazione formale con MQNU}: \(\Phi\) DEQ$^2$ come
  \emph{modulo} di un'ampiezza di probabilità nello spazio degli stati
  di MQNU? La connessione è sospettata ma non formalizzata. Potrebbe
  richiedere un terzo libro.
\end{enumerate}



\section{Inviti Espliciti alla
Falsificazione}\label{inviti-espliciti-alla-falsificazione}

\Hypoth{} \textbf{Sette modi diretti per attaccare la DEQ$^2$.}

{\def\LTcaptype{none} % do not increment counter
{\small\begin{longtable}[]{@{}
  >{\raggedright\arraybackslash}p{(\linewidth - 4\tabcolsep) * \real{0.3333}}
  >{\raggedright\arraybackslash}p{(\linewidth - 4\tabcolsep) * \real{0.3333}}
  >{\raggedright\arraybackslash}p{(\linewidth - 4\tabcolsep) * \real{0.3333}}@{}}
\toprule\noalign{}
\begin{minipage}[b]{\linewidth}\raggedright
\#
\end{minipage} & \begin{minipage}[b]{\linewidth}\raggedright
Critica
\end{minipage} & \begin{minipage}[b]{\linewidth}\raggedright
Risposta DEQ$^2$ (e dove cercare)
\end{minipage} \\
\midrule\noalign{}
\endhead
\bottomrule\noalign{}
\endlastfoot
1 & L'asimmetria \(\Phi\)-moltiplica-\(A\) vs \(\Phi\)-divide-\(\Omega\)
è arbitraria & Cap. 4.3.2: argomento termodinamico; alternative non
lineari da testare nel paper EN \\
2 & \(\Phi\) di Kuramoto non è applicabile a sistemi sociali & Cap. 4.1:
ipotesi di campo medio è esplicita; alternative (correlazione di
Pearson, mutual information, Granger causality) producono risultati
equivalenti per sistemi \(N\) grande \\
3 & Le 16 predizioni Cap. 11 sono troppo numerose per essere
genuinamente rischiose & Cap. 11.5: 5 sono \emph{strutturali}, ciascuna
individualmente capace di invalidare l'apparato \\
4 & La centralità di B1 (elezione brasiliana) è interessata & Cap. 9
\S{}9.0 + Profilo Utente: l'autore ha dichiarato esplicitamente la propria
origine brasiliana e ha mantenuto onestà analitica con tre fragilità
esplicitate \\
5 & Il modello sottostima eventi estremi (cigni neri) & Cap. 12.6:
\(S_{\mathcal{R}_-}\) è progettato esplicitamente per scenari peggiori
del \emph{worst} mediano \\
6 & L'apparato Estado Logístico è specificamente brasiliano e non
generalizzabile & Cap. 9.1: la struttura formale di
\(\Pi^{\text{Logístico}}\) multi-vettoriale è generalizzabile a
\emph{qualunque} attore con \(s_{iii} > s_{ii}\) relativo --- Brasile è
il caso paradigmatico, non l'unico possibile \\
7 & La DEQ$^2$ ignora variabili climatiche/ecologiche & Cap. 4.1 + Cap. 5:
\(\sigma\) esogena include shock climatici; COP30 è incorporata nel Cap.
9.2 come test di \(\Phi^{\text{esportabile}}\). Estensione formale
dell'asse climatico è agenda futura \\
\end{longtable}}
}

\Proved{} \textbf{Tutte le critiche sopra producono \emph{aggiornamento}
della teoria, non sua sospensione}. La DEQ$^2$ è progettata per crescere
via critica.



\section{Chiusura}\label{chiusura}

\Hypoth{} \textbf{Ultimo paragrafo del libro.}

Le predizioni di questo libro saranno verificate o falsificate fra
ottobre 2026 (B1 --- elezione brasiliana) e dicembre 2031 (G1, G2 ---
sistema globale). Nel frattempo il sistema globale opera, decide,
sceglie. La DEQ$^2$ non pretende di prescrivere quelle scelte --- pretende
di \emph{renderle leggibili}, di mostrare il loro spazio degli effetti,
di rendere trasparente il costo dell'inazione e l'asimmetria del
rischio. Se il suo apparato si dimostrerà robusto, sarà perché è stato
attaccato; se verrà superato, sarà da una teoria che fa meglio le stesse
cose. In entrambi i casi avrà fatto il suo lavoro.

Il sistema globale entra nel triennio 2028-2031 con coerenza di fase in
declino quasi ovunque, capacità materiale concentrata in attori
strutturalmente fragili, e una sola candidatura non degenere a
\(\Phi\)-engine globale. Quella candidatura sarà confermata o smentita
dal Brasile il 4 ottobre 2026. Da quel giorno in poi, ogni cittadino,
ogni organizzazione, ogni Stato avrà un dato in più per scegliere fra
strategie passive e strategie antifragili. Questo libro è scritto perché
quella scelta sia \emph{informata}, non perché sia
\emph{predeterminata}.



\section{Note Epistemiche di Chiusura del
Libro}\label{note-epistemiche-di-chiusura-del-libro}

{\def\LTcaptype{none} % do not increment counter
{\small\begin{longtable}[]{@{}
  >{\raggedright\arraybackslash}p{(\linewidth - 4\tabcolsep) * \real{0.3333}}
  >{\raggedright\arraybackslash}p{(\linewidth - 4\tabcolsep) * \real{0.3333}}
  >{\raggedright\arraybackslash}p{(\linewidth - 4\tabcolsep) * \real{0.3333}}@{}}
\toprule\noalign{}
\begin{minipage}[b]{\linewidth}\raggedright
\#
\end{minipage} & \begin{minipage}[b]{\linewidth}\raggedright
Affermazione
\end{minipage} & \begin{minipage}[b]{\linewidth}\raggedright
Status
\end{minipage} \\
\midrule\noalign{}
\endhead
\bottomrule\noalign{}
\endlastfoot
1 & Tesi centrale ricapitolata fedelmente & \Proved{} \\
2 & 21 risultati strutturali tracciati ai capitoli di origine &
\Proved{} aggregazione completa \\
3 & Posizione esplicita verso 10 tradizioni & \Hypoth{} mappa
concettuale \\
4 & Tre vettori di estensione già aperti (paper EN, articolo PT, saggio
multi-lingua) & \Proved{} pianificati \\
5 & Tre vettori R\&D futuri (web service, sub-attori non-statali,
integrazione MQNU) & \Conj{} prospettici \\
6 & Sette inviti espliciti alla falsificazione con risposte DEQ$^2$ &
\Hypoth{} esposizione metodologica \\
7 & Chiusura politica: la DEQ$^2$ come rendering visibile, non come
prescrizione & \Hypoth{} posizione esplicita \\
\end{longtable}}
}



\section{Ringraziamenti}\label{ringraziamenti}

A \textbf{Clayton Teixeira}, collaboratore principale del corpus
matematico (in particolare \emph{La Costante di Marcelino}, base
epistemica dell'apparato formale). Alla \textbf{tradizione critica
francofortese} che ha reso pensabile la formalizzazione della
comunicazione come oggetto matematico. Ai \textbf{lettori che vorranno
falsificare} anziché accogliere --- il loro lavoro è la condizione di
vita della teoria.




%% ═══════════════════════════════════════════════════════════════════════════
%% APPENDICI
%% ═══════════════════════════════════════════════════════════════════════════
\appendix

\chapter{\texorpdfstring{Derivazione Formale Completa di
\(T_s^{(2),\text{sys}}\)}{Derivazione Formale Completa di T\_s\^{}\{(2),\textbackslash text\{sys\}\}}}\label{derivazione-formale-completa-di-t_s2textsys}

\begin{quote}
\textbf{Scopo dell'appendice.} Esplicitare ogni passaggio matematico
dalla legge micro-comunicativa \(\Pi_j(t+1)\) della Topografia (\S{}5.7)
fino alla forma sistemica \(T_s^{(2),\text{sys}}\) del Cap. 4,
dichiarando \emph{tutte} le ipotesi operative e identificando i punti in
cui la derivazione ammette alternative. L'appendice è progettata per
essere autonomamente verificabile da un lettore con background in fisica
statistica (Kuramoto), teoria dei giochi ripetuti (Folk Theorem) e
Bayesian inference.
\end{quote}



\section{Setup Formale Esteso}\label{setup-formale-esteso}

\Hypoth{} \textbf{Definizione (Sistema sociale formalizzato).} Un
\emph{sistema sociale} \(\mathcal{X}\) è una tripla
\((\mathcal{N}, \{\Pi_j\}, \mathcal{K})\) dove:

\begin{itemize}
\tightlist
\item
  \(\mathcal{N} = \{1, 2, \ldots, N\}\) è l'insieme degli
  \emph{accoppiamenti elementari} --- ciascun \(j\) rappresenta un
  accoppiamento emittente-ricevente nel senso della Topografia (\S{}5.1);
\item
  \(\Pi_j : \mathbb{R}_{\geq 0} \to \mathbb{R}\) è la legge di payoff
  istantaneo dell'accoppiamento \(j\), della forma topografica
  \(\Pi_j(t) = M_{e,j}(t) \cdot Q_{0,j}(t) \cdot (1 - \delta_{R,j} z_{R,j})\);
\item
  \(\mathcal{K} : \mathcal{N} \times \mathcal{N} \to \mathbb{R}_{\geq 0}\)
  è la matrice di accoppiamento di Kuramoto, riducibile per ipotesi di
  campo medio a \(\mathcal{K}_{jk} \equiv K/N\) (\(K \geq 0\) scalare).
\end{itemize}

\Proved{} \textbf{Conseguenza notazionale.} Ciascun
\(j \in \mathcal{N}\) ha: - una \textbf{fase}
\(\theta_j(t) \in [0, 2\pi)\) --- posizione nel ciclo di
legittimazione/decoerenza; - una \textbf{frequenza naturale}
\(\omega_j\) con distribuzione \(g(\omega)\) unimodale simmetrica
intorno a \(\omega_0 = 0\) (riconducibile a questo caso per cambio di
sistema di riferimento); - un \textbf{payoff} \(\Pi_j(t)\) come sopra.



\section{Dinamica di Kuramoto e il Parametro
d'Ordine}\label{dinamica-di-kuramoto-e-il-parametro-dordine}

\Proved{} \textbf{Equazione di Kuramoto (1984).} L'evoluzione delle fasi
è governata da:

\[
\frac{d\theta_j}{dt} = \omega_j + \frac{K}{N} \sum_{k=1}^{N} \sin(\theta_k - \theta_j), \qquad j = 1, \ldots, N
\]

Definiamo il \textbf{parametro d'ordine complesso}:

\[
\Phi(t) e^{i\psi(t)} := \frac{1}{N} \sum_{j=1}^{N} e^{i\theta_j(t)}, \quad \Phi \in [0, 1], \quad \psi \in [0, 2\pi)
\]

\Proved{} \textbf{Riduzione di campo medio.} L'equazione di Kuramoto si
riscrive equivalentemente:

\[
\frac{d\theta_j}{dt} = \omega_j + K \Phi(t) \sin(\psi(t) - \theta_j)
\]

(passaggio standard, vedi Strogatz 2000 per la derivazione).

\Proved{} \textbf{Soglia critica \(K_c\).} Per \(g(\omega)\) unimodale
simmetrica, esiste

\[
K_c = \frac{2}{\pi g(0)}
\]

tale che: - \(K < K_c \Rightarrow \Phi \to 0\) asintoticamente (regime
di decoerenza completa); -
\(K > K_c \Rightarrow \Phi \to \Phi_\infty(K) > 0\) asintoticamente
(regime di sincronizzazione parziale); - \(\Phi_\infty(K) \to 1\) per
\(K \to \infty\) (sincronizzazione piena).

\Hypoth{} \textbf{Ipotesi A1 (Kuramoto applicabile).} La distribuzione
\(g(\omega)\) delle frequenze naturali degli accoppiamenti
emittente-ricevente in un sistema sociale è unimodale simmetrica.
\emph{Verifica empirica}: per cicli mediatici, elettorali, di mercato,
la distribuzione è approssimativamente lognormale; per simmetria
approssimata accettabile per \(N\) grande (limite centrale).



\section{\texorpdfstring{Operatore di Trasposizione
\(\langle \cdot \rangle_\Phi\)}{Operatore di Trasposizione \textbackslash langle \textbackslash cdot \textbackslash rangle\_\textbackslash Phi}}\label{operatore-di-trasposizione-langle-cdot-rangle_phi}

\Hypoth{} \textbf{Definizione (operatore di trasposizione).} Per una
grandezza locale \(X_j\) definita su \(\mathcal{N}\), la \emph{quantità
sistemica derivata} è:

\[
\langle X \rangle_\Phi := \Phi \cdot \frac{1}{N} \sum_{j=1}^{N} X_j = \Phi \cdot \langle X \rangle
\]

dove \(\langle X \rangle := N^{-1} \sum_j X_j\) è la media aritmetica
locale.

\Proved{} \textbf{Proprietà dell'operatore.}

{\def\LTcaptype{none} % do not increment counter
{\small\begin{longtable}[]{@{}
  >{\raggedright\arraybackslash}p{(\linewidth - 2\tabcolsep) * \real{0.5000}}
  >{\raggedright\arraybackslash}p{(\linewidth - 2\tabcolsep) * \real{0.5000}}@{}}
\toprule\noalign{}
\begin{minipage}[b]{\linewidth}\raggedright
Proprietà
\end{minipage} & \begin{minipage}[b]{\linewidth}\raggedright
Verifica
\end{minipage} \\
\midrule\noalign{}
\endhead
\bottomrule\noalign{}
\endlastfoot
Linearità:
\(\langle aX + bY \rangle_\Phi = a\langle X\rangle_\Phi + b\langle Y\rangle_\Phi\)
& Diretta da linearità di \(\sum\) \\
\(\langle\text{cost}\rangle_\Phi = \Phi \cdot \text{cost}\) & Diretta \\
\(\langle X \rangle_\Phi \to \langle X \rangle\) per \(\Phi \to 1\) &
Diretta \\
\(\langle X \rangle_\Phi \to 0\) per \(\Phi \to 0\) & Diretta \\
\(\partial \langle X \rangle_\Phi / \partial \Phi = \langle X \rangle\)
& Diretta \\
\end{longtable}}
}

\Hypoth{} \textbf{Interpretazione.} \(\langle X \rangle_\Phi\) è la
quantità \(X\) \emph{effettivamente esprimibile dal sistema};
\(\langle X \rangle\) è la quantità \(X\) \emph{potenzialmente
disponibile}. La differenza
\(\langle X \rangle - \langle X \rangle_\Phi = (1-\Phi)\langle X \rangle\)
misura quanto \(X\) è \emph{perduto per decoerenza}.



\section{\texorpdfstring{Trasposizione di \(A\) e
\(\Omega\)}{Trasposizione di A e \textbackslash Omega}}\label{trasposizione-di-a-e-omega}

\Hypoth{} \textbf{Postulato A2 (trasposizione di \(A\) per linearità di
Taylor).} L'antifragilità sistemica \(A_{\text{sys}}\) è la più semplice
forma consistente con le condizioni al contorno
\(A_{\text{sys}}(\Phi=1) = \langle A\rangle\) e
\(A_{\text{sys}}(\Phi=0) = 0\):

\[
A_{\text{sys}}(\Phi) = \Phi \cdot \langle A \rangle = \langle A \rangle_\Phi
\]

\Hypoth{} \textbf{Postulato A3 (trasposizione di \(\Omega\) per
amplificazione inversa).} L'involuzione sistemica
\(\Omega_{\text{sys}}\) è la più semplice forma consistente con
\(\Omega_{\text{sys}}(\Phi=1) = \langle\Omega\rangle\) e
\(\Omega_{\text{sys}}(\Phi \to 0) \to \infty\):

\[
\Omega_{\text{sys}}(\Phi) = \frac{\langle \Omega \rangle}{\Phi}
\]

\Proved{} \textbf{Asimmetria formale.} \(A_{\text{sys}}\) è
\emph{omogenea di grado +1 in \(\Phi\)}; \(\Omega_{\text{sys}}\) è
\emph{omogenea di grado --1 in \(\Phi\)}. Questa asimmetria di grado è la
firma matematica della dissimmetria termodinamica fra valore (necessita
trasporto) e spreco (si crea ovunque non sia inibito).

\Conj{} \textbf{Alternative non banali da testare empiricamente.}

{\def\LTcaptype{none} % do not increment counter
{\small\begin{longtable}[]{@{}
  >{\raggedright\arraybackslash}p{(\linewidth - 4\tabcolsep) * \real{0.3333}}
  >{\raggedright\arraybackslash}p{(\linewidth - 4\tabcolsep) * \real{0.3333}}
  >{\raggedright\arraybackslash}p{(\linewidth - 4\tabcolsep) * \real{0.3333}}@{}}
\toprule\noalign{}
\begin{minipage}[b]{\linewidth}\raggedright
Forma
\end{minipage} & \begin{minipage}[b]{\linewidth}\raggedright
Verifica condizioni al contorno
\end{minipage} & \begin{minipage}[b]{\linewidth}\raggedright
Razionale alternativo
\end{minipage} \\
\midrule\noalign{}
\endhead
\bottomrule\noalign{}
\endlastfoot
\(A_{\text{sys}} = (1 - e^{-\alpha\Phi})\langle A\rangle\) & \Proved{}
se \(\alpha\) tale che \(1-e^{-\alpha} = 1\) (cioè
\(\alpha \to \infty\), recupera lineare) & saturazione per \(\Phi\)
alto \\
\(A_{\text{sys}} = \Phi^\beta \langle A\rangle\) & \Proved{} per
qualunque \(\beta > 0\) & \(\beta < 1\): rendimenti decrescenti;
\(\beta > 1\): soglia di attivazione \\
\(\Omega_{\text{sys}} = \langle\Omega\rangle / \Phi^\gamma\) & \Proved{}
per qualunque \(\gamma > 0\) & \(\gamma > 1\): amplificazione esplosiva
della decoerenza \\
\end{longtable}}
}

\Proved{} \textbf{Posizione DEQ$^2$. } La forma lineare
\(A_{\text{sys}} = \Phi\langle A\rangle\) e razionale
\(\Omega_{\text{sys}} = \langle\Omega\rangle/\Phi\) sono la \emph{prima
approssimazione consistente}; la loro validità sarà testata sul dataset
50+ paesi del paper EN.



\section{\texorpdfstring{Costruzione di
\(T_s^{(2),\text{sys}}\)}{Costruzione di T\_s\^{}\{(2),\textbackslash text\{sys\}\}}}\label{costruzione-di-t_s2textsys}

\Proved{} \textbf{Punto di partenza (Cap. 3.4).} La forma locale
dell'equazione DEQ$^2$ è:

\[
T_s^{(2),\text{loc}} = \frac{s_{iii} \cdot z \cdot \delta \cdot (1+A)}{\sigma \cdot \varepsilon_{\text{dis}} \cdot (1+\Omega)}
\]

dove \(A, \Omega\) sono i valori locali del singolo accoppiamento o
sotto-attore.

\Hypoth{} \textbf{Passaggio alla scala sistemica.} Trattiamo
\(s_{iii}, z, \delta, \sigma, \varepsilon_{\text{dis}}\) come
\emph{quantità sistemiche aggregate} (non sottoposte a trasposizione
\(\langle \cdot \rangle_\Phi\) poiché sono parametri di campo esogeno o
quasi-esogeno alla scala temporale di interesse). Le sole quantità che
richiedono trasposizione sono \(A\) e \(\Omega\):

\[
A \to A_{\text{sys}} = \Phi \langle A\rangle, \qquad \Omega \to \Omega_{\text{sys}} = \langle\Omega\rangle/\Phi
\]

\Proved{} \textbf{Sostituzione diretta.}

\[
\boxed{T_s^{(2),\text{sys}} = \frac{s_{iii} \cdot z \cdot \delta \cdot (1 + \Phi \langle A\rangle)}{\sigma \cdot \varepsilon_{\text{dis}} \cdot \left(1 + \dfrac{\langle\Omega\rangle}{\Phi}\right)}}
\]

\Proved{} \textbf{Riscrittura in forma fisicamente trasparente.}
Moltiplichiamo numeratore e denominatore per \(\Phi\):

\[
T_s^{(2),\text{sys}} = \frac{\Phi \cdot s_{iii} \cdot z \cdot \delta \cdot (1 + \Phi \langle A\rangle)}{\Phi \cdot \sigma \cdot \varepsilon_{\text{dis}} + \sigma \cdot \varepsilon_{\text{dis}} \cdot \langle\Omega\rangle}
\]

Da questa forma è immediato leggere:

\begin{itemize}
\tightlist
\item
  numeratore \(\propto \Phi\): la stabilità sistemica scala
  \emph{linearmente} con la coerenza, plus correzione antifragile
  \(\Phi^2 \langle A\rangle\);
\item
  denominatore: la decoerenza ha due componenti additive --- una
  \emph{proporzionale a \(\Phi\)} (volatilità modulata dalla coerenza) e
  una \emph{indipendente da \(\Phi\)}
  (\(\sigma \varepsilon_{\text{dis}} \langle\Omega\rangle\),
  l'involuzione strutturale).
\end{itemize}



\section{Verifiche di Consistenza nei
Limiti}\label{verifiche-di-consistenza-nei-limiti}

\Proved{} \textbf{Tabella di verifica completa.}

{\def\LTcaptype{none} % do not increment counter
{\small\begin{longtable}[]{@{}
  >{\raggedright\arraybackslash}p{(\linewidth - 4\tabcolsep) * \real{0.3333}}
  >{\raggedright\arraybackslash}p{(\linewidth - 4\tabcolsep) * \real{0.3333}}
  >{\raggedright\arraybackslash}p{(\linewidth - 4\tabcolsep) * \real{0.3333}}@{}}
\toprule\noalign{}
\begin{minipage}[b]{\linewidth}\raggedright
Limite
\end{minipage} & \begin{minipage}[b]{\linewidth}\raggedright
\(T_s^{(2),\text{sys}}\)
\end{minipage} & \begin{minipage}[b]{\linewidth}\raggedright
Interpretazione
\end{minipage} \\
\midrule\noalign{}
\endhead
\bottomrule\noalign{}
\endlastfoot
\(\Phi \to 1\) &
\(\dfrac{s_{iii}z\delta(1+\langle A\rangle)}{\sigma\varepsilon_{\text{dis}}(1+\langle\Omega\rangle)}\)
& Recupero della forma locale Cap. 3.4 con \(A=\langle A\rangle\),
\(\Omega=\langle\Omega\rangle\) \Proved{} \\
\(\Phi \to 0\) &
\(\dfrac{s_{iii}z\delta \cdot 1}{\sigma\varepsilon_{\text{dis}} \cdot \infty} \to 0\)
& Decoerenza completa \(\Rightarrow\) collasso sistemico \Proved{} \\
\(\langle A\rangle \to \infty\) &
\(\propto \Phi^2 \cdot \langle A\rangle / [\sigma\varepsilon_{\text{dis}}(\Phi+\langle\Omega\rangle)] \to \infty\)
& Antifragilità infinita \(\Rightarrow\) stabilità infinita (caso
teorico) \Proved{} \\
\(\langle\Omega\rangle \to \infty\) & \(\to 0\) & Involuzione totale
\(\Rightarrow\) collasso \Proved{} \\
\(\sigma \to 0\) & \(\to \infty\) & Volatilità nulla \(\Rightarrow\)
stabilità infinita (caso teorico) \Proved{} \\
\(\sigma \to \infty\) & \(\to 0\) a meno che
\(\langle A\rangle \to \infty\) commensurabile & Volatilità infinita
\(\Rightarrow\) collasso \Proved{} \\
\(\varepsilon_{\text{dis}} \to 0\) & \(\to \infty\) & Disimpegno nullo
\(\Rightarrow\) stabilità infinita \Proved{} \\
\(\delta \to 0\) & \(\to 0\) & Miopia totale \(\Rightarrow\) collasso
\Proved{} \\
\(K \to K_c^-\) & \(\Phi \to 0 \Rightarrow T_s^{(2),\text{sys}} \to 0\)
& Sotto soglia di accoppiamento \(\Rightarrow\) collasso sistemico
inevitabile \Proved{} \\
\end{longtable}}
}

\Proved{} \textbf{Risultato.} Tutti i limiti hanno interpretazione
coerente con la fenomenologia attesa. Nessun limite produce
comportamento patologico (oscillazioni, salti, cambi di segno).



\section{\texorpdfstring{Forma della Superficie Critica
\(\mathcal{S}\)}{Forma della Superficie Critica \textbackslash mathcal\{S\}}}\label{forma-della-superficie-critica-mathcals}

\Hypoth{} \textbf{Definizione operativa.} La \emph{soglia critica di
coerenza} \(\Phi^*\) è il valore di \(\Phi\) per cui
\(T_s^{(2),\text{sys}} = 1\):

\[
s_{iii} z \delta (1 + \Phi^* \langle A\rangle) = \sigma \varepsilon_{\text{dis}} \left(1 + \frac{\langle\Omega\rangle}{\Phi^*}\right)
\]

\Proved{} \textbf{Risoluzione.} Riscrivendo:

\[
s_{iii} z \delta + s_{iii} z \delta \Phi^* \langle A\rangle = \sigma \varepsilon_{\text{dis}} + \sigma \varepsilon_{\text{dis}} \langle\Omega\rangle / \Phi^*
\]

Moltiplicando per \(\Phi^*\):

\[
s_{iii} z \delta \Phi^* + s_{iii} z \delta \langle A\rangle (\Phi^*)^2 = \sigma\varepsilon_{\text{dis}}\Phi^* + \sigma\varepsilon_{\text{dis}}\langle\Omega\rangle
\]

Riordinando:

\[
s_{iii} z \delta \langle A\rangle (\Phi^*)^2 + (s_{iii}z\delta - \sigma\varepsilon_{\text{dis}}) \Phi^* - \sigma\varepsilon_{\text{dis}}\langle\Omega\rangle = 0
\]

\Proved{} \textbf{Soluzione (formula quadratica).} Posto
\(a := s_{iii}z\delta\langle A\rangle\),
\(b := s_{iii}z\delta - \sigma\varepsilon_{\text{dis}}\),
\(c := -\sigma\varepsilon_{\text{dis}}\langle\Omega\rangle\):

\[
\boxed{\Phi^* = \frac{-b + \sqrt{b^2 - 4ac}}{2a}}
\]

(scelta della radice positiva perché \(\Phi^* \in [0,1]\)).

\Hypoth{} \textbf{Discussione.} La forma quadratica raffinata è più
precisa della formulazione lineare del Cap. 4.5.1. Quest'ultima era
ottenuta per linearizzazione attorno a \(\Phi^* \approx 1\) ---
appropriata per sistemi prossimi alla coerenza piena, ma sotto-stimata
per sistemi in regime di forte decoerenza (USA in Cap. 7).

\Proved{} \textbf{Limite consistente.} Per
\(\langle\Omega\rangle \to 0\) (\(c \to 0\)), \(\Phi^* \to 0\) se
\(b > 0\) --- recupero della prima formulazione lineare per sistemi
senza involuzione.



\section{Stabilità Lineare dell'Equilibrio
Sistemico}\label{stabilituxe0-lineare-dellequilibrio-sistemico}

\Hypoth{} \textbf{Tesi.} Il punto
\(\boldsymbol{p}^*(t) := (\Phi^*, A^*, \Omega^*, \varepsilon_{\text{dis}}^*)\)
corrispondente a \(T_s^{(2),\text{sys}} = 1\) è di \textbf{equilibrio
instabile} rispetto a perturbazioni in \(\Phi\).

\Proved{} \textbf{Prova (analisi locale).} Calcoliamo
\(\partial T_s^{(2),\text{sys}} / \partial \Phi\) a
\(\boldsymbol{p}^*\):

\[
\frac{\partial T_s^{(2),\text{sys}}}{\partial \Phi} = \frac{s_{iii}z\delta\langle A\rangle}{\sigma\varepsilon_{\text{dis}}(1+\langle\Omega\rangle/\Phi)} + \frac{s_{iii}z\delta(1+\Phi\langle A\rangle) \cdot \langle\Omega\rangle/\Phi^2}{\sigma\varepsilon_{\text{dis}}(1+\langle\Omega\rangle/\Phi)^2}
\]

Entrambi i termini sono positivi a \(\boldsymbol{p}^*\). Dunque
\(T_s^{(2),\text{sys}}\) è \emph{crescente} in \(\Phi\) in un intorno di
\(\Phi^*\).

\Hypoth{} \textbf{Conseguenza dinamica.} Perturbazioni
\(\Phi^* + \epsilon\) con \(\epsilon < 0\) portano
\(T_s^{(2),\text{sys}} < 1\) --- il sistema si auto-rinforza nella
decoerenza tramite R3 (\(\Phi\)-\(\varepsilon_{\text{dis}}\) negativo,
Cap. 5.1) e R2 (\(\Phi\)-\(\Omega\) negativo). Perturbazioni
\(\Phi^* + \epsilon\) con \(\epsilon > 0\) portano
\(T_s^{(2),\text{sys}} > 1\) --- il sistema si auto-rinforza nella
coerenza per gli stessi meccanismi.

\Proved{} \textbf{Risultato strutturale.} \(\boldsymbol{p}^*\) è un
\emph{separatore di destini}: piccole perturbazioni in \(\Phi\)
producono, attraverso le sei relazioni R1-R6, traiettorie che divergono
in \(\mathbb{V}^{(4)}\). Questo è la base formale del concetto di
\emph{transizione di fase} del Cap. 3.5.3 e della \emph{biforcazione}
del Cap. 9.5.



\section{Note sulle Ipotesi
Implicite}\label{note-sulle-ipotesi-implicite}

\Conj{} \textbf{Ipotesi A4 (separazione di scala temporale).} I
parametri esogeni \((s_{iii}, z, \delta, \sigma)\) variano su scale
temporali molto più lunghe (decennali per \(s_{iii}, z, \delta\); ciclo
geopolitico per \(\sigma\)) rispetto alle variabili di
\(\mathbb{V}^{(4)}\) (\(\Phi\) stagionale, \(A\) post-stress, \(\Omega\)
trimestrale, \(\varepsilon_{\text{dis}}\) semestrale). Questa
separazione di scala giustifica trattare i parametri come quasi-costanti
per finestre di osservazione di 5-10 anni.

\Conj{} \textbf{Ipotesi A5 (additività delle modulazioni).} Le forme
\((1+A)\) al numeratore e \((1+\Omega)\) al denominatore della
\(T_s^{(2),\text{loc}}\) sono additive --- alternative moltiplicative o
esponenziali producono comportamenti qualitativamente analoghi nei
limiti, ma differenze quantitative significative nel regime intermedio.
La scelta additiva è giustificata da minimalità (Occam) e da consistenza
con le forme di Bandura per il disimpegno morale.

\Conj{} \textbf{Ipotesi A6 (omogeneità della popolazione
\(\mathcal{N}\)).} Il campo medio Kuramoto assume omogeneità degli
accoppiamenti. Sistemi con forte clustering (sub-popolazioni con
coerenza interna alta ma fra-cluster bassa, es. polarizzazione USA)
richiedono il modello generalizzato di Kuramoto multi-cluster (Strogatz
2003). Estensione possibile, rinviata al paper EN.

\Proved{} \textbf{Conseguenza metodologica.} Le tre ipotesi A4-A6 sono
\emph{esplicitate}; ciascuna ammette estensione nota in letteratura; la
loro violazione produce \emph{correzioni} alla forma della
\(T_s^{(2),\text{sys}}\) ma non ne invalida la struttura asintotica.



\section{Sintesi dell'Appendice}\label{sintesi-dellappendice}

\Proved{} \textbf{Risultato centrale.} La derivazione formale completa
di \(T_s^{(2),\text{sys}}\) dal campo dei payoff topografici
\(\{\Pi_j(t)\}\) richiede sei ipotesi esplicite (A1-A6, di cui due come
postulati di trasposizione A2-A3), produce una forma chiusa con sei
verifiche di consistenza nei limiti e ammette una superficie critica
\(\mathcal{S}\) con forma quadratica esplicita per
\(\Phi^*(A, \Omega, \varepsilon_{\text{dis}}; s_{iii}, z, \delta, \sigma)\).
L'equilibrio \(\boldsymbol{p}^* \in \mathcal{S}\) è instabile rispetto a
\(\Phi\), giustificando formalmente la nozione di \emph{biforcazione
strategica} del Cap. 9 e dei \emph{quattro mondi possibili} del Cap. 10.

\Hypoth{} \textbf{Posizione epistemica.} L'apparato è \emph{minimale}
(sei ipotesi), \emph{consistente} (sei limiti verificati), \emph{aperto}
(alternative non lineari elencate per test empirico) e
\emph{tracciabile} (ogni passaggio è linkato al capitolo corrispondente
del libro).



\emph{Appendice tecnica. Pre-requisiti: fisica statistica (Strogatz
2000, Kuramoto 1984), teoria dei giochi ripetuti (Folk Theorem,
Fudenberg-Tirole 1991), Bayesian inference (Gelman et al., }Bayesian
Data Analysis\emph{, 3a ed.).}

\input{chapters-it/AppB-calibrazione-sigma}
\chapter{Protocolli di Misurazione dei Parametri
DEQ$^2$}\label{protocolli-di-misurazione-dei-parametri-dequxb2}

\begin{quote}
\textbf{Scopo dell'appendice.} Trasformare le sette variabili
dell'apparato \((s_{iii}, z, \delta, \sigma\) esogene;
\(\Phi, A, \Omega, \varepsilon_{\text{dis}}\) in \(\mathbb{V}^{(4)}\))
in \emph{quantità misurabili in modo riproducibile da terzi
indipendenti}, su base trimestrale o annuale, a partire da fonti dati
pubbliche. L'appendice è la \emph{base tecnica} per il repository GitHub
\texttt{bilanciere-quantistico/predizioni-deq2} del Cap. 11.4 e per la
dashboard SPC-DEQ del Cap. 10.
\end{quote}



\section{Principi del Protocollo}\label{principi-del-protocollo}

\Hypoth{} \textbf{Cinque principi operativi.}

\begin{enumerate}
\def\labelenumi{\arabic{enumi}.}
\tightlist
\item
  \textbf{Replicabilità}: ogni misurazione deve essere ricalcolabile da
  terzi a partire da fonti pubbliche;
\item
  \textbf{Trasparenza dei proxy}: dichiarazione esplicita di proxy
  primario, secondari e fallback;
\item
  \textbf{Gestione del missing data}: regole esplicite di imputazione e
  di esclusione;
\item
  \textbf{Aggiornamento periodico}: ogni variabile ha frequenza di
  rilevazione fissata, con scadenze pubbliche;
\item
  \textbf{Intervallo di confidenza}: ogni stima è accompagnata da CI
  95\% via bootstrap.
\end{enumerate}

\Proved{} \textbf{Conseguenza.} I protocolli sotto consentono il calcolo
del punto
\(\boldsymbol{p}(t) = (\Phi, A, \Omega, \varepsilon_{\text{dis}})\) per
qualunque Stato membro ONU per cui esistano i dataset elencati. La
dashboard SPC-DEQ non richiede dati proprietari.



\section{\texorpdfstring{Protocollo per \(\Phi\) (Coerenza di
Fase)}{Protocollo per \textbackslash Phi (Coerenza di Fase)}}\label{protocollo-per-phi-coerenza-di-fase}

\Hypoth{} \textbf{Definizione operativa.} \(\Phi\) è il parametro
d'ordine di Kuramoto del campo \(\{\theta_j\}\) degli accoppiamenti
emittente-ricevente. Operativizzato come \emph{indice composito
multi-proxy}.

{\def\LTcaptype{none} % do not increment counter
{\small\begin{longtable}[]{@{}
  >{\raggedright\arraybackslash}p{(\linewidth - 8\tabcolsep) * \real{0.2000}}
  >{\raggedright\arraybackslash}p{(\linewidth - 8\tabcolsep) * \real{0.2000}}
  >{\raggedright\arraybackslash}p{(\linewidth - 8\tabcolsep) * \real{0.2000}}
  >{\raggedright\arraybackslash}p{(\linewidth - 8\tabcolsep) * \real{0.2000}}
  >{\raggedright\arraybackslash}p{(\linewidth - 8\tabcolsep) * \real{0.2000}}@{}}
\toprule\noalign{}
\begin{minipage}[b]{\linewidth}\raggedright
Componente
\end{minipage} & \begin{minipage}[b]{\linewidth}\raggedright
Proxy primario
\end{minipage} & \begin{minipage}[b]{\linewidth}\raggedright
Fonte
\end{minipage} & \begin{minipage}[b]{\linewidth}\raggedright
Frequenza
\end{minipage} & \begin{minipage}[b]{\linewidth}\raggedright
Peso
\end{minipage} \\
\midrule\noalign{}
\endhead
\bottomrule\noalign{}
\endlastfoot
\(\Phi^{\text{istituzionale}}\) & V-Dem Liberal Democracy Index &
\url{https://v-dem.net} & annuale (aprile) & 0.40 \\
\(\Phi^{\text{discorsivo}}\) & Pew Research \emph{common ground across
issues} (sei aree) & \url{https://pewresearch.org} & annuale & 0.25 \\
\(\Phi^{\text{fiduciario}}\) & Pew \emph{trust in federal government} &
\url{https://pewresearch.org} & annuale & 0.20 \\
\(\Phi^{\text{mediatico}}\) & Reuters Institute \emph{news trust index}
& \url{https://reutersinstitute.politics.ox.ac.uk} & annuale & 0.15 \\
\end{longtable}}
}

\Proved{} \textbf{Calcolo composito.} Ciascun sotto-indice è
normalizzato in \([0, 1]\) e combinato per media pesata:

\[
\Phi(t) = \sum_k w_k \cdot \tilde{\phi}_k(t)
\]

con \(\tilde{\phi}_k\) valore normalizzato e \(w_k\) peso da tabella.
Per Stati senza dataset Pew/Reuters, fallback a
\(\Phi^{\text{istituzionale}}\) con riduzione del peso totale
(rinormalizzazione).

\Hypoth{} \textbf{Pseudocodice (Python).}

\begin{verbatim}
def phi(country, year):
    vdem = fetch_vdem_lib_dem(country, year)
    pew_cg = fetch_pew_common_ground(country, year)
    pew_tr = fetch_pew_trust(country, year)
    rt_news = fetch_reuters_news_trust(country, year)
    weights = [0.40, 0.25, 0.20, 0.15]
    values = [vdem, pew_cg, pew_tr, rt_news]
    available = [v for v in values if v is not None]
    aw = [w for v, w in zip(values, weights) if v is not None]
    return sum(v*w for v, w in zip(available, aw)) / sum(aw)
\end{verbatim}

\Conj{} \textbf{Caveat.} Lo sbilanciamento Pew/Reuters verso paesi
anglofoni è una limitazione nota. Per Stati non-OECD, sostituire con
World Values Survey + Edelman Trust Barometer +
Afrobarometer/Latinobarómetro/Asiabarometer regionali.



\section{\texorpdfstring{Protocollo per \(A\)
(Antifragilità)}{Protocollo per A (Antifragilità)}}\label{protocollo-per-a-antifragilituxe0}

\Hypoth{} \textbf{Definizione operativa.}
\(A = \partial^2\Pi/\partial\sigma^2\) misurata via \emph{stress-test
retrospettivo} su tre crisi recenti: 2008-2009 (finanziaria), 2020
(pandemica), 2022-2023 (energetica/inflazionistica). Per ciascuna crisi:

\begin{itemize}
\tightlist
\item
  \(\Pi^{\text{baseline}}\): produttività totale media nei 3 anni
  pre-crisi;
\item
  \(\Pi^{\text{stress}}\): produttività media nei 3 anni post-crisi;
\item
  \(\sigma^{\text{shock}}\): deviazione standard di GDP, equity index,
  currency nei 12 mesi della crisi acuta.
\end{itemize}

\Hypoth{} \textbf{Stima.} Per ciascuna crisi \(i\):

\[
A_i = \frac{\Pi^{\text{stress}}_i - \Pi^{\text{baseline}}_i}{(\sigma^{\text{shock}}_i)^2}
\]

\(A^{\text{aggregato}} = \text{mediana}(\{A_1, A_2, A_3\})\) (più
robusto della media a outlier).

{\def\LTcaptype{none} % do not increment counter
{\small\begin{longtable}[]{@{}
  >{\raggedright\arraybackslash}p{(\linewidth - 2\tabcolsep) * \real{0.5000}}
  >{\raggedright\arraybackslash}p{(\linewidth - 2\tabcolsep) * \real{0.5000}}@{}}
\toprule\noalign{}
\begin{minipage}[b]{\linewidth}\raggedright
Fonte primaria
\end{minipage} & \begin{minipage}[b]{\linewidth}\raggedright
Per
\end{minipage} \\
\midrule\noalign{}
\endhead
\bottomrule\noalign{}
\endlastfoot
World Bank WDI (Total Factor Productivity) & \(\Pi\) \\
IMF International Financial Statistics & \(\sigma\) (volatilità
currency, equity) \\
OECD Productivity Database & TFP cross-country \\
\end{longtable}}
}

\Conj{} \textbf{Caveat.} \(A\) è la variabile più difficile in real-time
perché richiede \emph{eventi di stress} per manifestarsi. Per stime in
periodi senza crisi recenti, usare \emph{stress-test simulati} (Monte
Carlo controfattuali su shock di \(\sigma\) standardizzato) --- vedi
protocollo \S{}C.7.

\Hypoth{} \textbf{Pseudocodice.}

\begin{verbatim}
def A(country, crises=[(2008, 2009), (2020, 2020), (2022, 2023)]):
    A_values = []
    for c in crises:
        pre = mean_TFP(country, c[0]-3, c[0]-1)
        post = mean_TFP(country, c[1]+1, c[1]+3)
        shock = std_GDP(country, c[0], c[1])
        A_values.append((post - pre) / shock**2)
    return median(A_values)
\end{verbatim}



\section{\texorpdfstring{Protocollo per \(\Omega\)
(Involuzione)}{Protocollo per \textbackslash Omega (Involuzione)}}\label{protocollo-per-omega-involuzione}

\Hypoth{} \textbf{Definizione operativa.}
\(\Omega = \Delta E_{\text{dissipata}}/\Delta E_{\text{utile}} - 1\).
Operativizzato come \emph{complemento della Total Factor Productivity in
declino persistente}, con correzione per settori specifici dominanti.

{\def\LTcaptype{none} % do not increment counter
{\small\begin{longtable}[]{@{}
  >{\raggedright\arraybackslash}p{(\linewidth - 6\tabcolsep) * \real{0.2500}}
  >{\raggedright\arraybackslash}p{(\linewidth - 6\tabcolsep) * \real{0.2500}}
  >{\raggedright\arraybackslash}p{(\linewidth - 6\tabcolsep) * \real{0.2500}}
  >{\raggedright\arraybackslash}p{(\linewidth - 6\tabcolsep) * \real{0.2500}}@{}}
\toprule\noalign{}
\begin{minipage}[b]{\linewidth}\raggedright
Componente
\end{minipage} & \begin{minipage}[b]{\linewidth}\raggedright
Proxy
\end{minipage} & \begin{minipage}[b]{\linewidth}\raggedright
Fonte
\end{minipage} & \begin{minipage}[b]{\linewidth}\raggedright
Frequenza
\end{minipage} \\
\midrule\noalign{}
\endhead
\bottomrule\noalign{}
\endlastfoot
\(\Omega^{\text{TFP}}\) & \(-\Delta(\ln \text{TFP})/\Delta t\) in
finestra mobile 8 anni & OECD Productivity Database & annuale \\
\(\Omega^{\text{settore-critico}}\) & per Cina: contrazione property +
PPI; per USA: budget federale primario; per Brasile: deficit primario &
NBS (CN), BEA (US), IBGE (BR) & trimestrale \\
\(\Omega^{\text{governance}}\) & costo medio di rinegoziazione
legislativa (proxy: numero di vetoes, cloture votes, judicial overrides)
& sistemi legislativi nazionali & annuale \\
\end{longtable}}
}

\Proved{} \textbf{Combinazione.}
\(\Omega = \max(\Omega^{\text{TFP}}, \Omega^{\text{settore}}, \Omega^{\text{governance}})\)
--- l'involuzione \emph{si sente} dove è più alta, non dove è media.

\Hypoth{} \textbf{Esempio applicativo (Cina 2025).}
\(\Omega^{\text{TFP}} \approx 0{,}3\) (declino lento),
\(\Omega^{\text{settore}} \approx 1{,}2\) (property + deflazione PPI 10
trim.), \(\Omega^{\text{governance}} \approx 0{,}1\) (nessun vincolo
legislativo). \(\Omega^{\text{Cina}}_{2025} = 1{,}2\). Coerente con
stima Cap. 8 e Cap. 5.2.



\section{\texorpdfstring{Protocollo per \(\varepsilon_{\text{dis}}\)
(Disimpegno
Morale)}{Protocollo per \textbackslash varepsilon\_\{\textbackslash text\{dis\}\} (Disimpegno Morale)}}\label{protocollo-per-varepsilon_textdis-disimpegno-morale}

\Hypoth{} \textbf{Definizione operativa.}
\(\varepsilon_{\text{dis}} \in [0,1]\) è il coefficiente di disimpegno
morale collettivo (Bandura/Milgram). Indice composito.

{\def\LTcaptype{none} % do not increment counter
{\small\begin{longtable}[]{@{}
  >{\raggedright\arraybackslash}p{(\linewidth - 8\tabcolsep) * \real{0.2000}}
  >{\raggedright\arraybackslash}p{(\linewidth - 8\tabcolsep) * \real{0.2000}}
  >{\raggedright\arraybackslash}p{(\linewidth - 8\tabcolsep) * \real{0.2000}}
  >{\raggedright\arraybackslash}p{(\linewidth - 8\tabcolsep) * \real{0.2000}}
  >{\raggedright\arraybackslash}p{(\linewidth - 8\tabcolsep) * \real{0.2000}}@{}}
\toprule\noalign{}
\begin{minipage}[b]{\linewidth}\raggedright
Componente
\end{minipage} & \begin{minipage}[b]{\linewidth}\raggedright
Proxy
\end{minipage} & \begin{minipage}[b]{\linewidth}\raggedright
Fonte
\end{minipage} & \begin{minipage}[b]{\linewidth}\raggedright
Frequenza
\end{minipage} & \begin{minipage}[b]{\linewidth}\raggedright
Peso
\end{minipage} \\
\midrule\noalign{}
\endhead
\bottomrule\noalign{}
\endlastfoot
\(\varepsilon^{\text{partisan}}\) & Pew ``altro partito come minaccia
per la nazione'' & \url{https://pewresearch.org} & annuale & 0.30 \\
\(\varepsilon^{\text{turnout}}\) &
\(1 - \text{turnout elettorale normalizzato}\) & IDEA International &
per ciclo elettorale & 0.20 \\
\(\varepsilon^{\text{civic}}\) & \(1 - \text{volontariato pro capite}\)
& Charities Aid Foundation World Giving Index & annuale & 0.20 \\
\(\varepsilon^{\text{media-trust}}\) & \(1 - \text{trust mediatico}\) &
Edelman Trust Barometer & annuale & 0.15 \\
\(\varepsilon^{\text{intergroup}}\) & indice di dissociazione fra gruppi
(proxy: segregazione residenziale, omogamia educativa) & OECD Family
Database & quinquennale & 0.15 \\
\end{longtable}}
}

\Hypoth{} \textbf{Calcolo.} Identico a \(\Phi\) (media pesata,
normalizzazione, gestione missing data).



\section{\texorpdfstring{Protocolli per Variabili Esogene
\((s_{iii}, z, \delta, \sigma)\)}{Protocolli per Variabili Esogene (s\_\{iii\}, z, \textbackslash delta, \textbackslash sigma)}}\label{protocolli-per-variabili-esogene-s_iii-z-delta-sigma}

{\def\LTcaptype{none} % do not increment counter
{\small\begin{longtable}[]{@{}
  >{\raggedright\arraybackslash}p{(\linewidth - 6\tabcolsep) * \real{0.2500}}
  >{\raggedright\arraybackslash}p{(\linewidth - 6\tabcolsep) * \real{0.2500}}
  >{\raggedright\arraybackslash}p{(\linewidth - 6\tabcolsep) * \real{0.2500}}
  >{\raggedright\arraybackslash}p{(\linewidth - 6\tabcolsep) * \real{0.2500}}@{}}
\toprule\noalign{}
\begin{minipage}[b]{\linewidth}\raggedright
Variabile
\end{minipage} & \begin{minipage}[b]{\linewidth}\raggedright
Proxy primario
\end{minipage} & \begin{minipage}[b]{\linewidth}\raggedright
Fonte
\end{minipage} & \begin{minipage}[b]{\linewidth}\raggedright
Frequenza
\end{minipage} \\
\midrule\noalign{}
\endhead
\bottomrule\noalign{}
\endlastfoot
\(s_{iii}\) (legittimità etico-discorsiva) & Brand Finance Global Soft
Power Index &
\url{https://brandfinance.com/insights/global-soft-power-index} &
annuale \\
\(z\) (proiezione asse dialettico) & media ponderata Topografia \S{}5.1
(asse \(z\) del campo \(Q_0\)) & dataset corpus discorso pubblico &
annuale (proxy: indice complessità sintattica + diversità lessicale di
dichiarazioni esecutive) \\
\(\delta\) (orizzonte temporale) & OECD long-term planning indicator +
tasso di sconto sociale (Stern review) & OECD + Stern + AEI &
quinquennale \\
\(\sigma\) (volatilità ambientale) & std GDP + std CDS sovrani + std
petrolio + std eventi climatici estremi & IMF + Bloomberg + EM-DAT &
trimestrale \\
\end{longtable}}
}

\Conj{} \textbf{Caveat per \(z\).} La misurazione del \emph{contenuto
dialettico} del discorso pubblico richiede analisi NLP avanzata ---
pipeline non standardizzata. Proposta: usare il rapporto fra
\emph{propositions argumentate} e \emph{propositions assertive} in
corpus di dichiarazioni governative ufficiali (modello GPT-class come
strumento di codifica).



\section{\texorpdfstring{Stress-Test Simulato (per \(A\) in periodi
senza
crisi)}{Stress-Test Simulato (per A in periodi senza crisi)}}\label{stress-test-simulato-per-a-in-periodi-senza-crisi}

\Hypoth{} \textbf{Procedura Monte Carlo.} Quando la finestra delle tre
crisi recenti non è disponibile (paese giovane, dataset incompleto),
usare:

\begin{verbatim}
def A_simulated(country, year, n_simulations=5000):
    baseline = current_TFP(country, year)
    sigma_hist = historical_volatility(country, year)
    A_estimates = []
    for _ in range(n_simulations):
        shock = random_shock(distribution='lognormal', mean=0, sd=sigma_hist)
        # Modello dinamico: applicare shock, simulare risposta basata su:
        # - struttura federale (federalismo aumenta A)
        # - diversificazione produttiva (HHI inverso)
        # - opzionalità giuridica (numero di forme societarie)
        response = simulate_response(country, shock)
        A_estimates.append((response - baseline) / shock**2)
    return median(A_estimates), CI_95(A_estimates)
\end{verbatim}

\Proved{} \textbf{Validazione.} Calibrare il modello
\(simulate\_response\) su paesi con dataset crisi completo (USA, UE,
Giappone) e verificare che
\(A^{\text{simulato}} \approx A^{\text{empirico}}\) per quei paesi. Se
OK, applicare ad altri Stati.



\section{Tabella di Riepilogo
Operativo}\label{tabella-di-riepilogo-operativo}

{\def\LTcaptype{none} % do not increment counter
{\small\begin{longtable}[]{@{}
  >{\raggedright\arraybackslash}p{(\linewidth - 8\tabcolsep) * \real{0.2000}}
  >{\raggedright\arraybackslash}p{(\linewidth - 8\tabcolsep) * \real{0.2000}}
  >{\raggedright\arraybackslash}p{(\linewidth - 8\tabcolsep) * \real{0.2000}}
  >{\raggedright\arraybackslash}p{(\linewidth - 8\tabcolsep) * \real{0.2000}}
  >{\raggedright\arraybackslash}p{(\linewidth - 8\tabcolsep) * \real{0.2000}}@{}}
\toprule\noalign{}
\begin{minipage}[b]{\linewidth}\raggedright
Variabile
\end{minipage} & \begin{minipage}[b]{\linewidth}\raggedright
Frequenza calcolo
\end{minipage} & \begin{minipage}[b]{\linewidth}\raggedright
Tempo di latenza dati
\end{minipage} & \begin{minipage}[b]{\linewidth}\raggedright
CI tipica
\end{minipage} & \begin{minipage}[b]{\linewidth}\raggedright
Osservabilità del trend
\end{minipage} \\
\midrule\noalign{}
\endhead
\bottomrule\noalign{}
\endlastfoot
\(\Phi\) & annuale & 4-6 mesi (V-Dem) & \(\pm 0{,}03\) & media (segnali
a 1-2 anni) \\
\(A\) & post-stress + simulato annualmente & 3-5 anni (per stress reali)
& \(\pm 0{,}1\) & bassa (richiede crisi) \\
\(\Omega\) & trimestrale & 2-3 mesi & \(\pm 0{,}05\) & alta (segnali a
1-2 trimestri) \\
\(\varepsilon_{\text{dis}}\) & semestrale & 6-12 mesi & \(\pm 0{,}04\) &
media \\
\(s_{iii}\) & annuale & 3 mesi (Brand Finance) & \(\pm 1\) punto &
media \\
\(z\) & annuale & 6-12 mesi (NLP custom) & \(\pm 0{,}1\) & bassa \\
\(\delta\) & quinquennale & 1-3 anni & \(\pm 0{,}05\) & bassissima
(variabile slow) \\
\(\sigma\) & trimestrale & 1 mese & \(\pm 0{,}02\) & altissima \\
\end{longtable}}
}

\Hypoth{} \textbf{Conseguenza per la dashboard.} Il pannello aggregato
\(P_0 = d(\boldsymbol{p}, \mathcal{S})\) del Cap. 10 può essere
aggiornato \emph{trimestralmente} sulla base di
\(\Omega, \sigma, s_{iii}\) con interpolazione di
\(\Phi, \varepsilon_{\text{dis}}\) fra gli aggiornamenti annuali. La
risoluzione temporale effettiva del cruscotto è quindi \emph{quattro
volte/anno}.



\section{Limiti Espliciti del
Protocollo}\label{limiti-espliciti-del-protocollo}

\Conj{} \textbf{Cinque limiti riconosciuti.}

\begin{enumerate}
\def\labelenumi{\arabic{enumi}.}
\tightlist
\item
  \textbf{Bias verso paesi OECD}: la maggior parte dei proxy ha
  copertura completa solo per OECD + grandi emergenti. Per Stati piccoli
  o paesi a basso reddito: copertura ridotta, fallback a sotto-indici.
\item
  \textbf{Aggregazione vs sotto-attori}: la dashboard misura \emph{Stati
  nazionali}. Per sotto-attori (regioni, città, aziende) i protocolli
  vanno re-calibrati (vedi Cap. 13.5 vettore 2 di R\&D).
\item
  \textbf{Lag dati}: la dashboard rappresenta \emph{passato recente},
  non presente; per uso predittivo, applicare il sistema dinamico Cap.
  5.4 con \(\boldsymbol{p}(t-\Delta t)\) come stato iniziale.
\item
  \textbf{Sensibilità ai pesi}: i pesi nei sotto-indici sono scelte
  parametriche; il paper EN testerà la robustezza via sensitivity
  analysis.
\item
  \textbf{NLP per \(z\)}: la pipeline non è standardizzata; la
  riproducibilità di \(z\) è inferiore alle altre variabili.
\end{enumerate}

\Proved{} \textbf{Mitigazione.} Tutti e cinque i limiti sono
\emph{trasparenti} --- il repository GitHub riporta esplicitamente quali
paesi/anni hanno copertura completa vs parziale, con flag di
affidabilità.



\section{Sintesi dell'Appendice}\label{sintesi-dellappendice}

\Proved{} \textbf{Risultato.} Le sette variabili dell'apparato DEQ$^2$ sono
operativizzabili tramite indici compositi multi-proxy basati su dataset
pubblici (V-Dem, Pew, Reuters, OECD, IMF, Brand Finance, Edelman).
Frequenza di aggiornamento: trimestrale per \(\Omega\) e \(\sigma\),
semestrale per \(\varepsilon_{\text{dis}}\), annuale per
\(\Phi, s_{iii}, z, A\), quinquennale per \(\delta\). Latenza tipica
dato-stima: 1-12 mesi a seconda della variabile. Tutto il protocollo è
riproducibile in pseudocodice Python e implementabile come web service.

\Hypoth{} \textbf{Posizione strategica.} Senza un protocollo
riproducibile, le predizioni del Cap. 11 sarebbero tecnicamente non
falsificabili (le valutazioni dipenderebbero dalla scelta dei
misuratori). Con il protocollo C.1-C.9, le predizioni \emph{diventano}
falsificabili da terzi indipendenti, soddisfacendo il criterio (3) del
Cap. 11.0.



\emph{Appendice tecnica. Implementazione di riferimento: repository
\texttt{bilanciere-quantistico/predizioni-deq2} (in setup,
post-pubblicazione libro).}

\chapter{Dialogo Testuale con le
Tradizioni}\label{dialogo-testuale-con-le-tradizioni}

\begin{quote}
\textbf{Scopo dell'appendice.} Espandere la mappa concettuale del Cap.
13.3 in dieci dialoghi brevi e diretti con le tradizioni che la DEQ$^2$
eredita, modifica o supera. Ciascun dialogo è strutturato come
\emph{obiezione attesa} dalla tradizione + \emph{risposta DEQ$^2$} +
\emph{concessione esplicita} di ciò che la DEQ$^2$ \emph{non} fa rispetto a
quella tradizione. La forma è quella del dialogo immaginario: non
pretende di parlare per gli autori citati, pretende di \emph{farsi
parlare} dalle loro categorie.
\end{quote}



\section{\texorpdfstring{Con Giovanni Arrighi (\emph{Il lungo XX
secolo},
1994)}{Con Giovanni Arrighi (Il lungo XX secolo, 1994)}}\label{con-giovanni-arrighi-il-lungo-xx-secolo-1994}

\textbf{Arrighi:} ``Lei ha individuato una transizione 2028-2031 senza
riconoscere che è la \emph{fase autunnale} del ciclo finanziario USA ---
già descritta nel mio ciclo egemonico
genovese-olandese-britannico-americano. Non sta scoprendo nulla, sta
misurando ciò che io ho narrato.''

\textbf{DEQ$^2$:} ``La concessione è piena: Lei ha narrato la \emph{forma}
della transizione; io provo a \emph{misurarne il bordo}. La
\(T_s^{(2),\text{sys}} \to 0\) del Cap. 7 sull'USA è la traduzione
formale del Suo \emph{signal crisis} (anni '70-'80) seguito da
\emph{terminal crisis} (anni 2020-2030). Quello che aggiungo: la
transizione \emph{non genera automaticamente} un nuovo egemone --- può
generare invece un'asimmetria multi-attore in cui \(\Phi\) esportabile
risiede in un Bilanciere (Brasile) anziché in un nuovo egemone (Cina).
Lei pensava in termini di \emph{successione}; io penso in termini di
\emph{biforcazione}.''

\textbf{Concessione.} Senza Arrighi, il Cap. 7 e il Cap. 13.7 sarebbero
illeggibili. La DEQ$^2$ \emph{eredita} l'orizzonte temporale e l'ipotesi di
transizione; \emph{aggiunge} la quantificazione e la possibilità di un
esito non-egemonico.



\section{\texorpdfstring{Con Peter Turchin (\emph{Secular Cycles}, 2009;
\emph{End Times},
2023)}{Con Peter Turchin (Secular Cycles, 2009; End Times, 2023)}}\label{con-peter-turchin-secular-cycles-2009-end-times-2023}

\textbf{Turchin:} ``La Sua \(\Omega\) è la mia \emph{elite
overproduction}; il Suo \(\varepsilon_{\text{dis}}\) è la mia
\emph{Popular Immiseration}. Ha riformulato in notazione greca quello
che ho già modellato cliodinamicamente.''

\textbf{DEQ$^2$:} ``Ha ragione su due variabili su quattro. Ma manca
\(\Phi\). Lei modella il \emph{denominatore} della stabilità sistemica
(involuzione + immiserazione \(\Rightarrow\) instabilità) ma non il
\emph{numeratore} (coerenza di fase \(\times\) antifragilità
\(\Rightarrow\) stabilità positiva). Senza \(\Phi\), la cliodinamica
predice \emph{quando} un sistema crolla, ma non \emph{cosa lo riempie}.
La differenza è cruciale per il 2031: due sistemi possono avere la
stessa \(\Omega + \varepsilon_{\text{dis}}\) e finire in M1 (decoerenza
con \(\Phi\)-engine) o in M2 (decoerenza degenere). La cliodinamica è
agnostica fra M1 e M2; la DEQ$^2$ no.''

\textbf{Concessione.} Le serie temporali storiche di Turchin sono
essenziali per calibrare \(A\) via stress-test retrospettivo (App. C.3).
La DEQ$^2$ rinuncia a fare previsioni a tre secoli (orizzonte cliodinamico)
--- si limita a un \emph{ciclo} (2026-2031).



\section{\texorpdfstring{Con Immanuel Wallerstein (\emph{The Modern
World-System},
1974-2011)}{Con Immanuel Wallerstein (The Modern World-System, 1974-2011)}}\label{con-immanuel-wallerstein-the-modern-world-system-1974-2011}

\textbf{Wallerstein:} ``Il Brasile è semi-periferia ascendente del
sistema-mondo capitalista. Lei lo presenta come \(\Phi\)-engine --- cioè
come un attore con prerogative \emph{centrali} di produzione di ordine.
È un'illusione: la semi-periferia \emph{non} genera coerenza globale, è
generata da essa.''

\textbf{DEQ$^2$:} ``Concedo che lo schema centro-periferia ha potere
descrittivo notevole per il periodo 1500-2000. Ma la sua linearità
(centro \(\to\) semi-periferia \(\to\) periferia) è un \emph{vincolo}
che la DEQ$^2$ rifiuta di importare. La \(\Phi\) non è correlata alla
posizione nel sistema-mondo: dipende dalla struttura interna del sistema
sociale di un paese (Cap. 4.5, Cap. 5.2). Il Brasile può essere
\emph{semi-periferico per output} e \emph{centrale per \(\Phi\)
esportabile} --- è esattamente la posizione strutturale del Cap. 9. Lei
vincola; io disvincolo.''

\textbf{Concessione.} Lo schema centro-periferia resta utile per
descrivere il \emph{flusso materiale} (capitale, manodopera, risorse).
Ma il \emph{flusso di coerenza} segue regole diverse --- l'errore
strutturale del wallersteinismo è confondere i due flussi.



\section{\texorpdfstring{Con Amado Cervo (\emph{Inserção internacional},
2008)}{Con Amado Cervo (Inserção internacional, 2008)}}\label{con-amado-cervo-inseruxe7uxe3o-internacional-2008}

\textbf{Cervo:} ``L'Estado Logístico è una categoria che ho costruito
per descrivere il Brasile post-2003. Lei lo importa, ma la matematica
DEQ$^2$ rischia di trasformare una \emph{categoria politica vivente} in un
\emph{parametro statico}. La politica estera brasiliana non è un
coefficiente.''

\textbf{DEQ$^2$:} ``Concessione totale. La
\(\Pi^{\text{Logístico}} = \sum_k M_e^{(k)} Q_0^{(k)} (1 - \delta_R^{(k)} z_R^{(k)})\)
del Cap. 9.1 \emph{non sostituisce} la Sua categoria --- la
\emph{traduce} in linguaggio operativo per la dashboard SPC-DEQ. La
traduzione è utile \emph{operativamente} (permette monitoraggio in tempo
reale) ma è \emph{parziale concettualmente} (perde la dimensione storica
accumulativa che Lei evidenzia). Per questo l'Estado Logístico nella
DEQ$^2$ è \emph{condizionale} alla continuità politica --- esattamente la
fragilità del Cap. 9.5 + Cap. 12.5.''

\textbf{Concessione.} La DEQ$^2$ rinuncia a teorizzare il \emph{processo
storico} di formazione dell'Estado Logístico (1990-2003). Si limita a
operativizzarne i parametri \emph{date le condizioni del 2025-2031}.



\section{\texorpdfstring{Con Nassim Nicholas Taleb (\emph{Antifragile},
2012)}{Con Nassim Nicholas Taleb (Antifragile, 2012)}}\label{con-nassim-nicholas-taleb-antifragile-2012}

\textbf{Taleb:} ``Lei prende la mia antifragilità individuale e la
incolla a un sistema egemonico globale. La mia teoria è
\emph{anti-sistemica}: l'antifragilità \emph{vive nei dettagli e nelle
code}, non nelle medie aggregate. Il Suo \(\langle A\rangle\) è
esattamente l'errore che combatto.''

\textbf{DEQ$^2$:} ``Concessione parziale. Sì, \(\langle A\rangle\) è una
media --- ma è modulata da \(\Phi\), e l'asimmetria
\(\Phi\)-moltiplica-\(A\) vs \(\Phi\)-divide-\(\Omega\) è esattamente la
firma matematica della Sua intuizione: il sistema \emph{amplifica} la
fragilità più di quanto amplifichi l'antifragilità. La DEQ$^2$ formalizza
la convessità che Lei ha portato come euristica. Inoltre il Cap. 12.6
(\(S_{\mathcal{R}_-}\)) è interamente \emph{taleb-iano}: opzionalità,
ridondanza, disaccoppiamento dalla traiettoria sistemica. La DEQ$^2$
accoglie la Sua critica trasformandola in protocollo.''

\textbf{Concessione.} La DEQ$^2$ \emph{non} è una teoria della
\emph{previsione} nel senso strong (rifiutato da Taleb); è una teoria
dello \emph{spazio degli scenari} con probabilità soggettive esplicite.
Sotto questo profilo è coerente con \emph{Black Swan} e \emph{Skin in
the Game}.



\section{\texorpdfstring{Con Antonio Gramsci (\emph{Quaderni del
carcere}, 1929-1935) e Robert Cox (\emph{Production, Power, and World
Order},
1987)}{Con Antonio Gramsci (Quaderni del carcere, 1929-1935) e Robert Cox (Production, Power, and World Order, 1987)}}\label{con-antonio-gramsci-quaderni-del-carcere-1929-1935-e-robert-cox-production-power-and-world-order-1987}

\textbf{Gramsci/Cox:} ``L'egemonia non è \(T_s^{(2),\text{sys}}\). È
rapporto fra forze materiali, istituzionali e culturali in cui il
consenso e la coercizione si combinano \emph{organicamente}. Lei riduce
l'organico al numerico.''

\textbf{DEQ$^2$:} ``Concedo che \(T_s^{(2),\text{sys}}\) è una
\emph{misura} dell'egemonia, non la \emph{spiegazione} dell'egemonia. La
spiegazione resta gramsciana: il consenso (\(s_{iii}\) DEQ$^2$) e la
coercizione (\(\sigma \cdot \varepsilon_{\text{dis}}\)) si combinano in
proporzioni che variano per blocco storico. Quello che la DEQ$^2$ aggiunge
è la \emph{condizione di stabilità}: per qualunque blocco storico,
\(T_s^{(2),\text{sys}} > 1\) è condizione necessaria di durabilità. La
DEQ$^2$ è dunque \emph{interna} alla teoria gramsciana dell'egemonia, non
alternativa.''

\textbf{Concessione.} La nozione di \emph{blocco storico} gramsciano
contiene una temporalità \emph{qualitativa} (cesura, conservazione,
rivoluzione passiva, restaurazione) che la DEQ$^2$ non cattura. La
transizione di fase del Cap. 3.5.3 corrisponde solo alla \emph{cesura}
gramsciana; le altre forme richiedono estensione.



\section{\texorpdfstring{Con Edward Luttwak (\emph{Strategy: The Logic
of War and Peace},
1987-2001)}{Con Edward Luttwak (Strategy: The Logic of War and Peace, 1987-2001)}}\label{con-edward-luttwak-strategy-the-logic-of-war-and-peace-1987-2001}

\textbf{Luttwak:} ``La grande strategia è \emph{paradossale}: ciò che
funziona si esaurisce, ciò che non funziona viene ripreso. La Sua
matematica è \emph{lineare} nel tempo. Si perde il paradosso.''

\textbf{DEQ$^2$:} ``Concessione. Le forme \(T_s^{(2),\text{sys}}\) sono
monotoniche (Cap. 3.5.1); la dinamica del Cap. 5.4 ammette però
retroazioni non lineari (R6 \(\Omega \to \varepsilon_{\text{dis}}\), R3
\(\Phi \to \varepsilon_{\text{dis}}\) negativa) che producono
comportamento \emph{paradossale} nell'evoluzione di
\(\boldsymbol{p}(t)\): la stessa azione può essere antifragile in M1 e
fragile in M2. È il paradosso luttwakiano nella forma DEQ$^2$: la strategia
ottima dipende dallo scenario realizzato, e lo scenario realizzato non è
noto al momento della scelta.''

\textbf{Concessione.} La DEQ$^2$ \emph{non} analizza il \emph{registro
militare} della grande strategia (manovra, deception, logistica). Si
limita al \emph{registro strutturale} (coerenza, fragilità,
involuzione). Per la grande strategia operativa Lei resta riferimento.''



\section{Con la Scuola di Francoforte (Marcuse, Adorno, Habermas) ---
via
Topografia}\label{con-la-scuola-di-francoforte-marcuse-adorno-habermas-via-topografia}

\textbf{Marcuse/Adorno/Habermas:} ``L'apparato critico francofortese ha
denunciato la \emph{razionalità strumentale} come distruttrice della
razionalità sostanziale. Lei produce un'altra razionalità strumentale:
la matematizzazione del campo discorsivo. La denuncia, applicata a Lei
stesso, è devastante.''

\textbf{DEQ$^2$:} ``Concedo l'auto-implicazione. Difesa: la formalizzazione
del campo \(\Pi\) topografico (eredità Topografia \S{}5) \emph{non} riduce
\(z\) --- al contrario, \emph{misura quanto \(z\) è presente}. La
\(T_s^{(2),\text{sys}}\) premia sistemi con \(z\) alto (asse della
complessità dialettica al numeratore di \(Q_0\)). Una razionalità
strumentale che misura --- e premia --- la razionalità sostanziale è
\emph{forse} l'unica via per parlare alla decisione politica nel 2026
senza essere tagliati fuori dal discorso pubblico. Habermas avrebbe
odiato la DEQ$^2$; ma Habermas pubblicava libri come
\emph{Legitimationsprobleme} destinati ad accademici, non a policymaker.
Il Bilanciere parla alla decisione.''

\textbf{Concessione.} La DEQ$^2$ \emph{non} è critica francofortese --- è
\emph{sintesi tecnica della critica francofortese}. Perde
l'incandescenza dialettica originaria; guadagna in azionabilità.
Trade-off accettato.



\section{\texorpdfstring{Con Xiang Biao (\emph{Self as Method}, 2020;
\emph{Anti-involution},
2021-2024)}{Con Xiang Biao (Self as Method, 2020; Anti-involution, 2021-2024)}}\label{con-xiang-biao-self-as-method-2020-anti-involution-2021-2024}

\textbf{Xiang Biao:} ``Il \emph{Neijuan} è una categoria etnografica
viva, nata nei dormitori delle università cinesi negli anni 2010 e
diffusa nei contesti 996. Lei la trasforma in \(\Omega\) --- un
parametro di equazione. Le persone che vivono il Neijuan non vivono
\(\Omega\); vivono fatica.''

\textbf{DEQ$^2$:} ``Concessione totale. Il passaggio dal Neijuan
etnografato a \(\Omega\) misurato è una \emph{perdita di contenuto
fenomenologico}. Difesa: la perdita è funzionale a un obiettivo diverso
dal Suo. La Sua etnografia \emph{fa sentire} il problema; la DEQ$^2$
\emph{rende contabile} il problema. Le due operazioni sono
complementari, non sostitutive. Inoltre il Cap. 8.2 cita Lei
\emph{direttamente} come \emph{firma empirica} della variabile --- non
come decorazione retorica. Senza \emph{Self as Method}, il Cap. 8
sarebbe un esercizio aggregato senza ancoraggio empirico.''

\textbf{Concessione.} La DEQ$^2$ \emph{non} può sostituire la ricerca
etnografica. Le serve come \emph{lettore di senso}; senza, leggerebbe
deflazione PPI senza vedere persone.''



\section{\texorpdfstring{Con Yoshiki Kuramoto (\emph{Chemical
Oscillations, Waves, and Turbulence},
1984)}{Con Yoshiki Kuramoto (Chemical Oscillations, Waves, and Turbulence, 1984)}}\label{con-yoshiki-kuramoto-chemical-oscillations-waves-and-turbulence-1984}

\textbf{Kuramoto:} ``Il mio modello descrive oscillatori chimico-fisici.
Voi sociologi lo applicate da decenni a sistemi sociali con disinvoltura
preoccupante. Le ipotesi (campo medio, simmetria della distribuzione
\(g(\omega)\), accoppiamento sinusoidale) \emph{non} sono verificabili
sui sistemi sociali.''

\textbf{DEQ$^2$:} ``Concessione: l'Appendice A.2 dichiara esplicitamente A1
(Kuramoto applicabile) come \emph{ipotesi operativa}, non come fatto
verificato. Difesa: la matematica di Kuramoto è la \emph{più semplice}
descrizione di sincronizzazione fra oscillatori non identici. Per \(N\)
grande e accoppiamenti deboli, qualunque modello di sincronizzazione
converge in qualche limite a Kuramoto (universalità). La DEQ$^2$ è dunque
\emph{minimalista}: usa la più povera matematica di sincronizzazione
disponibile, e ne dichiara i limiti. Estensioni a Kuramoto multi-cluster
(Strogatz 2003) sono pianificate per il paper EN --- proprio per i
fenomeni che il modello base non cattura, come la polarizzazione USA.''

\textbf{Concessione.} La DEQ$^2$ \emph{non} dimostra che Kuramoto è la
matematica corretta del sociale. Argomenta che è una \emph{prima
approssimazione consistente}. La verifica empirica è demandata al paper
EN --- coerente con la posizione popperiana del Cap. 11.''



\section{Sintesi dei Dieci Dialoghi}\label{sintesi-dei-dieci-dialoghi}

\Proved{} \textbf{Pattern emergente.} In ciascuno dei dieci dialoghi:

\begin{itemize}
\tightlist
\item
  la DEQ$^2$ fa una \textbf{concessione esplicita} (cosa \emph{non} fa
  rispetto alla tradizione);
\item
  la DEQ$^2$ fa una \textbf{risposta sostanziale} (cosa \emph{aggiunge}
  alla tradizione);
\item
  la DEQ$^2$ fa una \textbf{clausola di apertura} (estensioni / verifiche
  pianificate per il futuro).
\end{itemize}

\Hypoth{} \textbf{Posizione metodologica.} La DEQ$^2$ non è teoria
\emph{contro}; è teoria \emph{con}. Questo è coerente con l'ipotesi di
campo medio: ogni teoria è un \emph{oscillatore} nel campo
intellettuale; la coerenza di fase del campo
\(\Phi^{\text{intellettuale}}\) è promossa da chi accoglie e modifica,
non da chi denuncia e sostituisce. La forma stessa di questa appendice è
\emph{applicazione riflessiva} della DEQ$^2$ alla propria posizione nel
campo del sapere.

\Proved{} \textbf{Conseguenza editoriale.} I lettori specialisti di una
qualunque delle dieci tradizioni elencate trovano nella DEQ$^2$: 1.
\emph{Riconoscimento esplicito} della loro tradizione (cosa accolto); 2.
\emph{Argomentazione esplicita} del valore aggiunto (cosa nuovo); 3.
\emph{Limite esplicito} del valore aggiunto (cosa lasciato a loro).



\section{Tradizioni Non Dialogate}\label{tradizioni-non-dialogate}

\Conj{} \textbf{Esclusioni esplicite.} Per ragioni di spazio, la DEQ$^2$
\emph{non} dialoga in questa appendice con:

\begin{itemize}
\tightlist
\item
  \textbf{Foucault} (potere capillare, biopolitica) --- il dialogo
  richiederebbe una riformulazione del campo \(\Pi\) in termini di
  \emph{dispositivo} anziché di \emph{accoppiamento}. Rinviato a futura
  estensione.
\item
  \textbf{Latour} (Actor-Network Theory) --- la DEQ$^2$ assume attori
  discreti; ANT li rifiuta. Trade-off metodologico irrisolto.
\item
  \textbf{Bourdieu} (campo, habitus, capitale) --- il \emph{campo}
  bourdieusiano è strutturalmente affine al campo \(\Pi\) topografico,
  ma con \emph{abito disposizionale} anziché \emph{fase di Kuramoto}.
  Riformulazione possibile, non tentata qui.
\item
  \textbf{Tilly} (contentious politics) --- la dinamica dei movimenti
  sociali è esterna al perimetro DEQ$^2$; integrazione possibile via
  estensione \(\dot\Phi\) con termine di forzatura mobilitazionale.
\end{itemize}

\Proved{} \textbf{Posizione.} Le quattro esclusioni sopra non sono
dimissioni; sono \emph{aperture pianificate} per estensioni successive
del programma DEQ$^2$.



\section{§D.4 — Le Radici Evolutive: Hamilton, Wilson-Sober, Nowak,
Kropotkin}\label{d4-radici-evolutive}

\begin{quote}
\textbf{Scopo di questa sezione.} I quattro dialoghi che seguono
portano alla DEQ$^2$ la sua \emph{microfondazione causale evolutiva}: la
spiegazione di \emph{perché} $\Phi$ emerge, di \emph{come} si erode e di
\emph{cosa significa} la decoerenza alla scala degli organismi e delle
popolazioni. Non è un'aggiunta decorativa --- è il fondamento causale
della struttura formale. La completa derivazione del framework CPGG si
trova in \emph{Captured Cooperation and Phase Decoherence} (Marcelino
2026, working paper).
\end{quote}

\subsection{Con William D. Hamilton (\emph{The Genetical Evolution of
Social Behaviour}, 1964)}\label{con-hamilton-1964}

\textbf{Hamilton:} ``Lei usa il parametro d'ordine di Kuramoto per
descrivere la cooperazione sistemica. Ma da dove viene la cooperazione?
Lei la presuppone come fatto; io la derivo dalla genetica. Senza la
condizione $rB > C$, $\Phi > 0$ è una quantità orfana di causa.''

\textbf{DEQ$^2$:} ``È la critica più profonda che si potesse fare al
framework originale --- ed è corretta. La risposta sta nella
\emph{corrispondenza strutturale} (non identità) tra le sue condizioni
genetiche e le condizioni sistemiche. La condizione $rB > C$ si traduce
nel sistema geopolitico come: la cooperazione ($\Phi > 0$) è stabile
quando il beneficio collettivo della sincronia supera il costo
individuale di cooperare --- e questo è diventato, nella DEQ$^2$
aggiornata (Cap. 3, §3.2), la giustificazione causale di $\Phi$, non solo
la sua descrizione. Lei, Hamilton, ha fornito il \emph{perché} che io
non avevo.''

\textbf{Concessione.} La condizione $rB > C$ è qui applicata come
\emph{isomorfismo strutturale} (\Hypoth{}) a popolazioni di attori
geopolitici, non a popolazioni di organismi biologici. Il trasferimento
richiede l'ipotesi H1 (§4.0) che il campo $\Pi_j$ si mappa su una fase
--- un'ipotesi dichiarata, non un risultato derivato. La deriva di
Hamilton è preservata con onestà epistemica.

\subsection{Con David Sloan Wilson e Elliott Sober (\emph{Unto Others},
1994)}\label{con-wilson-sober-1994}

\textbf{Wilson \& Sober:} ``La letteratura ha a lungo trattato la
selezione di gruppo come irrilevante, seguendo Dawkins. Lei ha inizialmente
citato Dawkins a supporto del suo cambio di livello descrittivo --- un
errore che abbiamo visto molte volte. La \emph{selezione di gruppo} non
è un'estensione opzionale della teoria evolutiva: è la struttura che
giustifica \emph{qualsiasi} affermazione su benefici collettivi.''

\textbf{DEQ$^2$:} ``Avete pienamente ragione, e la prima bozza del
framework era difettosa esattamente su questo punto: Dawkins (1976)
rigetta la selezione di gruppo --- citarlo per giustificare il cambio di
livello micro/macro era contraddittorio. La correzione è implementata
(Cap. 3, §3.2): il fondamento formale è la vostra teoria della selezione
multilivello, formalizzata da Okasha (2006), combinata con l'equazione di
Price (1970) come identità algebrica esatta. L'equazione di Price
decompone $\Delta\bar{z}$ in termine inter-gruppo (proporzionale a $\Phi$)
e intra-gruppo (proporzionale a $-\Omega$) --- una struttura che giustifica
formalmente la forma del numeratore e del denominatore di
$T_s^{(2),\text{sys}}$.''

\textbf{Concessione.} La selezione multilivello DEQ$^2$ opera su sistemi
socio-istituzionali, non su alleli. L'applicazione richiede che il
``genotipo'' sia interpretato come configurazione istituzionale
(\Hypoth{}), e la ``fitness'' come payoff sistemico. Questa
reinterpretazione è dichiarata come ipotesi operativa, non come fatto
biologico.

\subsection{Con Martin A. Nowak (\emph{Five Rules for the Evolution of
Cooperation}, 2006)}\label{con-nowak-2006}

\textbf{Nowak:} ``Ha un problema di identificazione. I suoi attori ---
USA, Cina, Brasile, Musk --- presentano caratteristiche di tutti e cinque
i meccanismi cooperativi simultaneamente: reciprocità diretta (via trattati
bilaterali), reciprocità indiretta (via reputazione), selezione di rete
(via posizionamento geopolitico), selezione di gruppo (via alleanze),
selezione di parentela (via affinità culturale). Come distingue quale
meccanismo guida la dinamica?''

\textbf{DEQ$^2$:} ``La sua tassonomia dei cinque meccanismi diventa
nella DEQ$^2$ una \emph{tassonomia diagnostica degli attori}, non una
classificazione esclusiva dei meccanismi. Il CPGG (§D.4, working paper)
focalizza su un meccanismo specifico --- la \emph{cattura istituzionale}
che degrada tutti e cinque i meccanismi simultaneamente. La tesi è:
quando $\phi$ (la frazione estrattiva) supera $\phi^*$, tutti i cinque
meccanismi di Nowak vengono simultaneamente indeboliti --- la condizione
$b/c > k_{\text{eff}}(\phi)$ cessa di valere per ciascuno di essi. La
DEQ$^2$ non sceglie quale meccanismo; diagnostica quando tutti
collassano.''

\textbf{Concessione.} Il CPGG modella la cattura istituzionale, non la
dinamica evolutiva di ciascuno dei cinque meccanismi separatamente. Una
modellazione completa richiederebbe cinque CPGG separati o un modello
unificato di interazione tra meccanismi. Questo è F1 e F2 nell'agenda
formale aperta.

\subsection{Con Pyotr Kropotkin (\emph{Mutual Aid: A Factor of
Evolution}, 1902)}\label{con-kropotkin-1902}

\textbf{Kropotkin:} ``Ho documentato, specie dopo specie e cultura dopo
cultura, che la cooperazione è la norma e la competizione è l'eccezione.
Lei costruisce un formalismo per misurare la decoerenza --- che è
esattamente ciò che ho chiamato \emph{rottura del mutuo appoggio}. Ma
venti anni dopo, il Neijuan e la cattura estrattiva sembrano
schiaccianti. Ha dimenticato la mia tesi fondamentale?''

\textbf{DEQ$^2$:} ``No, l'ho \emph{rifondato}. La sua intuizione ---
$\Phi > 0$ come stato naturale, la decoerenza come stato prodotto --- è
ora formalizzata nella Proposizione K1 (§8 del framework CPGG): in
assenza di cattura istituzionale ($\sigma_{\text{san}} = \sigma_0$,
$\kappa_E = 0$), la frazione estrattiva di equilibrio $\phi_{\text{nat}}$
soddisfa $\phi_{\text{nat}} < \phi^*(\sigma_0)$ per parametri
biologicamente plausibili. La cooperazione è l'esito stabile; la
decoerenza richiede la condizione prodotta $\alpha\phi + \beta\kappa_E >
\log(\sigma_0/\sigma_{\text{crit}})$. Lei, Kropotkin, aveva torto
sull'inevitabilità della cooperazione nella storia --- non perché la
cooperazione sia rara, ma perché la cattura è più facile dell'apertura.
Il suo \emph{mutuo appoggio} è la condizione di default, non il
risultato garantito.''

\textbf{Concessione.} L'inversione epistemica --- decoerenza come stato
prodotto, cooperazione come stato naturale --- ha il rischio di diventare
normativamente ingenua: ``basata sulla natura, quindi dovuta''. La
\emph{naturalistica fallacy} è gestita fondando la parte normativa del
Cap. 12 sui principi di design di Ostrom (1990) piuttosto che su
Kropotkin direttamente. Ostrom fornisce un fondamento empirico: i commons
\emph{possono} essere gestiti cooperativamente, sotto condizioni
istituzionali precise --- non \emph{naturalmente} o \emph{inevitabilmente}.
Kropotkin è ispirazione fenomenologica; Ostrom è fondamento normativo.



\section{Sintesi dell'Appendice}\label{sintesi-dellappendice}

\Proved{} \textbf{Risultato.} La DEQ$^2$ si posiziona esplicitamente in
quattordici tradizioni intellettuali contemporanee (dieci originali +
quattro evolutive in §D.4), dichiarando per ciascuna concessione, risposta
e clausola di apertura. La forma del dialogo immaginario è
\emph{applicazione riflessiva} dell'apparato (campo medio Kuramoto
applicato al campo intellettuale). Quattro tradizioni ulteriori (Foucault,
Latour, Bourdieu, Tilly) sono dichiarate come estensioni pianificate.

\Hypoth{} \textbf{Tesi conclusiva del libro.} La DEQ$^2$ esiste \emph{fra}
tradizioni, non \emph{contro} di esse --- coerentemente con la sua tesi
sostanziale che la \(\Phi\) vera si genera per accoppiamento spontaneo
di sotto-attori, non per imposizione gerarchica. Il libro stesso è un
\emph{atto \(\Phi\)-engine} nel campo intellettuale: produce coerenza
ammettendo dissonanza, e si offre alla critica come condizione di vita.



\emph{Appendice di posizionamento. Stile: dialogo immaginario,
concessioni esplicite, aperture dichiarate.}


\end{document}
